\chapter{同调镜像对称的数学}\label{chap:8}

\SourceParagraph{ch08-intro-001}{585}
本章的目标，是以数学上精确的形式陈述 Kontsevich 的同调镜像对称（HMS）猜想。主要将较详细地解释所涉及的代数结构，即 $A_\infty$ 代数与范畴。作为例子，将在最简单的椭圆曲线情形，给出一个部分证明。最后，简要讨论 Seidel 最近对四次 K3 曲面的 HMS 所作的证明。

\SourceParagraph{ch08-intro-002}{585}
粗略地说，HMS 断言：若 $X,\check X$ 是互为镜像的 Calabi--Yau 流形，则 $X$ 上 A 膜的范畴，等价于 $\check X$ 上 B 膜的范畴。数学上，\BookChapterRef{5}已说明，B 膜范畴为 $\check X$ 上凝聚层的导出范畴 $D^b(\check X)$，其严格定义见\BookSectionRef{sec:4.5}{4.5}。\BookSectionRef{sec:3.6}{3.6}讨论过，$X$ 上的 A 膜范畴为 Fukaya 范畴 $\operatorname{Fuk}(X)$；粗略地说，其对象为 $X$ 中的 Lagrange 子流形。因此，HMS 应给出 $D^b(\check X)$ 与 $\operatorname{Fuk}(X)$ 之间的范畴等价。

\SourceParagraph{ch08-intro-003}{585}
这一等价存在一个基本问题。$\operatorname{Fuk}(X)$ 实际上不是通常意义的范畴，而是所谓 $A_\infty$ 范畴。特别地，态射的复合不严格满足结合律，只在“同伦意义下”结合。另一方面，$D^b(\check X)$ 为三角范畴，$\operatorname{Fuk}(X)$ 却不是。因此，在 HMS 成为精确陈述之前，还需要进一步理解这些结构。

\SourceParagraph{ch08-intro-004}{585}
首先，必须理解什么是 $A_\infty$ 范畴，这需要理解相当复杂的一组 $A_\infty$ 关系。

\section{\texorpdfstring{$A_\infty$}{A∞} 代数与范畴}\label{sec:8.1}
\subsection{引言：\texorpdfstring{$A_\infty$}{A∞} 空间}\label{sec:8.1.1}

\SourceParagraph{ch08-8.1.1-001}{585}
$A_\infty$ 代数与范畴的定义，其动机不太容易说明。因此，先考察 $A_\infty$ 空间的思想，它由 Stasheff\BookCite{ref430}提出，\BookSectionRef{sec:2.5}{2.5}也已提及。设有拓扑空间 $X$ 及“乘法映射” $m_2\colon X\times X\to X$。此映射可能满足结合律，也可能不满足；结合律是一个额外条件。$A_\infty$ 空间要求较弱的结构：$m_2$ 在同伦意义下满足结合律，并且还有相应的“高阶”版本。确实，存在非常标准的情形，具有自然的乘法映射，虽然不满足严格结合律，却满足下面将描述的较弱条件。

\SourceParagraph{ch08-8.1.1-002}{586}
标准例子是另一个空间 $M$ 的基点环路空间 $X$。选定基点 $m_0\in M$，则
\[
X=\{x\colon[0,1]\to M\mid x\text{ 连续},\ x(0)=x(1)=m_0\}.
\]
定义环路复合为
\[
(x_2x_1)(t)=\begin{cases}
x_2(2t),&0\leq t\leq\tfrac12,\\
x_1(2t-1),&\tfrac12\leq t\leq1.
\end{cases}
\]
这一复合不满足严格结合律，但 $x_3(x_2x_1)$ 与 $(x_3x_2)x_1$ 为同伦环路。\BookRef{fig:8.1}{图1}示意了这一同伦。
\begin{figure}[H]
\centering
\BookImage{01_Numbered_Figures/Chapter_08/Ch08_Fig01_PDFp586_Bookp576.pdf}
\BookNumberOnlyCaption\label{fig:8.1}
\end{figure}

\SourceParagraph{ch08-8.1.1-003}{586}
左侧，在时间零至 $1/2$ 先走过 $x_3$，在 $1/2$ 至 $3/4$ 走过 $x_2$，再在 $3/4$ 至一走过 $x_1$。右侧，在零至 $1/4$ 先走过 $x_3$，在 $1/4$ 至 $1/2$ 走过 $x_2$，再在 $1/2$ 至一走过 $x_1$。连续改变这些时间，就能将一个环路同伦到另一个。这一同伦可表示为映射
\[
m_3\colon[0,1]\times X\times X\times X\longrightarrow X,
\]
在 $\{0\}\times X^3$ 上为
\[
(x_3,x_2,x_1)\longmapsto m_2(x_3,m_2(x_2,x_1)),
\]
在 $\{1\}\times X^3$ 上则为
\[
(x_3,x_2,x_1)\longmapsto m_2(m_2(x_3,x_2),x_1).
\]

\SourceParagraph{ch08-8.1.1-004}{586}
现在有一个问题：对于 $X$ 的四个元素 $x_1,\ldots,x_4$，其乘积有若干种不同加括号方式，$m_3$ 给出的同伦将这些方式联系起来。事实上，可以用两条不同路径连接 $((x_4x_3)x_2)x_1$ 与 $x_4(x_3(x_2x_1))$：
\[
((x_4x_3)x_2)x_1\sim(x_4x_3)(x_2x_1)\sim x_4(x_3(x_2x_1)),
\]
以及
\[
\begin{aligned}
((x_4x_3)x_2)x_1
&\sim(x_4(x_3x_2))x_1\\
&\sim x_4((x_3x_2)x_1)\sim x_4(x_3(x_2x_1)).
\end{aligned}
\]
每个 $\sim$ 都表示 $m_3$ 给出的同伦。示意地，可以用多边形 $\overline S_4$ 表示；每个顶点标记乘积 $x_4x_3x_2x_1$ 的一种加括号方式，边表示它们之间的同伦，见\BookRef{fig:8.2}{图2}。
\begin{figure}[H]
\centering
\BookImage{01_Numbered_Figures/Chapter_08/Ch08_Fig02_PDFp587_Bookp577.pdf}
\BookNumberOnlyCaption\label{fig:8.2}
\end{figure}

\SourceParagraph{ch08-8.1.1-005}{587}
换言之，$m_3$ 的同伦给出映射 $\partial\overline S_4\times X^4\to X$，在边界的每条边上，通过适当组合 $m_2,m_3$ 定义。例如，在连接 $((x_4x_3)x_2)x_1$ 与 $(x_4(x_3x_2))x_1$ 的边上，此映射为
\[
(s,x_4,\ldots,x_1)\longmapsto m_2(m_3(s,x_4,x_3,x_2),x_1).
\]

\SourceParagraph{ch08-8.1.1-006}{587}
随后可以对目前的结构施加一个新条件：要求此映射延伸到整个 $\overline S_4$，给出
\[
m_4\colon\overline S_4\times X^4\longrightarrow X.
\]
再看乘积 $x_5x_4x_3x_2x_1$ 的各种加括号方式。用\BookRef{fig:8.3}{图3}中的 $\overline S_5$ 表示这些不同选择。
\begin{figure}[H]
\centering
\BookImage{01_Numbered_Figures/Chapter_08/Ch08_Fig03_PDFp588_Bookp578.pdf}
\BookNumberOnlyCaption\label{fig:8.3}
\end{figure}

\SourceParagraph{ch08-8.1.1-007}{587}
同样，每个顶点表示一种加括号方式，边表示它们之间的同伦。二维面有两种：五边形显然是 $\overline S_4$ 的副本；四边形则来自以两种不同顺序应用两个结合变换，可以看作 $\overline S_3\times\overline S_3$，其中 $\overline S_3=[0,1]$。注意，每个面都对应于给 $x_5\cdots x_1$ 选一对括号，排除平凡选择 $(x_5\cdots x_1)$ 和只括住单个元素的 $x_5\cdots(x_i)\cdots x_1$。例如，左上前方的五边形对应于 $5(4321)$ 的加括号方式，左侧四边形对应于 $5(432)1$，右上前方的五边形对应于 $543(21)$。

\SourceParagraph{ch08-8.1.1-008}{588}
与先前一样，由此定义映射 $\partial\overline S_5\times X^5\to X$。例如，左上五边形上的映射为
\[
(s,x_5,\ldots,x_1)\longmapsto m_2(x_5,m_4(s,x_4,x_3,x_2,x_1));
\]
左侧四边形写作 $\overline S_3\times\overline S_3$，其上的映射为
\[
((s,t),x_5,\ldots,x_1)\longmapsto
m_3(s,x_5,m_3(t,x_4,x_3,x_2),x_1).
\]
右上五边形上则有
\[
(s,x_5,\ldots,x_1)\longmapsto m_4(s,x_5,x_4,x_3,m_2(x_2,x_1)).
\]
由于 $m_4$ 在 $\partial\overline S_4$ 上本身限制为 $m_2,m_3$ 的复合，可以验证，各面上定义的映射彼此相容，确实给出 $\partial\overline S_5\times X^5\to X$。下一阶“高阶结合律”关系，若存在，就将它延伸为
\[
m_5\colon\overline S_5\times X^5\longrightarrow X.
\]

\SourceParagraph{ch08-8.1.1-009}{588}
Stasheff\BookCite{ref430}构造了这类“结合多面体”的无限序列
\[
\overline S_2=\{\mathrm{pt}\},\qquad\overline S_3=[0,1],
\qquad\overline S_4,\ldots.
\]
$\overline S_d$ 的顶点对应于将 $x_d\cdots x_1$ 完全加括号的方式，其极大真面，与选取一对非平凡括号的方式一一对应。具体地，括号
\[
x_d\cdots x_{p+q+1}(x_{p+q}\cdots x_{q+1})x_q\cdots x_1
\]
对应的面同构于 $\overline S_{d-p+1}\times\overline S_p$。

\SourceParagraph{ch08-8.1.1-010}{589}
$\overline S_d$ 的一种描述特别有用：将其看作带状树的模空间。
\BookStatementNumber{1}
\begin{definition}\label{def:8.1}
度量带状树，是具有有限个顶点与边、没有二价顶点的连通树，并在每个顶点赋予相接边的循环次序，以及为每条边指定属于 $(0,\infty]$ 的长度。
\end{definition}

\SourceParagraph{ch08-8.1.1-011}{589}
假设这棵树有 $d+1$ 个外部顶点，也叫叶顶点，即一价顶点。选择其中一个作为出顶点，其余为入顶点。循环次序连同出顶点的选择，允许按与循环次序相容的方式，唯一地将入顶点标记为 $v_1,\ldots,v_d$。将树定向，使所有边指向出顶点。\BookRef{fig:8.4}{图4}给出一个例子，半圆形箭头表示循环次序。若忘掉边长，并且所有内部顶点均为三价，就可以将树看作 $x_d\cdots x_1$ 的一种加括号方式。
\begin{figure}[H]
\centering
\BookImage{01_Numbered_Figures/Chapter_08/Ch08_Fig04_PDFp589_Bookp579.pdf}
\BookNumberOnlyCaption\label{fig:8.4}
\end{figure}

\SourceParagraph{ch08-8.1.1-012}{589}
\BookStatementNumber{2}
\begin{definition}\label{def:8.2}
$\overline S_d$ 是具有 $d+1$ 个外部顶点的带状树的模空间，要求每条与外部顶点相接的边，即外部边，长度为无穷。它是所有内部边长有限的带状树模空间 $S_d\subseteq\overline S_d$ 的紧化。树按保持出顶点、循环次序及边长的同构识别；零长内部边通过收缩识别。
\end{definition}

\SourceParagraph{ch08-8.1.1-013}{590}
显然，$\overline S_2=\{\mathrm{pt}\}$，由\BookRef{fig:8.5}{图5}的树构成。
\begin{figure}[H]
\centering
\BookImage{01_Numbered_Figures/Chapter_08/Ch08_Fig05_PDFp590_Bookp580.pdf}
\BookNumberOnlyCaption\label{fig:8.5}
\end{figure}
一般而言，具有 $d+1$ 个外部顶点、内部顶点均为三价的带状树，有 $d-2$ 条内部边，每条边长属于 $(0,\infty]$。不过，可以考虑某条边长趋于零的极限，并将其看作收缩该边。因此，内部边长属于 $[0,\infty]$ 的任意这种三价树，都给出一个带状树，从而在 $\overline S_d$ 中给出立方体 $[0,\infty]^{d-2}$。所以，$\overline S_d$ 有一个立方体细分，各立方体由三价树的组合类型标记。例如，$d=3$ 时，得到\BookRef{fig:8.6}{图6}中的区间 $\overline S_3$ 的细分；$d=4$ 时，则得到\BookRef{fig:8.7}{图7}所示的细分。
\begin{figure}[H]
\centering
\BookImage{01_Numbered_Figures/Chapter_08/Ch08_Fig06_PDFp590_Bookp580.pdf}
\BookNumberOnlyCaption\label{fig:8.6}
\end{figure}
\begin{figure}[H]
\centering
\BookImage{01_Numbered_Figures/Chapter_08/Ch08_Fig07_PDFp591_Bookp581.pdf}
\BookNumberOnlyCaption\label{fig:8.7}
\end{figure}

\SourceParagraph{ch08-8.1.1-014}{590}
现在，$\overline S_d$ 的面对应于具有无穷长内部边的树。一般可以将 $\overline S_d$ 看作多胞体，每个极大真面，像先前一样由一对非平凡括号给出，即
\setcounter{equation}{0}
\begin{equation}\label{eq:8.1}
x_d\cdots x_{p+q+1}(x_{p+q}\cdots x_{q+1})x_q\cdots x_1.
\end{equation}
它对应于一对树 $(T,T')\in\overline S_{d-p+1}\times\overline S_p$：将 $T'$ 的出边与 $T$ 的第 $q+1$ 条入边识别，并令此内部边长度为无穷，就得到 $\overline S_d$ 中的树，如\BookRef{fig:8.8}{图8}所示。
\begin{figure}[H]
\centering
\BookImage{01_Numbered_Figures/Chapter_08/Ch08_Fig08_PDFp591_Bookp581.pdf}
\BookNumberOnlyCaption\label{fig:8.8}
\end{figure}

\SourceParagraph{ch08-8.1.1-015}{591}
现在可以给出 $A_\infty$ 空间的一般定义。
\begin{definition*}\label{def:8.ainfty-space}
$A_\infty$ 空间为带有一列映射 $m_d\colon\overline S_d\times X^d\to X$、$d\geq2$ 的拓扑空间 $X$，这些映射在 $\overline S_d$ 的边界上，满足由边界组合结构指定的相容条件。此条件推广了上述小 $d$ 的例子：将 $m_d$ 限制到\BookEqRef{eq:8.1}{8.1}的括号确定的边界面时，得到映射
\[
\overline S_{d-p+1}\times\overline S_p\times X^d\longrightarrow X,
\]
其表达式为
\[
\begin{aligned}
(T,T',x_d,\ldots,x_1)\longmapsto
m_{d-p+1}\bigl(&T,x_d,\ldots,x_{p+q+1},\\
&m_p(T',x_{p+q},\ldots,x_{q+1}),x_q,\ldots,x_1\bigr).
\end{aligned}
\]
\end{definition*}

\SourceParagraph{ch08-8.1.1-016}{592}
\BookStatementNumber{3}
\begin{note}[numbered]\label{note:8.3}
$\overline S_d$ 的另一种解释，是亏格零、带 $d+1$ 个标记点的稳定曲线模空间 $\overline{\mathcal M}_{0,d+1}$ 的实点集的一个子集。其开部分 $\mathcal M_{0,d+1}$ 的实点集，包含标记点 $x_1,\ldots,x_d,y$ 均位于赤道、且彼此不同的 Riemann 球面。$S_d$ 可以识别为这些点按给定循环次序排列的分支，$\overline S_d$ 则为其闭包。也可以将此模空间看作带 $d+1$ 个标记点的圆盘模空间，标记点 $x_1,\ldots,x_d,y$ 按循环次序位于圆盘边界上。在 $\overline S_d\setminus S_d$ 处，圆盘发生冒泡；例如，$d=4$ 的情形见\BookRef{fig:8.9}{图9}。
\end{note}
\begin{figure}[H]
\centering
\BookImage{01_Numbered_Figures/Chapter_08/Ch08_Fig09_PDFp592_Bookp582.pdf}
\caption{与\BookRef{fig:8.7}{图7}中的树比较。}\label{fig:8.9}
\end{figure}

\subsection{\texorpdfstring{$A_\infty$}{A∞} 代数与范畴}\label{sec:8.1.2}

\SourceParagraph{ch08-8.1.2-001}{593}
现在定义上述 $A_\infty$ 空间的代数及范畴论对应物。有了前面对空间的描述，这些定义应该更自然。基本思想是，$A_\infty$ 代数为带有高阶乘法的非结合代数，高阶乘法衡量结合律不成立的程度。

\SourceParagraph{ch08-8.1.2-002}{593}
以下始终固定底域 $k$。

\SourceParagraph{ch08-8.1.2-003}{593}
\BookStatementNumber{4}
\begin{definition}\label{def:8.4}
非幺 $A_\infty$ 代数 $A$ 是 $\mathbb Z$ 分次的 $k$ 向量空间 $A=\bigoplus_{p\in\mathbb Z}A^p$，配有对每个 $d\geq1$ 定义的 $k$ 线性分次映射
\[
m_d\colon A^{\otimes d}\longrightarrow A,
\]
次数为 $2-d$，并且对每个 $d\geq1$ 及齐次元素 $a_1,\ldots,a_d$，满足关系
\begin{equation}\label{eq:8.2}
\begin{aligned}
\sum_{\substack{1\leq p\leq d\\0\leq q\leq d-p}}
(-1)^{\deg a_1+\cdots+\deg a_q-q}
m_{d-p+1}\bigl(&a_d,\ldots,a_{p+q+1},\\
&m_p(a_{p+q},\ldots,a_{q+1}),a_q,\ldots,a_1\bigr)=0.
\end{aligned}
\end{equation}
\end{definition}

\SourceParagraph{ch08-8.1.2-004}{593}
对于小 $d$，试着展开这一定义。先取 $d=1$：$m_1$ 为次数一的映射，\BookEqRef{eq:8.2}{8.2}化为 $m_1(m_1(a_1))=0$。因此，$m_1$ 将 $A$ 变为一个上链复形。

\SourceParagraph{ch08-8.1.2-005}{593}
$m_2$ 的次数为零，即 $\deg m_2(a_2,a_1)=\deg a_2+\deg a_1$，应将其看作乘法映射。$d=2$ 时，\BookEqRef{eq:8.2}{8.2}中的 $(p,q)$ 可取 $(1,0)$、$(1,1)$、$(2,0)$，得到
\begin{equation}\label{eq:8.3}
\begin{aligned}
m_2(a_2,m_1(a_1))
&+(-1)^{\deg a_1-1}m_2(m_1(a_2),a_1)\\
&+m_1(m_2(a_2,a_1))=0.
\end{aligned}
\end{equation}

\SourceParagraph{ch08-8.1.2-006}{593}
除符号外，这就是说，$m_1$ 相对于乘法为分次导子。更准确地，设
\[
\partial(a)=(-1)^{\deg a}m_1(a),\qquad
a_2\cdot a_1=(-1)^{\deg a_1}m_2(a_2,a_1),
\]
则\BookEqRef{eq:8.3}{8.3}可以改写为
\begin{equation}\label{eq:8.4}
\partial(a_2\cdot a_1)
=(\partial a_2)\cdot a_1+(-1)^{\deg a_2}a_2\cdot(\partial a_1).
\end{equation}
若乘法满足结合律，就得到所谓微分分次代数。不过，考察 $d=3$ 时的关系，此时可取
\[
(p,q)=(1,0),(1,1),(1,2),(2,0),(2,1),(3,0).
\]

\SourceParagraph{ch08-8.1.2-007}{594}
由此得到
\begin{equation}\label{eq:8.5}
\begin{aligned}
0={}&m_3(a_3,a_2,m_1(a_1))
+(-1)^{\deg a_1-1}m_3(a_3,m_1(a_2),a_1)\\
&+(-1)^{\deg a_1+\deg a_2-2}m_3(m_1(a_3),a_2,a_1)\\
&+m_2(a_3,m_2(a_2,a_1))
+(-1)^{\deg a_1-1}m_2(m_2(a_3,a_2),a_1)\\
&+m_1(m_3(a_3,a_2,a_1)).
\end{aligned}
\end{equation}
第四、第五项为
\[
(-1)^{\deg a_2}a_3\cdot(a_2\cdot a_1)
-(-1)^{\deg a_2}(a_3\cdot a_2)\cdot a_1,
\]
因此，此表达式告诉我们，乘法偏离结合律的程度。

\SourceParagraph{ch08-8.1.2-008}{594}
特别地，用 $H^*(A)$ 表示 $A$ 相对于 $\partial$，等价地相对于 $m_1$ 的上同调。\BookEqRef{eq:8.4}{8.4}说明，若 $\partial a_2=\partial a_1=0$，且 $[a_2],[a_1]\in H^*(A)$ 分别由它们代表，则乘法 $[a_2]\cdot[a_1]=[a_2\cdot a_1]$ 定义良好。若 $i=1,2,3$ 时均有 $\partial a_i=0$，则\BookEqRef{eq:8.5}{8.5}给出
\[
a_3\cdot(a_2\cdot a_1)-(a_3\cdot a_2)\cdot a_1
=\pm\partial(m_3(a_3,a_2,a_1)),
\]
所以 $H^*(A)$ 上的乘法满足结合律。

\SourceParagraph{ch08-8.1.2-009}{594}
对于任意 $d$，得到的 $A_\infty$ 关系显然为两类项之和。$1<p<d$ 的项，对应于给 $a_d\cdots a_1$ 加一对非平凡括号的方式，从而与 $\overline S_d$ 的余维一面一一对应。$p=1$ 与 $p=d$ 的项则给出
\[
\begin{aligned}
&\pm m_d(a_d,\ldots,m_1(a_1))\pm\cdots\\
&\pm m_d(m_1(a_d),\ldots,a_1)
\pm m_1(m_d(a_d,\ldots,a_1)).
\end{aligned}
\]
这给出 $m_d$ 与低阶乘法之间的关系。可以将其看作 $A_\infty$ 空间定义的对应物：在那里，$m_d$ 将由低阶乘法在 $\partial\overline S_d$ 上定义的映射，延伸到整个 $\overline S_d$。

\SourceParagraph{ch08-8.1.2-010}{594}
还希望定义 $A_\infty$ 代数的态射。
\BookStatementNumber{5}
\begin{definition}\label{def:8.5}
$A_\infty$ 代数态射 $f\colon A\to B$，为一族次数 $1-d$ 的分次映射 $f_d\colon A^{\otimes d}\to B$，对每个 $d\geq1$ 满足
\begin{equation}\label{eq:8.6}
\begin{aligned}
&\sum_{r}\sum_{s_1,\ldots,s_r}
m_r^B\bigl(f_{s_r}(a_d,\ldots,a_{d-s_r+1}),\ldots,
f_{s_1}(a_{s_1},\ldots,a_1)\bigr)\\
&\quad=\sum_{p,q}(-1)^{\deg a_1+\cdots+\deg a_q-q}\\
&\qquad\quad\cdot f_{d-p+1}\bigl(a_d,\ldots,a_{p+q+1},
m_p^A(a_{p+q},\ldots,a_{q+1}),a_q,\ldots,a_1\bigr).
\end{aligned}
\end{equation}
左侧对 $r\geq1$ 及正整数分拆 $s_1+\cdots+s_r=d$ 求和，右侧对 $1\leq p\leq d$、$0\leq q\leq d-p$ 求和。若 $f_1$ 为通常上链复形意义下的拟同构，则称 $f$ 为拟同构。
\end{definition}

\SourceParagraph{ch08-8.1.2-011}{595}
再次展开这一定义。若所有 $n\geq2$ 都有 $f_n=0$，则对每个 $d$，只剩关系
\begin{equation}\label{eq:8.7}
m_d^B(f_1(a_d),\ldots,f_1(a_1))
=f_1(m_d^A(a_d,\ldots,a_1)).
\end{equation}
此时，称态射严格。

\SourceParagraph{ch08-8.1.2-012}{595}
若态射不严格，高阶 $f_n$ 衡量\BookEqRef{eq:8.7}{8.7}不成立的程度。特别地，对小 $d$，\BookEqRef{eq:8.6}{8.6}给出
\[
d=1:\qquad m_1^B(f_1(a_1))=f_1(m_1^A(a_1)),
\]
所以，$f_1$ 为上链复形态射。$d=2$ 时，则有
\begin{equation}\label{eq:8.8}
\begin{aligned}
m_2^B(f_1(a_2),f_1(a_1))+m_1^B(f_2(a_2,a_1))
={}&f_2(a_2,m_1^A(a_1))\\
&+(-1)^{\deg a_1-1}f_2(m_1^A(a_2),a_1)\\
&+f_1(m_2^A(a_2,a_1)).
\end{aligned}
\end{equation}

\SourceParagraph{ch08-8.1.2-013}{595}
可以复合 $A_\infty$ 代数态射 $f_d\colon A^{\otimes d}\to B$ 与 $g_d\colon B^{\otimes d}\to C$，定义为
\[
\begin{aligned}
(g\circ f)_d(a_d,\ldots,a_1)
=\sum_r\sum_{s_1,\ldots,s_r}
g_r\bigl(&f_{s_r}(a_d,\ldots,a_{d-s_r+1}),\ldots,\\
&f_{s_1}(a_{s_1},\ldots,a_1)\bigr),
\end{aligned}
\]
其中 $r\geq1$、$s_i\geq1$、$s_1+\cdots+s_r=d$。

\SourceParagraph{ch08-8.1.2-014}{595}
上述概念容易推广到范畴。
\BookStatementNumber{6}
\begin{definition}\label{def:8.6}
非幺 $A_\infty$ 范畴 $\mathcal A$ 由对象集合 $\operatorname{Ob}\mathcal A$，每对对象 $X_0,X_1$ 对应的 $\mathbb Z$ 分次 $k$ 向量空间 $\operatorname{Hom}_{\mathcal A}(X_0,X_1)$，以及对每个 $d\geq1$ 定义的 $k$ 线性复合映射
\[
\begin{aligned}
m_d\colon\operatorname{Hom}_{\mathcal A}(X_{d-1},X_d)
\otimes\cdots\otimes\operatorname{Hom}_{\mathcal A}(X_0,X_1)
\longrightarrow\operatorname{Hom}_{\mathcal A}(X_0,X_d)
\end{aligned}
\]
构成。这些映射的次数为 $2-d$，并满足
\[
\begin{aligned}
\sum_{p,q}(-1)^{\deg a_1+\cdots+\deg a_q-q}
m_{d-p+1}\bigl(&a_d,\ldots,a_{p+q+1},\\
&m_p(a_{p+q},\ldots,a_{q+1}),a_q,\ldots,a_1\bigr)=0,
\end{aligned}
\]
其中求和范围与\BookRef{def:8.4}{定义8.4}相同。

\SourceParagraph{ch08-8.1.2-015}{595}
两个 $A_\infty$ 范畴之间的 $A_\infty$ 函子，由对象映射 $F\colon\operatorname{Ob}\mathcal A\to\operatorname{Ob}\mathcal B$，以及次数 $1-d$ 的映射
\[
\begin{aligned}
F_d\colon\operatorname{Hom}_{\mathcal A}(X_{d-1},X_d)
\otimes\cdots\otimes\operatorname{Hom}_{\mathcal A}(X_0,X_1)
\longrightarrow\operatorname{Hom}_{\mathcal B}(F(X_0),F(X_d))
\end{aligned}
\]
构成，满足
\[
\begin{aligned}
&\sum_r\sum_{s_1,\ldots,s_r}
m_r^B\bigl(F_{s_r}(a_d,\ldots,a_{d-s_r+1}),\ldots,
F_{s_1}(a_{s_1},\ldots,a_1)\bigr)\\
&\quad=\sum_{p,q}(-1)^{\deg a_1+\cdots+\deg a_q-q}\\
&\qquad\quad\cdot F_{d-p+1}\bigl(a_d,\ldots,a_{p+q+1},
m_p^A(a_{p+q},\ldots,a_{q+1}),a_q,\ldots,a_1\bigr),
\end{aligned}
\]
各指标的范围与\BookRef{def:8.5}{定义8.5}相同。

\SourceParagraph{ch08-8.1.2-016}{596}
若对每对 $X_0,X_1$，映射
\[
F_1\colon\operatorname{Hom}_{\mathcal A}(X_0,X_1)
\longrightarrow\operatorname{Hom}_{\mathcal B}(F(X_0),F(X_1))
\]
均为拟同构，则称 $F$ 在态射复形上为拟同构，即拟全忠实。若对所有 $d>1$，均有 $F_d=0$，则称 $F$ 严格。
\end{definition}

\SourceParagraph{ch08-8.1.2-017}{596}
给定 $A_\infty$ 范畴 $\mathcal A$，还可以得到其他几个“范畴”。$H^*(\mathcal A)$ 的对象为 $\operatorname{Ob}\mathcal A$，态射为
\[
\operatorname{Hom}_{H^*(\mathcal A)}(X_0,X_1)
=H^*(\operatorname{Hom}_{\mathcal A}(X_0,X_1)).
\]
它几乎是通常意义的范畴：现在复合满足结合律，但未必存在恒等态射。此外，还有 $H^0(\mathcal A)$，对象同样为 $\operatorname{Ob}\mathcal A$，态射为
\[
\operatorname{Hom}_{H^0(\mathcal A)}(X_0,X_1)
=H^0(\operatorname{Hom}_{\mathcal A}(X_0,X_1)).
\]
同样，在未施加幺性条件时，它可以不具有恒等态射。

% 第8.2节；首次阅读翻译，按用户要求不进行额外复核。
\section{例子与构造}\label{sec:8.2}
\setcounter{equation}{8}
\setcounter{figure}{9}
\subsection{凝聚层}\label{sec:8.2.1}

\SourceParagraph{ch08-8.2.1-001}{596}
这里先考虑一个最容易描述的 $A_\infty$ 范畴例子，并借此引出构造 $A_\infty$ 结构的一般方法。

\SourceParagraph{ch08-8.2.1-002}{596}
设 $X$ 为域 $k$ 上的非奇异代数簇，带有仿射开覆盖 $\mathcal U=\{U_i\}_{i\in I}$，其中 $I$ 为有序指标集。用 $D^b_\infty(X)$ 表示以局部自由层的有界复形为对象的范畴。为定义态射，先定义复形 $\mathcal{H}om^\bullet_{\mathcal O_X}(\mathcal E^\bullet,\mathcal F^\bullet)$，其中
\[
\mathcal{H}om^q_{\mathcal O_X}(\mathcal E^\bullet,\mathcal F^\bullet)
=\bigoplus_m\mathcal{H}om_{\mathcal O_X}(\mathcal E^m,\mathcal F^{m+q})
\]
由次数为 $q$ 的分次映射组成，而
\[
\delta\colon
\mathcal{H}om^q_{\mathcal O_X}(\mathcal E^\bullet,\mathcal F^\bullet)
\longrightarrow
\mathcal{H}om^{q+1}_{\mathcal O_X}(\mathcal E^\bullet,\mathcal F^\bullet)
\]
将 $(f_m\colon\mathcal E^m\to\mathcal F^{m+q})_m$ 映为
\[
(\delta(f))_m
=d_{\mathcal F}\circ f_m-(-1)^q f_{m+1}\circ d_{\mathcal E}
\colon\mathcal E^m\longrightarrow\mathcal F^{m+q+1}.
\]
接着，取 $\mathcal{H}om^\bullet_{\mathcal O_X}(\mathcal E^\bullet,\mathcal F^\bullet)$ 的 Čech 复形的全复形：
\[
\operatorname{Hom}^n_{D^b_\infty(X)}(\mathcal E^\bullet,\mathcal F^\bullet)
=\bigoplus_{p+q=n}
\check C^p\bigl(\mathcal U,
\mathcal{H}om^q_{\mathcal O_X}(\mathcal E^\bullet,\mathcal F^\bullet)\bigr).
\]

\SourceParagraph{ch08-8.2.1-003}{597}
这里
\[
\check C^p(\mathcal U,\mathcal G)
=\bigoplus_{i_0<\cdots<i_p}\Gamma(U_{i_0\cdots i_p},\mathcal G),
\qquad
U_{i_0\cdots i_p}=U_{i_0}\cap\cdots\cap U_{i_p}.
\]
在复形 $\operatorname{Hom}^\bullet_{D^b_\infty(X)}(\mathcal E^\bullet,\mathcal F^\bullet)$ 上，可以利用全微分
\[
\partial=d+(-1)^p\delta
\]
定义微分 $\partial$，其中 $d\colon\check C^p(\mathcal U,\mathcal G)\to\check C^{p+1}(\mathcal U,\mathcal G)$ 为通常的 Čech 微分，其定义为
\[
(d\alpha)_{i_0\cdots i_{p+1}}
=\sum_{j=0}^{p+1}(-1)^j
\alpha_{i_0\cdots\widehat{i_j}\cdots i_{p+1}}
\big|_{U_{i_0\cdots i_{p+1}}}.
\]

\SourceParagraph{ch08-8.2.1-004}{597}
这一复形与导出范畴中通常的态射概念之间，具有如下关系：
\[
H^i\bigl(\operatorname{Hom}^\bullet_{D^b_\infty(X)}
(\mathcal E^\bullet,\mathcal F^\bullet)\bigr)
=\operatorname{Ext}^i(\mathcal E^\bullet,\mathcal F^\bullet).
\]
见\BookSectionRef{sec:4.5}{4.5}。暂且不取上同调，而把 $D^b_\infty(X)$ 看成一个 $A_\infty$ 范畴。按照已经选定的符号约定，定义
\[
m_1(f)=(-1)^{\deg f}\partial f.
\]
复合的定义如下。设
\[
\begin{aligned}
f&=(f_{n,i_0\cdots i_p})\in
\check C^p\bigl(\mathcal U,
\mathcal{H}om^q_{\mathcal O_X}(\mathcal E^\bullet,\mathcal F^\bullet)\bigr),\\
g&=(g_{n',j_0\cdots j_{p'}})\in
\check C^{p'}\bigl(\mathcal U,
\mathcal{H}om^{q'}_{\mathcal O_X}(\mathcal F^\bullet,\mathcal G^\bullet)\bigr).
\end{aligned}
\]
则复合 $g\circ f$ 为
\[
\begin{aligned}
(g\circ f)_{n,i_0\cdots i_{p+p'}}
&=\left((-1)^{q'p}
 g_{n+q,i_0\cdots i_{p'}}\circ
 f_{n,i_{p'}\cdots i_{p+p'}}\right)
 \big|_{U_{i_0\cdots i_{p+p'}}},\\
(g\circ f)_{n,i_0\cdots i_{p+p'}}&\colon
\mathcal E^n\big|_{U_{i_0\cdots i_{p+p'}}}
\longrightarrow
\mathcal G^{n+q+q'}\big|_{U_{i_0\cdots i_{p+p'}}}.
\end{aligned}
\]
再定义
\[
m_2(g,f)=(-1)^{\deg f}g\circ f.
\]
容易检验\BookEqRef{eq:8.3}{8.3}成立，等价地，\BookEqRef{eq:8.4}{8.4}成立。此外，层态射的复合本来就满足结合律，因此可令 $m_n=0$，$n\geq3$，从而得到一个 $A_\infty$ 范畴。

\SourceParagraph{ch08-8.2.1-005}{597}
当然，这比 $A_\infty$ 范畴的条件更强：它实际上是一个 DG 范畴，即微分分次范畴，其高阶乘法映射全部为零。既然这里的态射尚不是我们想要的态射，$A_\infty$ 结构也并不有趣，那么它为何会是我们所需要的东西？

\SourceParagraph{ch08-8.2.1-006}{597}
答案来自以下代数结论，最初由 Kadeishvili\BookCite{ref279}提出。我们将给出它在范畴情形中的一个版本，这一版本归于 Kontsevich 和 Soibelman\BookCite{ref311}。

\SourceParagraph{ch08-8.2.1-007}{597}
\BookStatementNumber{7}
\begin{theorem}\label{thm:8.7}
给定 $A_\infty$ 代数 $A$，在 $H^*(A)$ 上存在一个满足 $m_1=0$ 的 $A_\infty$ 代数结构，以及 $A_\infty$ 代数的拟同构 $i\colon H^*(A)\to A$。
\end{theorem}

\SourceParagraph{ch08-8.2.1-008}{598}
下面对任意 $A_\infty$ 范畴 $\mathcal A$ 解释一个更一般的构造。

\SourceParagraph{ch08-8.2.1-009}{598}
对每一对 $X,Y\in\operatorname{Ob}(\mathcal A)$，需要选择一个投影算子，即次数为零的链映射
\[
\Pi\colon\operatorname{Hom}_{\mathcal A}(X,Y)
\longrightarrow\operatorname{Hom}_{\mathcal A}(X,Y),
\qquad\Pi^2=\Pi,
\]
以及次数为 $-1$ 的链同伦
\[
H\colon\operatorname{Hom}_{\mathcal A}(X,Y)
\longrightarrow\operatorname{Hom}_{\mathcal A}(X,Y)[-1],
\]
满足
\[
\Pi-1=m_1H+Hm_1.
\]
这特别说明，$\Pi$ 的像是一个与 $\operatorname{Hom}_{\mathcal A}(X,Y)$ 同伦等价的复形。我们特别感兴趣的情形，是 $\Pi$ 的像实际上同构于 $H^*(\operatorname{Hom}_{\mathcal A}(X,Y))$；不过，无论是否如此，构造都适用。由此得到新的 $A_\infty$ 范畴 $\Pi\mathcal A$，其对象与态射为
\[
\operatorname{Ob}(\Pi\mathcal A)=\operatorname{Ob}(\mathcal A),
\qquad
\operatorname{Hom}_{\Pi\mathcal A}(X,Y)
=\Pi\bigl(\operatorname{Hom}_{\mathcal A}(X,Y)\bigr).
\]
接着作如下定义。
\begin{enumerate}
\item $m_1^{\Pi\mathcal A}=m_1^{\mathcal A}$。由于 $\Pi$ 是链映射，此定义有意义。
\item $m_2^{\Pi\mathcal A}=\Pi\circ m_2^{\mathcal A}$。
\item $m_d^{\Pi\mathcal A}=\sum_Tm_{d,T}$，其中 $T$ 遍历具有 $d$ 个输入和一个输出的所有带状树，树的边上不附加长度。将 $m_n$ 配置到每个价数为 $n+1$ 的内部顶点，将 $H$ 配置到每条内部边，将 $\Pi$ 配置到输出顶点，并将 $m_{d,T}$ 定义为由该图自然给出的逐次复合。例如，对\BookRef{fig:8.10}{图10}所示的树 $T$，得到
\[
m_{d,T}(a_5,\ldots,a_1)
=\Pi\bigl(m_2(a_5,H(m_3(a_4,a_3,H(m_2(a_2,a_1)))))\bigr).
\]
\end{enumerate}

\SourceParagraph{ch08-8.2.1-010}{598}
此外，还可以定义一个 $A_\infty$ 函子 $i\colon\Pi\mathcal A\to\mathcal A$。在对象层面，$i$ 是恒等映射；$i_1$ 就是链复形的包含映射，而 $i_d$ 的定义与 $m_d$ 类似，写成
\[
i_d=\sum_Ti_{d,T},
\]
其中 $T$ 遍历相同的树，将输出处的 $\Pi$ 换成 $H$，就得到 $i_{d,T}$ 的定义。

\begin{figure}[H]
\centering
\BookImage{01_Numbered_Figures/Chapter_08/Ch08_Fig10_PDFp599_Bookp589.pdf}
\BookNumberOnlyCaption\label{fig:8.10}
\end{figure}

\SourceParagraph{ch08-8.2.1-011}{599}
下面说明这一构造为何成立。首先考虑 $m_2$ 的关系，记住 $m_1^{\Pi\mathcal A}=m_1$：
\[
\begin{aligned}
&m_2^{\Pi\mathcal A}(a_2,m_1(a_1))
+(-1)^{\deg a_1-1}m_2^{\Pi\mathcal A}(m_1(a_2),a_1)
+m_1\bigl(m_2^{\Pi\mathcal A}(a_2,a_1)\bigr)\\
&\quad=\Pi\bigl(m_2(a_2,m_1(a_1))
+(-1)^{\deg a_1-1}m_2(m_1(a_2),a_1)\bigr)
+m_1\bigl(\Pi(m_2(a_2,a_1))\bigr)\\
&\quad=0.
\end{aligned}
\]
这里使用了 $m_2$ 的关系；在最后一项中，还使用了 $\Pi$ 与 $m_1$ 可交换以及 $\Pi^2=\Pi$。

\SourceParagraph{ch08-8.2.1-012}{599}
检验 $d\geq3$ 时的 $m_d^{\Pi\mathcal A}$ 满足 $A_\infty$ 关系，并检验 $\{i_d\}$ 定义了一个 $A_\infty$ 函子，是一个很好的练习。下面按照\BookCite{ref311}概述论证，将逐一检验符号这一相当费力的工作留给读者。

\SourceParagraph{ch08-8.2.1-013}{599}
定义新映射
\[
\begin{aligned}
\widehat m_d^{\Pi\mathcal A}\colon{}&
\operatorname{Hom}_{\Pi\mathcal A}(X_{d-1},X_d)
\otimes\cdots\otimes
\operatorname{Hom}_{\Pi\mathcal A}(X_0,X_1)\\
&\longrightarrow\operatorname{Hom}_{\Pi\mathcal A}(X_0,X_d)[3-d],
\end{aligned}
\]
其定义为
\[
\widehat m_d^{\Pi\mathcal A}(a_d,\ldots,a_1)
=\sum_{(T,e)}\pm\widehat m_{(T,e)}(a_d,\ldots,a_1),
\]
其中求和遍历上述所有树 $T$，以及在 $T$ 中选定的一条边 $e$。$\widehat m_{(T,e)}$ 的定义与以前相同，但若 $e$ 是内部边，就将该边上的标记 $H$ 换为 $Hm_1+m_1H=\Pi-1$；若 $e$ 是外部边，就在这条边上加上一个 $m_1$，见\BookRef{fig:8.11}{图11}。每棵树必须选择一个依赖于 $e$ 的符号；我们将确定使下述论证成立的符号这一任务留给读者。

\begin{figure}[H]
\centering
\BookImage{01_Numbered_Figures/Chapter_08/Ch08_Fig11_PDFp600_Bookp590.pdf}
\BookNumberOnlyCaption\label{fig:8.11}
\end{figure}

\SourceParagraph{ch08-8.2.1-014}{600}
现在用两种不同的方法计算 $\widehat m_d^{\Pi\mathcal A}$。首先使用表达式 $\Pi-1$，于是得到对 $\widehat m_d^{\Pi\mathcal A}$ 的三类不同贡献：
\[
\widehat m_d^{\Pi\mathcal A}
=m_d^{\Pi\mathcal A,m_1}
+m_d^{\Pi\mathcal A,\Pi}
-m_d^{\Pi\mathcal A,1}.
\]
这里，$m_d^{\Pi\mathcal A,m_1}$ 由 $e$ 为外部边的各对 $(T,e)$ 所给出的贡献组成。其余两项通过对 $e$ 为内部边的各对 $(T,e)$ 求和定义：在这两种情形中，分别将边上的 $H$ 换为 $\Pi$ 和 $1$。符号因子与 $\widehat m_d^{\Pi\mathcal A}$ 定义中的符号因子相同。

\SourceParagraph{ch08-8.2.1-015}{600}
另一方面，可以在每条内部边上使用表达式 $m_1H+Hm_1$。这样，$e$ 为内部边的每一对 $(T,e)$ 都产生两个贡献。现在考虑任意内部顶点 $v$。对每条以 $v$ 为端点的边，$\widehat m_d^{\Pi\mathcal A}$ 中都恰有一项在该边紧邻 $v$ 的位置带有一个 $m_1$，如\BookRef{fig:8.12}{图12}所示。注意，外部边上不配置 $H$。因此，对这样的顶点，可以应用 $m_n^{\mathcal A}$ 的 $A_\infty$ 关系，将这个和写成\BookRef{fig:8.13}{图13}所示形式的各项之和；这里的符号实际上由 $m_n$ 的 $A_\infty$ 关系中的符号决定。由此可见，$\widehat m_d^{\Pi\mathcal A}$ 可以写为 $m_d^{\Pi\mathcal A,1}$ 中出现的各棵树的贡献之和。仔细处理符号，实际上得到
\[
\widehat m_d^{\Pi\mathcal A}=-m_d^{\Pi\mathcal A,1}.
\]
于是
\[
m_d^{\Pi\mathcal A,m_1}+m_d^{\Pi\mathcal A,\Pi}=0.
\]
这恰好就是 $m_d^{\Pi\mathcal A}$ 的 $A_\infty$ 关系。

\SourceParagraph{ch08-8.2.1-016}{600}
\begin{exercise*}
请看看这一论证对 $m_3$ 如何进行！
\end{exercise*}

\begin{figure}[H]
\centering
\BookImage{01_Numbered_Figures/Chapter_08/Ch08_Fig12_PDFp601_Bookp591.pdf}
\BookNumberOnlyCaption\label{fig:8.12}
\end{figure}
\begin{figure}[H]
\centering
\BookImage{01_Numbered_Figures/Chapter_08/Ch08_Fig13_PDFp601_Bookp591.pdf}
\BookNumberOnlyCaption\label{fig:8.13}
\end{figure}

\SourceParagraph{ch08-8.2.1-017}{601}
证明 $\{i_d\}_{d\geq1}$ 构成一个 $A_\infty$ 函子，论证是类似的。先检验次数为一和二的情形，再概述一般情形。对于 $d=1$，只需
\[
m_1^{\mathcal A}(i_1(a_1))=i_1(m_1^{\Pi\mathcal A}(a_1)),
\]
而 $i_1$ 是包含映射，所以这显然成立。对于 $d=2$，注意 $i_d$ 的定义给出
\[
i_2(a_2,a_1)=Hm_2^{\mathcal A}(a_2,a_1).
\]
考虑表达式
\[
\begin{aligned}
&Hm_2^{\mathcal A}(a_2,m_1(a_1))
+(-1)^{\deg a_1-1}Hm_2^{\mathcal A}(m_1(a_2),a_1)\\
&\hspace{5em}+Hm_1(m_2^{\mathcal A}(a_2,a_1)).
\end{aligned}
\]
利用 $m_2^{\mathcal A}$ 的关系，知道它实际上为零。不过，$Hm_1=-m_1H+\Pi-1$，因此，在前三项中使用 $i_2$ 的定义，在第四项中使用 $m_2^{\Pi\mathcal A}$ 的定义，就得到
\[
\begin{aligned}
0={}&i_2(a_2,m_1^{\Pi\mathcal A}(a_1))
+(-1)^{\deg a_1-1}i_2(m_1^{\Pi\mathcal A}(a_2),a_1)\\
&-m_1^{\mathcal A}(i_2(a_2,a_1))
+i_1m_2^{\Pi\mathcal A}(a_2,a_1)
-m_2^{\mathcal A}(i_1(a_2),i_1(a_1)),
\end{aligned}
\]
这就化为\BookEqRef{eq:8.8}{8.8}。

\SourceParagraph{ch08-8.2.1-018}{602}
一般情形的处理与前面基本相同。定义
\[
\widehat i_d=\sum_{(T,e)}\pm\widehat i_{(T,e)},
\]
求和仍遍历具有 $d$ 个输入的树，以及在树中选定的一条边；$\widehat i_{(T,e)}$ 的定义与 $\widehat m_{(T,e)}^{\Pi\mathcal A}$ 相同，只是将输出顶点所配置的算子 $\Pi$ 换成 $H$。符号的选择与 $\widehat m_d^{\Pi\mathcal A}$ 定义中的选择相同。在内部边上配置 $\Pi-1$，就可以写成
\[
\widehat i_d=i_d^{m_1}+i_d^\Pi-i_d^1,
\]
其中 $i_d^{m_1}$、$i_d^\Pi$ 和 $i_d^1$ 的定义与 $m_d$ 论证中完全相同，只是输出顶点上的 $\Pi$ 仍换成 $H$。若用 $m_1H+Hm_1$ 代替 $\Pi-1$ 来计算 $\widehat i_d$，就和前面一样得到
\[
\widehat i_d=-i_d^1,
\]
从而
\[
i_d^{m_1}+i_d^\Pi=0.
\]
现在可以看出，$i_d^\Pi$ 产生如下带有适当符号的和：
\[
\sum_{\substack{1<p<d\\0\leq q\leq d-p}}
\pm i_{d-p+1}\bigl(a_d,\ldots,a_{p+q+1},
 m_p^{\Pi\mathcal A}(a_{p+q},\ldots,a_{q+1}),
 a_q,\ldots,a_1\bigr).
\]
注意，这个和不包括 $p=1,d$ 的情形。$i_d^{m_1}$ 中选定的边为输入边的各项，给出 $p=1$ 的项：
\[
\sum_{\substack{p=1\\0\leq q\leq d-p}}
\pm i_{d-p+1}\bigl(a_d,\ldots,a_{q+2},
 m_1^{\Pi\mathcal A}(a_{q+1}),a_q,\ldots,a_1\bigr).
\]

\SourceParagraph{ch08-8.2.1-019}{602}
为得到余下的各项，考虑 $i_d^{m_1}$ 中选定的边为输出边的各项；它们给出形如 $Hm_1(\cdots)$ 的贡献。利用 $Hm_1=-m_1H+\Pi-1$，可见这些项中来自 $-m_1H$ 的贡献为 $-m_1^{\mathcal A}(i_d(a_d,\ldots,a_1))$，来自 $\Pi$ 的贡献为 $i_1m_d^{\Pi\mathcal A}(a_d,\ldots,a_1)$，而来自 $-1$ 的贡献稍微复杂一些：\BookRef{fig:8.14}{图14}所示的树贡献
\[
-m_r^{\mathcal A}\bigl(i_{s_r}(a_d,\ldots,a_{d-s_r+1}),
\ldots,i_{s_1}(a_{s_1},\ldots,a_1)\bigr).
\]
这涵盖了 $i_d$ 的关系中的全部项；再次仔细处理符号，就能得到 $i_d$ 的关系。

\begin{figure}[H]
\centering
\BookImage{01_Numbered_Figures/Chapter_08/Ch08_Fig14_PDFp603_Bookp593.pdf}
\BookNumberOnlyCaption\label{fig:8.14}
\end{figure}

\SourceParagraph{ch08-8.2.1-020}{603}
显然，$\{i_d\}_{d\geq1}$ 是 $A_\infty$ 范畴的拟同构，因为这只是对 $i_1$ 的一个条件，而链同伦 $H$ 的存在保证该条件成立。

\SourceParagraph{ch08-8.2.1-021}{603}
现在回到范畴 $D^b_\infty(X)$。上述构造的意义，在于能够给 $H^*(D^b_\infty(X))$ 配置一个 $A_\infty$ 结构。具体做法是：对每一对 $\mathcal E^\bullet,\mathcal F^\bullet$，选择复形 $\operatorname{Ext}^\bullet(\mathcal E^\bullet,\mathcal F^\bullet)$ 到 $\operatorname{Hom}^\bullet_{D^b_\infty(X)}(\mathcal E^\bullet,\mathcal F^\bullet)$ 中的包含映射 $i$，使其为拟同构；前一个复形在次数 $i$ 处为 $H^i(\operatorname{Hom}^\bullet_{D^b_\infty(X)}(\mathcal E^\bullet,\mathcal F^\bullet))$，边界映射为零。还要选择从 $\operatorname{Hom}^\bullet_{D^b_\infty(X)}(\mathcal E^\bullet,\mathcal F^\bullet)$ 投影到 $i(\operatorname{Ext}^\bullet(\mathcal E^\bullet,\mathcal F^\bullet))$ 的投影算子 $\Pi$ 和同伦 $H$。然后利用上述构造，在 $H^*(D^b_\infty(X))$ 上得到一个 $A_\infty$ 结构。这个分次范畴的零次上同调范畴就是通常的 $D^b(X)$，其各次态射则为导出范畴中的相应 $\operatorname{Ext}$ 群，所以这给出了 $D^b(X)$ 的一个 $A_\infty$ 结构。微分 $m_1$ 为零，但高阶乘法映射并不为零。此外，$D^b(X)$ 上的这一 $A_\infty$ 结构，给出一个与 $D^b_\infty(X)$ 拟同构的 $A_\infty$ 范畴。因此，尽管最初定义的高阶乘法映射为零，仍应将 $D^b_\infty(X)$ 上的 $A_\infty$ 结构看成一个极不平凡的结构。

\SourceParagraph{ch08-8.2.1-022}{603}
\BookStatementNumber{8}
\begin{example}\label{ex:8.8}
这里给出 $D^b(\mathbb P^1)$ 中非平凡 $m_3$ 的一个尽可能简单的例子，它出现于
\[
\begin{aligned}
m_3\colon{}&\operatorname{Ext}^1(\mathcal O_{\mathbb P^1}(2),\mathcal O_{\mathbb P^1})
\otimes\operatorname{Hom}(\mathcal O_{\mathbb P^1}(1),\mathcal O_{\mathbb P^1}(2))\\
&\otimes\operatorname{Hom}(\mathcal O_{\mathbb P^1},\mathcal O_{\mathbb P^1}(1))
\longrightarrow\operatorname{Hom}(\mathcal O_{\mathbb P^1},\mathcal O_{\mathbb P^1}).
\end{aligned}
\]

\SourceParagraph{ch08-8.2.1-023}{604}
为进行计算，需要下列层的 Čech 消解：
\[
\begin{aligned}
\mathcal{H}om(\mathcal O_{\mathbb P^1},\mathcal O_{\mathbb P^1}(1))
&=\mathcal{H}om(\mathcal O_{\mathbb P^1}(1),\mathcal O_{\mathbb P^1}(2))
=\mathcal O_{\mathbb P^1}(1),\\
\mathcal{H}om(\mathcal O_{\mathbb P^1},\mathcal O_{\mathbb P^1}(2))
&=\mathcal O_{\mathbb P^1}(2),\\
\mathcal{H}om(\mathcal O_{\mathbb P^1}(2),\mathcal O_{\mathbb P^1})
&=\mathcal O_{\mathbb P^1}(-2),\\
\mathcal{H}om(\mathcal O_{\mathbb P^1}(1),\mathcal O_{\mathbb P^1})
&=\mathcal O_{\mathbb P^1}(-1),\\
\mathcal{H}om(\mathcal O_{\mathbb P^1},\mathcal O_{\mathbb P^1})
&=\mathcal O_{\mathbb P^1}.
\end{aligned}
\]

\SourceParagraph{ch08-8.2.1-024}{604}
采用 $\mathbb P^1$ 的标准开覆盖 $U_0,U_1$，其中 $U_0=\mathbb C$，$U_1=(\mathbb C\setminus\{0\})\cup\{\infty\}$，$U_0$ 和 $U_1$ 上的坐标分别为 $z,z^{-1}$。下面给出这五个层的 Čech 复形，以及在需要时所选择的 $H$ 和 $\Pi$。在这种情形中，$H$ 仅在从次数一到次数零时非零。
\begin{table}[H]
\centering
\begingroup\small\setlength{\tabcolsep}{4pt}
\begin{tabular}{ll}
\toprule
$\mathcal O_{\mathbb P^1}(1)$
& $\mathbb C[z]\oplus z\mathbb C[z^{-1}]\longrightarrow\mathbb C[z,z^{-1}]$\\
$m_1$ & $(f_0,f_1)\longmapsto f_1-f_0$\\
\midrule
$\mathcal O_{\mathbb P^1}(2)$
& $\mathbb C[z]\oplus z^2\mathbb C[z^{-1}]\longrightarrow\mathbb C[z,z^{-1}]$\\
$m_1$ & $(f_0,f_1)\longmapsto f_1-f_0$\\
$\Pi_0$ & $(a_0+a_1z+a_2z^2+z^3f_0(z),z^2f_1(z^{-1}))$\\
& $\qquad\longmapsto(a_0+a_1z+a_2z^2,a_0+a_1z+a_2z^2)$\\
$\Pi_1$ & $0$\\
$H$ & $(z^3f_0(z)+z^2f_1(z^{-1}))\longmapsto(z^3f_0(z),-z^2f_1(z^{-1}))$\\
\midrule
$\mathcal O_{\mathbb P^1}(-2)$
& $\mathbb C[z]\oplus z^{-2}\mathbb C[z^{-1}]\longrightarrow\mathbb C[z,z^{-1}]$\\
$m_1$ & $(f_0,f_1)\longmapsto f_1-f_0$\\
\midrule
$\mathcal O_{\mathbb P^1}(-1)$
& $\mathbb C[z]\oplus z^{-1}\mathbb C[z^{-1}]\longrightarrow\mathbb C[z,z^{-1}]$\\
$m_1$ & $(f_0,f_1)\longmapsto f_1-f_0$\\
$\Pi$ & $0$\\
$H$ & $(f_0(z)+z^{-1}f_1(z^{-1}))\longmapsto(f_0(z),-z^{-1}f_1(z^{-1}))$\\
\midrule
$\mathcal O_{\mathbb P^1}$
& $\mathbb C[z]\oplus\mathbb C[z^{-1}]\longrightarrow\mathbb C[z,z^{-1}]$\\
$m_1$ & $(f_0,f_1)\longmapsto f_1-f_0$\\
$\Pi_0$ & $(a_0+zf_0(z),f_1(z^{-1}))\longmapsto(a_0,a_0)$\\
$\Pi_1$ & $0$\\
$H$ & $(zf_0(z)+f_1(z^{-1}))\longmapsto(zf_0(z),-f_1(z^{-1}))$\\
\bottomrule
\end{tabular}
\endgroup
\end{table}

\SourceParagraph{ch08-8.2.1-025}{604}
对 $m_3$ 有贡献的项只有三个，见\BookRef{fig:8.15}{图15}；不过，第一项无需考虑，因为在 Čech 复形层面没有 $m_3$。此外，
\[
\Pi m_2(a_3,Hm_2(a_2,a_1))=0,
\]
因为对 $a_2\in\operatorname{Hom}(\mathcal O_{\mathbb P^1}(1),\mathcal O_{\mathbb P^1}(2))$、$a_1\in\operatorname{Hom}(\mathcal O_{\mathbb P^1},\mathcal O_{\mathbb P^1}(1))$，$m_2(a_2,a_1)$ 的次数为零，而 $H$ 只在次数一上起作用。对于 $\Pi m_2(Hm_2(a_3,a_2),a_1)$，得到下表。
\begin{table}[H]
\centering
\begingroup\footnotesize\setlength{\tabcolsep}{3pt}
\begin{tabular}{cccccc}
\toprule
$a_3$ & $a_2$ & $a_1$ & $m_2(a_3,a_2)$
& $Hm_2(a_3,a_2)$ & $\Pi m_2(Hm_2(a_3,a_2),a_1)$\\
\midrule
$z^{-1}$ & $(1,1)$ & $(1,1)$ & $z^{-1}$ & $(0,-z^{-1})$ & $(0,0)$\\
$z^{-1}$ & $(z,z)$ & $(1,1)$ & $1$ & $(1,0)$ & $(1,1)$\\
$z^{-1}$ & $(1,1)$ & $(z,z)$ & $z^{-1}$ & $(0,-z^{-1})$ & $(0,0)$\\
$z^{-1}$ & $(z,z)$ & $(z,z)$ & $1$ & $(1,0)$ & $(0,0)$\\
\bottomrule
\end{tabular}
\endgroup
\end{table}

\SourceParagraph{ch08-8.2.1-026}{605}
最后一列描述了 $m_3$；特别地，它不为零。

\begin{figure}[H]
\centering
\BookImage{01_Numbered_Figures/Chapter_08/Ch08_Fig15_PDFp605_Bookp595.pdf}
\BookNumberOnlyCaption\label{fig:8.15}
\end{figure}
\end{example}

\section{Morse 与 Fukaya \texorpdfstring{$A_\infty$}{A∞} 范畴}\label{sec:8.3}

\SourceParagraph{ch08-8.3-intro-001}{605}
现在我们将逐步引出 Fukaya 范畴的定义。它的特点与此前考虑的范畴有着根本差别。作为准备，我们先给出一个较简单的例子，即 Morse 范畴。它与 Fukaya 范畴具有相似的特点，并且在某些情形下与后者存在联系。

\subsection{Morse 范畴}\label{sec:8.3.1}

\SourceParagraph{ch08-8.3.1-001}{605}
我们先回顾 Morse 理论的基本内容（例如，可参见 \BookBibText{ref242}{MS1} 第 3.4 节和第 10.5 节）。本小节中，$M$ 始终表示一个 $n$ 维紧致定向 Riemann 流形。若可微函数 $f\colon M\to\mathbb R$ 的临界点 $p$ 附近存在局部坐标，使得
\[
f=f(p)+x_1^2+\cdots+x_k^2-x_{k+1}^2-\cdots-x_n^2,
\]
则称该临界点\emph{非退化}。此时，我们称它的指标为 $k$，并记
\[
\deg p=k.
\]

\SourceParagraph{ch08-8.3.1-002}{605}
对于这样的临界点，有两个流形 $W^s(p)$ 和 $W^u(p)$。其中，$W^s(p)$ 为稳定流形，由在 $f$ 的\emph{负}梯度流下、当 $t\to+\infty$ 时趋于 $p$ 的点组成；$W^u(p)$ 为不稳定流形，由在 $f$ 的\emph{正}梯度流下、当 $t\to+\infty$ 时趋于 $p$ 的点组成。由临界点的局部描述，很容易验证
\[
\deg p=\dim W^s(p).
\]

\BookStatementNumber{9}
\begin{definition}\label{def:8.9}
\SourceParagraph{ch08-8.3.1-003}{605}
若 $f$ 只有有限个临界点，这些临界点均非退化，并且所有稳定子流形与不稳定子流形都彼此横截相交，则称 $f$ 为 \emph{Morse--Smale 函数}。
\end{definition}

\SourceParagraph{ch08-8.3.1-004}{606}
对于任意临界点 $p$，$W^s(p)$ 并不自带一个内在的定向。不过，我们可以考虑由 $W^s(p)$ 的两个可能定向生成的自由 Abel 群，再对关系
\[
[\Omega_p]+[-\Omega_p]=0
\]
取商，所得的群记为 $o_p$。这里，$[\Omega_p]$ 和 $[-\Omega_p]$ 表示 $W^s(p)$ 的两个相反定向。作为抽象群，有 $o_p\cong\mathbb Z$。

\SourceParagraph{ch08-8.3.1-005}{606}
现在可以定义 Morse 复形
\[
CM^i(f)=\bigoplus_{\substack{p\in\operatorname{Crit}(f)\\\deg(p)=i}}o_p.
\]
我们按照如下方式定义微分，并将它记为 $m_1$。实质上，$m_1$ 计数连接两个次数相差 $1$ 的临界点的梯度流线。具体来说，首先注意，$M$ 已定向，且 $W^s(p)$ 与 $W^u(p)$ 在 $p$ 处横截相交，维数互补，因此 $W^s(p)$ 的一个定向会诱导 $W^u(p)$ 的定向。若 $\deg q=\deg p+1$，则
\[
\dim W^u(p)=n-\deg p,\qquad \dim W^s(q)=\deg p+1,
\]
所以 $W^u(p)\cap W^s(q)$ 为一维，因而是一些线段的并：这些线段就是从 $p$ 到 $q$ 的 $-f$ 的梯度流轨道。此外，$W^u(p)$ 和 $W^s(q)$ 的定向决定了 $W^u(p)\cap W^s(q)$ 的定向。映射 $f$ 将这些线段与 $\mathbb R$ 中的线段对应起来，而 $\mathbb R$ 具有自然定向。我们令 $n(q,p)$ 等于 $W^u(p)\cap W^s(q)$ 中定向与 $\mathbb R$ 的定向一致的线段数，减去定向相反的线段数。由此定义
\[
m_1\colon CM^i(f)\longrightarrow CM^{i+1}(f),
\qquad
m_1([\Omega_p])=\sum_q n(q,p)[\Omega_q].
\]
注意，改变 $\Omega_p$ 或 $\Omega_q$ 的定向，会改变 $W^u(p)\cap W^s(q)$ 的定向符号，因而改变 $n(q,p)$ 的符号。因此 $m_1$ 定义良好。

\SourceParagraph{ch08-8.3.1-006}{606}
\BookBibText{ref242}{MS1} 第 10.5.4 节已经证明，映射 $m_1$ 满足 $m_1^2=0$。我们在这里的目标是定义具有高阶乘法映射的 $A_\infty$ 范畴，而 $m_1^2=0$ 将被包含在 $A_\infty$ 关系之中。

\SourceParagraph{ch08-8.3.1-007}{606}
现在来定义 \emph{Morse 范畴} $\operatorname{Morse}(M)$，其对象为 $M$ 上的所有光滑实值函数。Fukaya 在\BookCite{ref165}中描述过这一构造，不过这里遵循 Abouzaid 在\BookCite{ref002}中的阐述。

\SourceParagraph{ch08-8.3.1-008}{606}
$\operatorname{Morse}(M)$ 的对象为 $M$ 上的所有光滑实值函数。但是，态射空间和高阶乘法映射并非总是存在。因此，我们所定义的甚至还不是一个 $A_\infty$ 范畴，而是一个称为 $A_\infty$ \emph{预范畴}的结构；下一小节将更精确地说明它的含义。

\SourceParagraph{ch08-8.3.1-009}{607}
若两个函数 $f_0,f_1$ 满足 $f_1-f_0$ 是 $M$ 上的 Morse--Smale 函数，则令
\[
\operatorname{Hom}(f_0,f_1)=CM^\bullet(f_1-f_0),
\]
其中 $m_1$ 如上定义。给定满足某些附加条件的 $\{f_0,\ldots,f_d\}$——这些条件稍后会准确说明——我们将定义高阶乘法映射
\[
m_d\colon\operatorname{Hom}(f_{d-1},f_d)\otimes\cdots\otimes
\operatorname{Hom}(f_0,f_1)\longrightarrow\operatorname{Hom}(f_0,f_d).
\]
为此，进行如下构造。

\BookStatementNumber{10}
\begin{definition}\label{def:8.10}
\SourceParagraph{ch08-8.3.1-010}{607}
给定一列函数 $f_0,\ldots,f_d$，以及 $f_{i+1}-f_i$ 的临界点 $p_{i,i+1}$ 和 $f_d-f_0$ 的临界点 $p_{0,d}$，\emph{梯度树模空间}
\[
\mathcal S_d(p_{0,d};p_{0,1},\ldots,p_{d-1,d})
\]
定义为满足以下条件的映射 $\phi\colon S\to M$ 所组成的空间，其中 $S\in\mathcal S_d$，后者为带状树模空间。
\begin{enumerate}
\item 按如下方式为 $S$ 的每条边 $e$ 标记一个函数 $f_e$。若 $e$ 是标号为 $1\leq i\leq d$ 的外部输入边，则 $f_e=f_i-f_{i-1}$；否则，若 $e$ 从顶点 $v$ 输出，则 $f_e$ 等于所有输入 $v$ 的边所标记的函数之和。每条边 $e$ 的像均须为 $f_e$ 的一条梯度流线。
\item $e$ 的定向由 $-f_e$ 的梯度流给出。此外，若 $e$ 的长度为 $l$，则 $\phi$ 将由边的一个固定参数化 $[0,l]$ 给出的向量场，与 $-f_e$ 的梯度向量场对应起来。
\item 若 $e$ 是外部边，且 $f_e=f_i-f_j$，则 $\phi$ 的像在外部顶点处趋于 $p_{j,i}$。
\end{enumerate}
\end{definition}

\SourceParagraph{ch08-8.3.1-011}{607}
对于给定的函数列 $f_0,\ldots,f_d$，可以令
\[
\mathcal S(f_0,\ldots,f_d)=
\coprod_{p_{0,d},p_{0,1},\ldots,p_{d-1,d}}
\mathcal S(p_{0,d};p_{0,1},\ldots,p_{d-1,d}),
\]
这里的不交并遍历所有可能的临界点。

\SourceParagraph{ch08-8.3.1-012}{607}
允许边的长度趋于 $\infty$，便可将此模空间紧化；此时，一棵梯度树趋于若干梯度树的并，如\BookRef{fig:8.16}{图16}所示。这样得到紧化 $\overline{\mathcal S}(f_0,\ldots,f_d)$。

\setcounter{figure}{15}
\begin{figure}[H]
\centering
\BookImage{01_Numbered_Figures/Chapter_08/Ch08_Fig16_PDFp608_Bookp598.pdf}
\BookNumberOnlyCaption\label{fig:8.16}
\end{figure}

\SourceParagraph{ch08-8.3.1-013}{607}
下面说明如何构造 $\mathcal S(p_{0,d};p_{0,1},\ldots,p_{d-1,d})$。记 $p_i:=p_{i-1,i}$、$q=p_{0,d}$，并令
\[
\mathcal E(q;p_1,\ldots,p_d)=W^s_{f_d-f_0}(q)
\times\left(\prod_{i=1}^dW^u_{f_i-f_{i-1}}(p_i)\right)\times\mathcal S_d,
\]
其中 $W$ 的下标表示我们针对哪个函数取稳定流形或不稳定流形。对每个 $S\in\mathcal S_d$，在输出边上选取一个基点 $s\in S$，使它与最后一个内部顶点相距某个固定距离。于是，对于给定的 $S$，可以按如下方式定义映射
\[
\Phi_j^s\colon W^u_{f_j-f_{j-1}}(p_j)\longrightarrow M.
\]

\SourceParagraph{ch08-8.3.1-014}{608}
令 $v_j$ 为第 $j$ 条输入边所接到的内部顶点，$A$ 为 $S$ 中从 $v_j$ 到 $s$ 的路径。对于点 $y\in W^u_{f_j-f_{j-1}}(p_j)$，依次施加 $A$ 所经过的各条边上函数的梯度流，取 $y$ 在这组梯度流下的像；其中，$-f_e$ 的梯度流所经历的时间等于 $e\cap A$ 的长度。将这些映射与包含映射 $W^s_{f_d-f_0}(q)\hookrightarrow M$ 合在一起，得到具有 $d+1$ 个分量的映射
\[
\Phi\colon\mathcal E(q;p_1,\ldots,p_d)\longrightarrow M^{d+1}.
\]

\SourceParagraph{ch08-8.3.1-015}{608}
此时很明显，$\mathcal S(q;p_1,\ldots,p_d)$ 正是对角线 $\Delta\subseteq M^{d+1}$ 的原像。若每个 $f_j-f_i$ 在 $i<j$ 时均为 Morse--Smale 函数，且 $\Phi$ 与 $\Delta$ 横截，即 $\Phi^{-1}(\Delta)$ 是具有预期维数的光滑流形，则称函数组 $f_0,\ldots,f_d$ 为 \emph{Morse--Smale 函数组}。注意，这一预期维数为
\[
\begin{aligned}
&\dim\mathcal E(q;p_1,\ldots,p_d)-\operatorname{codim}_{M^{d+1}}\Delta\\
&\quad=\deg q+d\dim M-\sum_i\deg p_i+d-2-d\dim M\\
&\quad=\deg q+d-2-\sum_i\deg p_i.
\end{aligned}
\]
因此，特别地，当 $2-d+\sum_i\deg p_i=\deg q$ 时，$\overline{\mathcal S}(f_0,\ldots,f_d)$ 是零维的。

\SourceParagraph{ch08-8.3.1-016}{608}
这里有一个微妙之处：我们需要为 $\overline{\mathcal S}(f_0,\ldots,f_d)$ 赋予定向。这一步必须格外谨慎，才能使 $A_\infty$ 关系在正确的符号约定下成立。只有\BookCite{ref002}附录 B 详细完成了这项工作。这里略去这些细节。

\SourceParagraph{ch08-8.3.1-017}{609}
当 $2-d+\sum_i\deg p_i=\deg q$ 时，令 $n(q;p_1,\ldots,p_d)$ 为 $\mathcal S(q;p_1,\ldots,p_d)$ 中梯度树的带符号数目，然后用公式
\[
m_d([\Omega_{p_d}]\otimes\cdots\otimes[\Omega_{p_1}])
=\sum_q n(q;p_1,\ldots,p_d)[\Omega_q]
\]
定义 $m_d$。当然，定向 $\Omega_{p_1},\ldots,\Omega_{p_d},\Omega_q$ 会影响模空间 $\mathcal S(q;p_1,\ldots,p_d)$ 的定向，从而使此定义良好。

\SourceParagraph{ch08-8.3.1-018}{609}
还需要说明，如此定义的 $m_d$ 为什么满足 $A_\infty$ 关系。为此，考虑满足
\[
3-d+\sum_{i=1}^d\deg p_i=\deg q
\]
时的紧化 $\overline{\mathcal S}(q;p_1,\ldots,p_d)$。由这一次数条件可知，模空间是一维的。其边界由形如
\[
\overline{\mathcal S}(q;p_1,\ldots,p_j,r,p_{i+j+1},\ldots,p_d)
\times\overline{\mathcal S}(r;p_{j+1},\ldots,p_{i+j})
\]
的部分组成，其中 $1\leq i\leq d$、$0\leq j\leq d-i$，并且
\[
\deg r=2-i+\sum_{k=j+1}^{i+j}\deg p_k.
\]

\SourceParagraph{ch08-8.3.1-019}{609}
现在考虑 $d$ 次 $A_\infty$ 关系中出现的一项：
\[
\begin{aligned}
&m_{d-i+1}\bigl([\Omega_{p_d}],\ldots,[\Omega_{p_{i+j+1}}],
m_i([\Omega_{p_{i+j}}],\ldots,[\Omega_{p_{j+1}}]),
[\Omega_{p_j}],\ldots,[\Omega_{p_1}]\bigr)\\
&\quad=\sum_r n(r;p_{j+1},\ldots,p_{i+j})\,
m_{d-i+1}\bigl([\Omega_{p_d}],\ldots,[\Omega_{p_{i+j+1}}],
[\Omega_r],[\Omega_{p_j}],\ldots,[\Omega_{p_1}]\bigr)\\
&\quad=\sum_{q,r}n(r;p_{j+1},\ldots,p_{i+j})\,
n(q;p_1,\ldots,p_j,r,p_{i+j+1},\ldots,p_d)[\Omega_q].
\end{aligned}
\]

\SourceParagraph{ch08-8.3.1-020}{609}
因此，临界点 $p_1,\ldots,p_d$ 对 $A_\infty$ 关系中 $[\Omega_q]$ 的系数所作的贡献，恰好是 $\overline{\mathcal S}(q;p_1,\ldots,p_d)$ 的边界点数的带符号和。然而，一维带边界流形的边界点数的带符号和总是零，所以只要所有定向选择彼此一致，$A_\infty$ 方程就成立。\BookCite{ref002}描述了这种定向的选择。

\subsection{\texorpdfstring{$A_\infty$}{A∞} 预范畴}\label{sec:8.3.2}

\SourceParagraph{ch08-8.3.2-001}{609}
上一小节的构造存在一个缺点：并非任意两个对象之间都定义了态射空间，而且高阶乘法还受到进一步的限制。因此，我们甚至尚未定义一个真正的 $A_\infty$ 范畴，而只是得到了 Kontsevich 与 Soibelman 在\BookCite{ref311}中称为 $A_\infty$ \emph{预范畴}的结构。下面给出定义。

\BookStatementNumber{11}
\begin{definition}\label{def:8.11}
\SourceParagraph{ch08-8.3.2-002}{609}
一个（非幺）$A_\infty$ 预范畴 $\mathcal A$ 包含以下数据：
\begin{enumerate}
\item 一个对象类 $\operatorname{Ob}(\mathcal A)$。
\item 对每个 $n\geq2$，给定一族\emph{横截序列}
$\operatorname{Ob}_{\mathrm{tr}}^n(\mathcal A)\subseteq
\operatorname{Ob}(\mathcal A)^n$，即由对象的 $n$ 元序列组成的一个集合。
\item 对每个 $(X_0,X_1)\in\operatorname{Ob}_{\mathrm{tr}}^2(\mathcal A)$，给定一个 $\mathbb Z$ 分次链复形 $\operatorname{Hom}_{\mathcal A}(X_0,X_1)$。
\item 对每个 $(X_0,\ldots,X_d)\in\operatorname{Ob}_{\mathrm{tr}}^{d+1}(\mathcal A)$，给定映射
\[
\begin{aligned}
m_d\colon\operatorname{Hom}_{\mathcal A}(X_{d-1},X_d)\otimes\cdots
\otimes\operatorname{Hom}_{\mathcal A}(X_0,X_1)
\longrightarrow\operatorname{Hom}_{\mathcal A}(X_0,X_d)[2-d].
\end{aligned}
\]
此外，还要求：
\item 横截序列的每个子序列都是横截的。
\item 各映射 $m_d$ 满足 $A_\infty$ 关系。
\end{enumerate}
\end{definition}

\SourceParagraph{ch08-8.3.2-003}{610}
因此，$\operatorname{Morse}(M)$ 是一个 $A_\infty$ 预范畴的例子：它的横截序列为那些自身以及每个子序列都满足 Morse--Smale 条件的序列 $(f_0,\ldots,f_d)$。

\SourceParagraph{ch08-8.3.2-004}{610}
我们还需要 $A_\infty$ 预范畴之间的拟等价概念。首先定义 $A_\infty$ 预范畴之间的函子。

\BookStatementNumber{12}
\begin{definition}\label{def:8.12}
\SourceParagraph{ch08-8.3.2-005}{610}
设 $\mathcal A,\mathcal B$ 为 $A_\infty$ 预范畴。一个 $A_\infty$ 函子 $F\colon\mathcal A\to\mathcal B$ 包含：
\begin{enumerate}
\item 映射 $F\colon\operatorname{Ob}(\mathcal A)\to\operatorname{Ob}(\mathcal B)$，满足只要 $(X_1,\ldots,X_d)\in\operatorname{Ob}_{\mathrm{tr}}^d(\mathcal A)$，就有
\[
(F(X_1),\ldots,F(X_d))\in\operatorname{Ob}_{\mathrm{tr}}^d(\mathcal B).
\]
\item 对 $\mathcal A$ 的每个横截序列 $(X_0,\ldots,X_d)$，给定次数为 $1-d$ 的映射
\[
\begin{aligned}
F_d\colon\operatorname{Hom}_{\mathcal A}(X_{d-1},X_d)\otimes\cdots
\otimes\operatorname{Hom}_{\mathcal A}(X_0,X_1)
\longrightarrow\operatorname{Hom}_{\mathcal B}(F(X_0),F(X_d)),
\end{aligned}
\]
这些映射满足 $A_\infty$ 函子的通常公式。
\end{enumerate}
\end{definition}

\BookStatementNumber{13}
\begin{definition}\label{def:8.13}
\SourceParagraph{ch08-8.3.2-006}{610}
若 $A_\infty$ 预范畴 $\mathcal A,\mathcal B$ 之间的 $A_\infty$ 函子在 $H^*(\mathcal A)$、$H^*(\mathcal B)$ 层次诱导同构，则称它为\emph{拟等价}。
\end{definition}

\SourceParagraph{ch08-8.3.2-007}{610}
在许多情形下，例如即将在\BookSectionRef{sec:8.4}{8.4}节看到的情形，避开对象之间不横截时的定义通常比较方便；这一形式体系使我们能够这样处理。

\subsection{Fukaya 范畴}\label{sec:8.3.3}

\SourceParagraph{ch08-8.3.3-001}{610}
现在描述 Calabi--Yau 流形 $X$ 的 Fukaya 范畴，其中 $X$ 配备一个辛形式 $\omega$ 和一个处处非零的全纯 $n$ 形式 $\Omega$。我们将比本书前面的介绍（见\BookSectionRef{sec:3.6}{3.6}节）以及 \BookBibText{ref242}{MS1} 中的介绍稍微详细一些。不过，要尝试给出在技术上完全正确的定义，远远超出了这里所能完成的范围。有关更多技术细节，可查阅 Seidel 的书\BookCite{ref421}；这里紧密遵循其中的处理。

\SourceParagraph{ch08-8.3.3-002}{610}
首先描述 Fukaya 范畴 $\operatorname{Fuk}(X)$ 的对象。它们是四元组 $(L,\mathcal L,\widetilde\xi_L,S_L)$，其中：
\begin{itemize}
\item $L\subseteq X$ 是一个定向 Lagrange 子流形。
\item $\mathcal L$ 是 $L$ 上具有平坦联络的 $U(1)$ 丛。
\item 总存在一个函数
\[
\xi_L\colon L\longrightarrow S^1,\qquad
\xi_L=\frac{\Omega|_L}{|\Omega|_L|}.
\]
换言之，在任意点 $x\in L$ 处，选取 $L$ 的切向量的一组正定向基 $\eta_1,\ldots,\eta_n$，则
\[
\xi_L(x)=\frac{\Omega(\eta_1,\ldots,\eta_n)}
{|\Omega(\eta_1,\ldots,\eta_n)|}.
\]
很容易验证，$\xi_L$ 与基的选择无关。$\widetilde\xi_L$ 是 $L$ 的一个\emph{分次}，即一个提升
\[
\widetilde\xi_L\colon L\longrightarrow\mathbb R,
\qquad \xi_L=e^{\pi\mathrm i\widetilde\xi_L}.
\]

\SourceParagraph{ch08-8.3.3-003}{611}
注意，并非所有 Lagrange 子流形都具有分次。函数 $\xi_L$ 决定了 $H^1(L,\mathbb Z)$ 中的一个类，即沿 $\xi_L$ 拉回 $H^1(S^1,\mathbb Z)$ 的生成元所得的类，称为 $L$ 的 \emph{Maslov 类}，见\BookRef{fig:8.17}{图17}。只有当这个类为零时，分次才存在；分次并不唯一，可以加上任意偶整数。注意，改变 $L$ 的定向需要将分次平移 $1$，这可以理解为把膜变为反膜。若 $L$ 实际上是特殊 Lagrange 子流形，即对于某个固定的角 $\theta$，有 $\Omega|_L=e^{\mathrm i\theta}\operatorname{Vol}(L)$，则 $\xi_L$ 是常函数 $e^{\mathrm i\theta}$。因此，特殊 Lagrange 子流形总是具有分次。

\SourceParagraph{ch08-8.3.3-004}{611}
\item $S_L$ 是 $L$ 上的一个自旋结构。$L$ 上存在自旋结构的障碍是 $L$ 的第二 Stiefel--Whitney 类 $w_2\in H^2(L,\mathbb Z/2\mathbb Z)$；若 $w_2=0$，则自旋结构由 $H^1(L,\mathbb Z/2\mathbb Z)$ 分类。这一选择在我们的阐述中不会起很大的作用，但对于确定那些用来定义映射 $m_k$ 的 $J$ 全纯圆盘模空间的定向，它实际上很重要。由于这里将略去对这些符号的大部分讨论，读者不必过于担心这一问题。
\end{itemize}

\SourceParagraph{ch08-8.3.3-005}{611}
事实上，还需要对 Lagrange 子流形施加进一步的限制：$L$ 应当在\BookCite{ref164}的意义下\emph{无障碍}。这是一个极其精细而复杂的条件，涉及边界位于 $L$ 内的全纯圆盘。稍后会再作一点简短讨论，因为这一问题有大量技术细节。这里只需指出，在我们考虑的情形中，它仅在 $\dim X\geq3$ 时成为问题。

\SourceParagraph{ch08-8.3.3-006}{611}
假设给定 $\operatorname{Fuk}(X)$ 的两个对象，为方便起见，将它们简记为 $L_0$、$L_1$。假设 $L_0$ 与 $L_1$ 横截相交。我们希望定义 $\operatorname{Hom}_{\operatorname{Fuk}(X)}(L_0,L_1)$。即使仅将它定义为一个分次向量空间，也需要几个概念，更不用说把它定义为一个复形了。

\begin{figure}[H]
\centering
\BookImage{01_Numbered_Figures/Chapter_08/Ch08_Fig17_PDFp612_Bookp602.pdf}
\caption{曲面 $X$ 内的一条可缩闭环 $L$ 总具有非平凡的 Maslov 类：在图示例子中，切线的角度从 $0$ 变化到 $2\pi$。}\label{fig:8.17}
\end{figure}

\SourceParagraph{ch08-8.3.3-007}{611}
首先，需要给出交点 $x\in L_0\cap L_1$ 的指标的正确概念。具体如下。考虑辛向量空间 $V:=T_xX$，以 $\operatorname{Gr}(V)$ 表示 Lagrange Grassmann 流形，即 $V$ 的所有 Lagrange 子空间组成的空间。因此，对于 $i=0,1$，$\Lambda_i:=T_xL_i$ 给出了 $\operatorname{Gr}(V)$ 中的两个点。由横截性可知，$\Lambda_0\cap\Lambda_1=0$。若 $\operatorname{Gr}(V)$ 中连接 $\Lambda_0$ 与 $\Lambda_1$ 的路径 $\Lambda_t$（$0\leq t\leq1$）满足 $\Lambda_0\cap\Lambda_t=0$，$0<t\leq1$，且在 $t=0$ 处的\emph{穿越形式}负定，则称这条路径为\emph{无穿越路径}。这一 $\Lambda_0$ 上的二次型如下定义：对充分小的 $t$，先选择同构 $\phi_t\colon\Lambda_0\to\Lambda_t$，使 $\phi_0$ 为恒等映射，再令二次型为
\[
v\longmapsto\left.\frac{d}{dt}\omega(v,\phi_t(v))\right|_{t=0}.
\]

\SourceParagraph{ch08-8.3.3-008}{612}
例如，总可以取 $V=\mathbb C^n$，配备标准辛形式，令 $\Lambda_0=\mathbb R^n$，并假设
\[
\Lambda_1=e^{\mathrm i\pi c_1}\mathbb R\times\cdots\times
e^{\mathrm i\pi c_n}\mathbb R,\qquad c_k\in(-1,0).
\]
那么
\[
\Lambda_t=e^{\mathrm i\pi t c_1}\mathbb R\times\cdots\times
e^{\mathrm i\pi t c_n}\mathbb R
\]
就是连接 $\Lambda_0$ 与 $\Lambda_1$ 的一条无穿越路径。此时，总可以选择函数 $\widetilde a\colon[0,1]\to\mathbb R$，使
\[
e^{2\pi\mathrm i\widetilde a(t)}
=\left(\frac{\Omega|_{\Lambda_t}}{|\Omega|_{\Lambda_t}|}\right)^2.
\]
然后，将交点 $x\in L_0\cap L_1$ 的\emph{绝对 Maslov 指标}定义为
\[
I(x):=(\widetilde\xi_{L_1}(x)-\widetilde a(1))
-(\widetilde\xi_{L_0}(x)-\widetilde a(0)).
\]
这个数实际上总属于 $\mathbb Z$。\BookRef{fig:8.18}{图18}给出了一个例子。

\begin{figure}[H]
\centering
\BookImage{01_Numbered_Figures/Chapter_08/Ch08_Fig18_PDFp613_Bookp603.pdf}
\caption{$\mathbb C$ 中的两个例子。箭头表示定向，分次为常数 $\widetilde\xi\in(-1,1)$。第一种情形的指标为 $1$，第二种情形的指标为 $0$。}\label{fig:8.18}
\end{figure}

\SourceParagraph{ch08-8.3.3-009}{612}
其次，由于收敛性问题，不能只在基域 $\mathbb C$ 上工作，而要使用 \emph{Novikov 环} $\Lambda_{\mathrm{nov}}$。它定义为形如
\[
\sum_{i\in\mathbb Z}a_iq^{\lambda_i}
\]
的形式幂级数所组成的环，其中系数 $a_i\in\mathbb Z$ 对所有充分小的负整数 $i$ 均为零，而实数列 $\lambda_i$ 满足
\[
\lim_{i\to\infty}\lambda_i=\infty.
\]
后一条件保证这个环中的乘法有意义。

\SourceParagraph{ch08-8.3.3-010}{613}
当 $L_0$ 与 $L_1$ 横截相交时，现在可以定义 Floer 复形 $\operatorname{Hom}(L_0,L_1)$。暂且限制于局部系统 $\mathcal L_0,\mathcal L_1$ 平凡的情形，稍后再考虑一般情形。令
\[
\operatorname{Hom}^i(L_0,L_1)=
\bigoplus_{\substack{x\in L_0\cap L_1\\I(x)=i}}[x]\cdot\Lambda_{\mathrm{nov}}.
\]

\SourceParagraph{ch08-8.3.3-011}{613}
利用伪全纯带状域的模空间来定义微分 $m_1$。取一族与辛形式 $\omega$ 相容的殆复结构 $J:=\{J_t\}_{0\leq t\leq1}$，即要求 $\omega(\cdot,J_t\cdot)$ 在 $X$ 上定义 Riemann 度量。对于 $p,q\in L_0\cap L_1$，以 $\mathcal M(p,q,L_0,L_1,J)$ 表示满足下列条件的映射 $u\colon\mathbb R\times[0,1]\to X$ 的模空间：
\[
\begin{gathered}
u(\mathbb R\times\{0\})\subseteq L_0,\qquad
u(\mathbb R\times\{1\})\subseteq L_1,\\
\frac{\partial u}{\partial\tau}+J_t(u)\frac{\partial u}{\partial t}=0,\\
\lim_{\tau\to-\infty}u(\tau,t)=p,\qquad
\lim_{\tau\to+\infty}u(\tau,t)=q.
\end{gathered}
\]
这里，$\tau$ 是 $\mathbb R$ 上的坐标，$t$ 是 $[0,1]$ 上的坐标。第二行应理解为 Cauchy--Riemann 方程。对于一般选取的 $J_t$，预期有
\[
\dim\mathcal M(p,q,L_0,L_1,J)=I(q)-I(p).
\]
这是通过计算某些算子的指标得到的，属于伪全纯曲线理论中的标准计算。这里不提供这些技术细节。

\SourceParagraph{ch08-8.3.3-012}{613}
注意，$\mathbb R$ 通过平移坐标 $\tau$ 作用于此模空间：对于 $r\in\mathbb R$，把解 $u(\cdot,\cdot)$ 替换为解 $u(\cdot+r,\cdot)$。这个 $\mathbb R$ 作用只有在 $p=q$ 且 $u$ 为常值映射时才不是自由的。于是得到商空间
\[
\mathcal M(p,q,L_0,L_1,J)/\mathbb R,
\]
其维数为 $I(q)-I(p)-1$；当 $p=q$ 时，约定该商为空。可以为这一空间赋予定向，这与我们对 $L_i$ 的自旋结构所作的选择有关。因此，当 $I(q)=I(p)+1$ 时，$\mathcal M(p,q,L_0,L_1,J)/\mathbb R$ 是一个有限集合，每个点 $u$ 都带有符号 $(-1)^{s(u)}$。

\SourceParagraph{ch08-8.3.3-013}{614}
与 Morse 范畴的情形一样，我们利用上述模空间定义 $m_1$。具体而言，若 $p,r\in L_0\cap L_1$，且 $I(r)=I(p)+1$，则 $m_1([p])$ 中 $[r]$ 的系数为
\[
\sum_{u\in\mathcal M(p,r,L_0,L_1,J)/\mathbb R}
(-1)^{s(u)}q^{\int u^*\omega}\in\Lambda_{\mathrm{nov}}.
\]

\SourceParagraph{ch08-8.3.3-014}{614}
若要证明 $m_1^2=0$，理想的做法是仿照 Morse 情形，在 $I(q)=I(p)+2$ 时考虑一维模空间 $\mathcal M(p,q,L_0,L_1,J)/\mathbb R$，再通过研究其边界得到 $m_1^2=0$。然而，与 Morse 情形不同，全纯带状域可以以几种不同方式退化：它可以退化为一个从 $p$ 连接到 $r$ 的带状域与一个从 $r$ 连接到 $q$ 的带状域的并，其中 $I(r)=I(p)+1$；也可能有一个全纯圆盘冒泡出来，其边界包含在 $L_0$ 或 $L_1$ 中，见\BookRef{fig:8.19}{图19}。（熟悉伪全纯曲线通常情形的读者，可能还会担心球面冒泡，但那是实余维为二的现象。）第一类退化提供关于 $m_1^2$ 的信息，而第二类退化是新的，可能导致 $m_1^2=0$ 的障碍。

\begin{figure}[H]
\centering
\BookImage{01_Numbered_Figures/Chapter_08/Ch08_Fig19_PDFp614_Bookp604.pdf}
\caption{全纯带状域的两种可能退化。第二种情形的绘图只是示意：圆盘的边界仍位于其中一个 Lagrange 子流形上。}\label{fig:8.19}
\end{figure}

\SourceParagraph{ch08-8.3.3-015}{615}
\BookCite{ref164}对这个问题进行了十分深入的研究。若边界位于 Lagrange 子流形 $L_0$、$L_1$ 上的全纯圆盘模空间满足某些性质——\BookSectionRef{sec:3.6.2.3}{3.6.2.3}节曾提示这些性质——则可以修改 $m_1$ 的定义，使 $m_1^2=0$。若 Lagrange 子流形 $L$ 不满足相应的性质，就称它\emph{有障碍}，不应将其视为 Fukaya 范畴的对象。这是对 Fukaya 范畴对象的一项附加限制。

\SourceParagraph{ch08-8.3.3-016}{615}
一般而言，可能有许多全纯圆盘的边界位于某个 Lagrange 子流形上，这自然引出了开 Gromov--Witten 不变量。目前尚无完全令人满意的开 Gromov--Witten 不变量数学理论，不过在许多有趣的情形下已有相应理论，见\BookCite{ref297}、\BookCite{ref429}。这里不展开这个主题，只简要讨论这种圆盘模空间的预期维数。

\SourceParagraph{ch08-8.3.3-017}{615}
给定全纯圆盘 $u\colon D\to X$，满足 $u(\partial D)\subseteq L$。对于任意辛流形 $X$，都可以定义 $u$ 的 Maslov 类 $\mu(u)$。具体如下：将丛 $u^*(TX)$ 作为辛向量丛平凡化为 $D\times\mathbb C^n$，其中 $\mathbb C^n$ 带有标准辛结构。这就定义了映射 $\gamma\colon\partial D\to\operatorname{Gr}(\mathbb C^n)$：它把 $x\in\partial D$ 映为在此平凡化下的 $T_{u(x)}L\subseteq\mathbb C^n$。与前面一样，$\mathbb C^n$ 上的标准全纯 $n$ 形式 $\Omega$ 定义映射 $\operatorname{Gr}(\mathbb C^n)\to S^1$，将由 $\eta_1,\ldots,\eta_n$ 张成的子空间映为
\[
\left(\frac{\Omega(\eta_1,\ldots,\eta_n)}
{|\Omega(\eta_1,\ldots,\eta_n)|}\right)^2.
\]
$H^1(S^1,\mathbb Z)$ 的正生成元拉回至 $H^1(\operatorname{Gr}(\mathbb C^n),\mathbb Z)$ 所得的类，称为 \emph{Maslov 类} $\mu$；它生成 $H^1(\operatorname{Gr}(\mathbb C^n),\mathbb Z)$。于是，$u$ 的 Maslov 类为
\[
\mu(u):=\int_{\partial D}\gamma^*\mu\in\mathbb Z.
\]

\SourceParagraph{ch08-8.3.3-018}{615}
若 $J$ 是一般的殆复结构，以 $\mathcal M(L,J)$ 表示边界位于 $L$ 内的 $J$ 全纯映射 $u\colon D\to X$ 的空间。在暂不考虑多重覆盖问题时，预期 $\mathcal M(L,J)$ 中包含 $u$ 的分支具有维数 $\tfrac12\dim_{\mathbb R}X+\mu(u)$。在我们的情形中，若 $L$ 具有分次，则实际上 $\mu(u)=0$。由于 $\operatorname{PSL}_2(\mathbb R)$ 通过重新参数化作用于 $\mathcal M(L,J)$，可得
\[
\dim\mathcal M(L,J)/\operatorname{PSL}_2(\mathbb R)
=\tfrac12\dim_{\mathbb R}X-3.
\]
因此，在低维情况下不必担心这种情形；只有当 Calabi--Yau 流形 $X$ 的复维数至少为 $3$ 时，它才成为问题。

\SourceParagraph{ch08-8.3.3-019}{615}
从现在起，假设我们处于不会出现边界位于某个 Lagrange 子流形上的非恒定全纯圆盘的情形。当然，这排除了我们最感兴趣的三维情形，但我们仍将在这个假设下继续。此时便能够定义
\[
m_d\colon\operatorname{Hom}(L_{d-1},L_d)\otimes\cdots\otimes
\operatorname{Hom}(L_0,L_1)\longrightarrow\operatorname{Hom}(L_0,L_d)[2-d].
\]

\SourceParagraph{ch08-8.3.3-020}{616}
对于 $p_{i-1,i}\in L_{i-1}\cap L_i$（$1\leq i\leq d$）、$p_{0,d}\in L_0\cap L_d$，以及某个一般的殆复结构 $J$，定义
\[
\mathcal M(p_{0,d};p_{0,1},\ldots,p_{d-1,d};L_0,\ldots,L_d)
\]
为如下映射 $u\colon D\to X$ 的模空间：$D$ 是一个圆盘，其边界上有 $d+1$ 个标记点 $x_{0,1},\ldots,x_{d-1,d},x_{0,d}\in\partial D$，按逆时针循环顺序排列，并满足
\[
\begin{aligned}
u(x_{i-1,i})&=p_{i-1,i},&&1\leq i\leq d,\\
u(x_{0,d})&=p_{0,d},\\
u([x_{i-1,i},x_{i,i+1}])&\subseteq L_i,&&1\leq i\leq d-1,\\
u([x_{d-1,d},x_{0,d}])&\subseteq L_d,\\
u([x_{0,d},x_{0,1}])&\subseteq L_0.
\end{aligned}
\]
这里，对于 $D$ 上相邻的两个标记点 $x,y$，$[x,y]$ 表示 $\partial D$ 上连接它们的弧段。这是全纯多边形的模空间，见\BookRef{fig:8.20}{图20}。

\begin{figure}[H]
\centering
\BookImage{01_Numbered_Figures/Chapter_08/Ch08_Fig20_PDFp616_Bookp606.pdf}
\BookNumberOnlyCaption\label{fig:8.20}
\end{figure}

\SourceParagraph{ch08-8.3.3-021}{616}
对于给定的带有 $d+1$ 个标记点的圆盘，此模空间的预期维数为 $I(p_{0,d})-\sum_{i=1}^d I(p_{i-1,i})$。由于这类带有 $d+1$ 个标记点的圆盘的模空间为维数 $d-2$ 的 Stasheff 结合多面体 $\mathcal S_d$，得到
\[
\begin{aligned}
&\dim\mathcal M(p_{0,d};p_{0,1},\ldots,p_{d-1,d};L_0,\ldots,L_d)
/\operatorname{PSL}_2(\mathbb R)\\
&\hspace{4em}=I(p_{0,d})+d-2-\sum_{i=1}^d I(p_{i-1,i}).
\end{aligned}
\]
当这个维数为零时，$m_d([p_{d-1,d}],\ldots,[p_{0,1}])$ 中 $[p_{0,d}]$ 的系数为
\[
\sum_{u\in\mathcal M(p_{0,d};p_{0,1},\ldots,p_{d-1,d};L_0,\ldots,L_d)
/\operatorname{PSL}_2(\mathbb R)}
(-1)^{s(u)}q^{\int u^*\omega}\in\Lambda_{\mathrm{nov}}.
\]
这里，$s$ 仍表示那个尚未展开解释的符号选择，它依赖于自旋结构的选择。

\SourceParagraph{ch08-8.3.3-022}{617}
假设全纯圆盘冒泡不会造成问题，则这些 $m_k$ 满足 $A_\infty$ 关系。证明仍然采用通常的方法，即分析一维圆盘模空间的紧化方式。

\SourceParagraph{ch08-8.3.3-023}{617}
现在说明，当 Lagrange 子流形 $L_i$ 带有非平凡的平坦 $U(1)$ 丛 $\mathcal L_i$ 时，如何修改上述公式。此时使用复数系数的 Novikov 环。首先，Floer 复形更自然的写法是
\[
\operatorname{Hom}^i((L_0,\mathcal L_0),(L_1,\mathcal L_1))
=\bigoplus_{\substack{p\in L_0\cap L_1\\I(p)=i}}
\operatorname{Hom}(\mathcal L_{0,p},\mathcal L_{1,p})\otimes\Lambda_{\mathrm{nov}}.
\]
因此，给定 $p\in L_0\cap L_1$ 和 $t\in\operatorname{Hom}(\mathcal L_{0,p},\mathcal L_{1,p})$，便得到一个元素
\[
t_p\in\operatorname{Hom}^*((L_0,\mathcal L_0),(L_1,\mathcal L_1)).
\]

\SourceParagraph{ch08-8.3.3-024}{617}
其次，按照如下方式修改各映射 $m_d$。给定
\[
u\in\mathcal M(p_{0,d};p_{0,1},\ldots,p_{d-1,d};L_0,\ldots,L_d),
\]
对于 $1\leq i\leq d-1$，利用 $\mathcal L_i$ 的联络，沿 $u([x_{i-1,i},x_{i,i+1}])$ 作平行移动，得到对应
\[
H_i^u\colon\mathcal L_{i,p_{i-1,i}}\xrightarrow{\sim}\mathcal L_{i,p_{i,i+1}}.
\]
沿 $u([x_{0,d},x_{0,1}])$ 在 $\mathcal L_0$ 中作平行移动，以及沿 $u([x_{d-1,d},x_{0,d}])$ 在 $\mathcal L_d$ 中作平行移动，则分别得到
\[
H_0^u\colon\mathcal L_{0,p_{0,d}}\xrightarrow{\sim}\mathcal L_{0,p_{0,1}},
\qquad
H_d^u\colon\mathcal L_{d,p_{d-1,d}}\xrightarrow{\sim}\mathcal L_{d,p_{0,d}}.
\]
于是，将 $m_d(t_{p_{d-1,d}},\ldots,t_{p_{0,1}})$ 在点 $p_{0,d}$ 处的系数定义为
\[
\begin{aligned}
&\sum_u(-1)^{s(u)}
\bigl(H_d^u\circ t_{p_{d-1,d}}\circ H_{d-1}^u\circ
t_{p_{d-2,d-1}}\circ\cdots\circ H_0^u\bigr)
\otimes q^{\int u^*\omega}\\
&\hspace{3em}\in\operatorname{Hom}(\mathcal L_{0,p_{0,d}},\mathcal L_{d,p_{0,d}})
\otimes\Lambda_{\mathrm{nov}},
\end{aligned}
\]
其中，求和遍历 $\mathcal M(p_{0,d};p_{0,1},\ldots,p_{d-1,d};L_0,\ldots,L_d)
/\operatorname{PSL}_2(\mathbb R)$ 中的所有 $u$。同一个公式也适用于 $m_1$，以及秩更高的平坦丛。

\SourceParagraph{ch08-8.3.3-025}{617}
到目前为止，这些定义仅适用于 $L_0,\ldots,L_k$ 两两横截的情形，而一般情况下当然未必如此。有两种处理方式。第一种是只定义一个 $A_\infty$ 预范畴；此后若要证明同调镜像对称，就证明它作为预范畴与镜像上的凝聚层导出范畴拟等价。第二种是证明这一构造在 Hamilton 形变下不变。\BookCite{ref421}按如下方式利用这一点：对于每对 Lagrange 子流形 $L_0,L_1$，在 $X$ 上选择一族 Hamilton 函数 $H_t\colon X\to\mathbb R$，$t\in[0,1]$，使相应的时间 $t$ 的 Hamilton 流 $\phi^t$ 满足 $\phi^1(L_0)$ 与 $L_1$ 横截相交。然后仍按照前面的方式定义 $\operatorname{Hom}^*(L_0,L_1)$，但将 $L_0$ 替换为 $\phi^1(L_0)$。这样便能避开横截性问题。特别地，对于单个对象 $L$，也得到了 $\operatorname{Hom}^*(L,L)$。随后可证明，所得 $A_\infty$ 范畴在拟同构意义下与这些选择无关；细节见\BookCite{ref421}。

\SourceParagraph{ch08-8.3.3-026}{618}
虽然前面的一般性讨论较为粗略，但在 $\dim_{\mathbb R}X=2$ 时，我们希望说得更具体一些。对我们而言，这就是 $X$ 为椭圆曲线的情形；下一节将详细处理它。目前，先讨论\BookCite{ref421}、\BookCite{ref002}分析过的符号问题。

\SourceParagraph{ch08-8.3.3-027}{618}
当 $X$ 是椭圆曲线时，任意一维子流形都是 Lagrange 子流形。将 $X$ 写为 $\mathbb C/\Gamma$，其中 $\Gamma$ 为格，便可将 Lagrange 子流形提升到万有覆盖 $\widetilde X=\mathbb C$。因此，给定圆周 $L\subseteq X$ 及其一个提升 $\widetilde L\subseteq\widetilde X$，$\widetilde L$ 或者是圆周（当 $L$ 在 $H_1(X,\mathbb Z)$ 中表示零元时），或者是一条与 $\mathbb R$ 微分同胚的直线。前一种情形没有分次，如\BookRef{fig:8.17}{图17}所示。因此，可以假设 $\widetilde L$ 是一条定向直线。此时，在 $L$ 上取由以圆周 $L$ 为边界的抽象圆盘 $D$ 的唯一自旋结构所诱导的自旋结构。这就是 Neveu--Schwarz 自旋结构，见\BookRef{ex:2.3}{例题2.3}。\BookRef{fig:8.18}{图18}随后给出了交点的指标。

\SourceParagraph{ch08-8.3.3-028}{618}
计算 $A_\infty$ 范畴结构时出现的全纯圆盘，现在可以取为普通的全纯圆盘，也就是 $\mathbb C$ 中的多边形，其边界包含于 $\widetilde L_0\cup\cdots\cup\widetilde L_d$，见\BookRef{fig:8.21}{图21}。图中标出了交点 $p_{0,1},\ldots,p_{3,4}$ 以及 $p_{0,4}$ 的指标。这个图给出一个元素
\[
u\in\mathcal M(p_{0,4};p_{0,1},\ldots,p_{3,4};L_0,\ldots,L_4).
\]
当然，提升 $\widetilde L_i$ 的不同选择会产生不同的圆盘。符号 $(-1)^{s(u)}$ 按如下方式确定：对于每个指标为 $1$ 的点 $p_{i,j}$，$i<j$，有一个符号贡献 $(-1)^{s(p_{i,j})}$；当且仅当圆盘边界的自然定向与 $L_j$ 的给定定向一致时，这个符号为正。例如，在\BookRef{fig:8.21}{图21}中，指标为 $1$ 的交点 $p_{0,1}$ 贡献 $+1$，$p_{2,3}$ 贡献 $-1$，总符号为 $(-1)(+1)=-1$。（在利用带状域，或等价地利用二边形，定义 $\operatorname{Hom}(L_0,L_1)$ 中的 $m_1$ 时，规则是：若 $L_1$ 的给定定向与带状域边界诱导的定向一致，则符号为 $+1$；否则为 $-1$。）

\begin{figure}[H]
\centering
\BookImage{01_Numbered_Figures/Chapter_08/Ch08_Fig21_PDFp618_Bookp608.pdf}
\caption{图中标出了圆盘的标准定向。}\label{fig:8.21}
\end{figure}

\SourceParagraph{ch08-8.3.3-029}{619}
这里还应注意，零维模空间中的圆盘在交点处总有凸角。如\BookRef{fig:8.22}{图22}所示，圆盘的凹角在局部会在一维族中变化：第二幅图中，圆盘边界沿自身折返，而发生折返的点显然可以移动。

\begin{figure}[H]
\centering
\BookImage{01_Numbered_Figures/Chapter_08/Ch08_Fig22_PDFp619_Bookp609.pdf}
\BookNumberOnlyCaption\label{fig:8.22}
\end{figure}

\begin{example*}
\SourceParagraph{ch08-8.3.3-030}{619}
可以在一个简单例子中实施这一过程。取格 $\Gamma=\langle1,\tau\rangle\subseteq\mathbb C$，以及 $L=\mathbb R/\mathbb Z\subseteq X=\mathbb C/\Gamma$。$L$ 的一个提升 $\widetilde L$ 就是 $\mathbb R$。我们希望计算 $\operatorname{Hom}^*(L,L)$。按照上面的讨论，应取 $L_0=L$，而 $L_1$ 为 $L$ 的一个 Hamilton 扰动。在圆柱 $\mathbb C/\mathbb Z$ 中，把 $L$ 的一个小邻域视为圆周 $L$ 的余切丛的零截面邻域，则 $L$ 的一个可能的 Hamilton 形变就是 $L$ 上某个恰当微分的图像，例如
\[
L_1=\{(x,\varepsilon\sin(2\pi x))\mid x\in\mathbb R\},
\]
其中 $\varepsilon>0$ 充分小，使扰动位于这个邻域内。考虑到周期性，有 $L_0\cap L_1=\{p_0,p_1\}$，其中 $p_1=(0,0)$、$p_0=(1/2,0)$；$p_1$ 的指标为 $1$，$p_0$ 的指标为 $0$。\BookRef{fig:8.23}{图23}给出了两个全纯带状域，并标明了它们的符号贡献。因此，$\operatorname{Hom}^*(L_0,L_1)$ 是
\[
[p_0]\Lambda_{\mathrm{nov}}\longrightarrow[p_1]\Lambda_{\mathrm{nov}},
\]
并且
\[
m_1([p_0])=\left(-q^{\int u_1^*\omega}+q^{\int u_2^*\omega}\right)[p_1],
\]
其中，$u_1\colon D\to X$ 和 $u_2\colon D\to X$ 分别对应左边和右边的圆盘。注意到 $\int u_1^*\omega=\int u_2^*\omega$——任何类似的 Hamilton 形变都具有这一性质——可见实际上 $m_1=0$，从而恢复了 $L$ 的以 $\Lambda_{\mathrm{nov}}$ 为系数的上同调 $H^*(L,\Lambda_{\mathrm{nov}})$，正如所预期的那样。
\end{example*}

\begin{figure}[H]
\centering
\BookImage{01_Numbered_Figures/Chapter_08/Ch08_Fig23_PDFp620_Bookp610.pdf}
\BookNumberOnlyCaption\label{fig:8.23}
\end{figure}

\SourceParagraph{ch08-8.3.3-031}{619}
下一节将看到一些更复杂的例子，其中多边形参与椭圆曲线的高阶乘法映射。在此之前，我们最后讨论一下与 $A_\infty$ 范畴有关的代数问题。

\subsection{恒等态射与三角 \texorpdfstring{$A_\infty$}{A∞} 范畴}\label{sec:8.3.4}

\SourceParagraph{ch08-8.3.4-001}{620}
在比较 $X$ 上凝聚层的导出范畴与镜像 $\check X$ 的 Fukaya 范畴时，至少还会遇到两个代数问题，它们似乎使人预期二者不可能同构。第一个问题是，$D^b(X)$ 具有恒等态射，而 $\operatorname{Fuk}(\check X)$ 事先并不具有。第二个问题是，$D^b(X)$ 是三角范畴，而 $\operatorname{Fuk}(\check X)$ 事先并不是。第一个问题结果并不构成障碍：事实上，与最初的预期相反，$\operatorname{Fuk}(\check X)$ 与一个具有恒等态射的 $A_\infty$ 范畴拟同构。第二个问题稍复杂一些。大致而言，与普通范畴相比，三角范畴需要附加数据，即一组满足\BookSectionRef{sec:4.4.3}{4.4.3}节公理的正合三角形。然而，对 $A_\infty$ 范畴而言，成为三角范畴仅仅是一项性质。按照我们目前的定义，尚不知道 $\operatorname{Fuk}(\check X)$ 是否具有这一性质。不过，有一个代数过程可以把它替换为扭曲范畴 $\operatorname{Tw}(\operatorname{Fuk}(\check X))$，后者是包含 $\operatorname{Fuk}(\check X)$ 的三角范畴。介绍这一过程还需要一些工作。

\begin{definition*}
\SourceParagraph{ch08-8.3.4-002}{620}
先讨论恒等态射。设 $\mathcal A$ 是一个（非幺）$A_\infty$ 范畴。恒等态射有几种有用的概念。首先，若非幺的普通范畴 $H^*(\mathcal A)$ 实际上具有恒等态射，即 $H^*(\mathcal A)$ 实际上是一个通常意义的范畴，则称 $\mathcal A$ \emph{上同调幺}，或 \emph{$c$ 幺}。

\SourceParagraph{ch08-8.3.4-003}{620}
另一方面，若对每个对象 $X$ 都存在一个必然唯一的元素 $e_X\in\operatorname{Hom}_{\mathcal A}^0(X,X)$，满足
\[
\begin{gathered}
m_1(e_X)=0,\\
(-1)^{\deg a}m_2(e_{X_1},a)=a=m_2(a,e_{X_0}),
\qquad a\in\operatorname{Hom}_{\mathcal A}(X_0,X_1),\\
\text{任何含有 }e_X\text{ 的高阶乘法映射均为零，}
\end{gathered}
\]
则称 $\mathcal A$ \emph{严格幺}。
\end{definition*}

\SourceParagraph{ch08-8.3.4-004}{620}
首先，$\operatorname{Fuk}(X)$ 实际上是 $c$ 幺的。其原因是：在我们考虑的 Calabi--Yau 情形及本节排除圆盘冒泡的假设下，对于单个分次 Lagrange 子流形 $L$，Floer 上同调 $H^*(\operatorname{Hom}_{\operatorname{Fuk}(X)}(L,L))$ 同构于通常的上同调群 $H^*(L,\Lambda_{\mathrm{nov}})$，而 $1\in H^0(L,\Lambda_{\mathrm{nov}})$ 当然就是恒等元。Piunikhin、Salamon 与 Schwarz 在\BookCite{ref392}中证明了这一事实，虽然是在稍特殊的背景下。随后，Seidel 在\BookCite{ref421}推论 2.14 中证明，任何 $c$ 幺 $A_\infty$ 范畴都与一个严格幺 $A_\infty$ 范畴拟同构。这里不给出细节，但这说明我们不必担心这一问题。因此，从现在起，假设所有 $A_\infty$ 范畴均为严格幺的。

\SourceParagraph{ch08-8.3.4-005}{621}
接下来讨论三角结构。这需要引入若干新的概念。

\BookStatementNumber{14}
\begin{definition}\label{def:8.14}
\SourceParagraph{ch08-8.3.4-006}{621}
设 $\mathcal A$ 是一个 $A_\infty$ 范畴。$\mathcal A$ 上的一个（右）\emph{$A_\infty$ 模} $\mathcal M$，由每个 $X\in\operatorname{Ob}(\mathcal A)$ 对应的分次向量空间 $\mathcal M(X)$，以及对每个 $d\geq1$ 给定的下列乘法映射组成：
\setcounter{equation}{8}
\begin{equation}\label{eq:8.9}
\begin{aligned}
m_d^{\mathcal M}\colon{}&\mathcal M(X_{d-1})\otimes
\operatorname{Hom}_{\mathcal A}(X_{d-2},X_{d-1})\otimes\cdots
\otimes\operatorname{Hom}_{\mathcal A}(X_0,X_1)\\
&\longrightarrow\mathcal M(X_0)[2-d].
\end{aligned}
\end{equation}
这些映射应满足关系
\[
\begin{aligned}
&\sum_n(-1)^{\deg a_1+\cdots+\deg a_n-n}
m_{n+1}^{\mathcal M}\bigl(\\
&\qquad m_{d-n}^{\mathcal M}(b,a_{d-1},\ldots,a_{n+1}),a_n,\ldots,a_1\bigr)\\
&\quad+\sum_{m,n}(-1)^{\deg a_1+\cdots+\deg a_n-n}\,
m_{d-m+1}^{\mathcal M}\bigl(b,a_{d-1},\ldots,a_{n+m+1},\\
&\qquad\quad m_m^{\mathcal A}(a_{n+m},\ldots,a_{n+1}),a_n,\ldots,a_1\bigr)=0.
\end{aligned}
\]

\SourceParagraph{ch08-8.3.4-007}{621}
第二个和遍历 $n+m<d$ 的指标。例如，当 $d=1$ 时，这个关系成为 $(m_1^{\mathcal M})^2=0$，即 $m_1^{\mathcal M}$ 使 $\mathcal M(X)$ 成为链复形。当 $d=2$ 时，它为
\[
m_1^{\mathcal M}(m_2^{\mathcal M}(b,a_1))
+(-1)^{\deg a_1-1}m_2^{\mathcal M}(m_1^{\mathcal M}(b),a_1)
+m_2^{\mathcal M}(b,m_1^{\mathcal A}(a_1))=0,
\]
而当 $d=3$ 时，它为
\[
\begin{aligned}
&m_1^{\mathcal M}(m_3^{\mathcal M}(b,a_2,a_1))
+(-1)^{\deg a_1-1}m_2^{\mathcal M}(m_2^{\mathcal M}(b,a_2),a_1)\\
&\quad+(-1)^{\deg a_1+\deg a_2-2}
m_3^{\mathcal M}(m_1^{\mathcal M}(b),a_2,a_1)\\
&\quad+m_3^{\mathcal M}(b,a_2,m_1^{\mathcal A}(a_1))\\
&\quad+(-1)^{\deg a_1-1}m_3^{\mathcal M}(b,m_1^{\mathcal A}(a_2),a_1)
+m_2^{\mathcal M}(b,m_2^{\mathcal A}(a_2,a_1))=0.
\end{aligned}
\]
与原来的 $A_\infty$ 关系一样，后面这两个展开式中的第一个说明 $m_1^{\mathcal M}$ 对模乘法 $m_2^{\mathcal M}$ 满足 Leibniz 法则，而 $m_3^{\mathcal M}$ 则衡量这一模乘法偏离结合律的程度。
\end{definition}

\SourceParagraph{ch08-8.3.4-008}{621}
$\mathcal A$ 上的 $A_\infty$ 模本身也构成一个 $A_\infty$ 范畴 $\operatorname{Mod}(\mathcal A)$。需要定义 $A_\infty$ 模之间的态射概念，我们称之为\emph{预同态}。若 $\mathcal M_1,\mathcal M_2$ 是两个 $A_\infty$ 模，则次数为 $\deg t$ 的预同态 $t\colon\mathcal M_1\to\mathcal M_2$，由对每个 $d\geq1$ 给定的一列映射组成：
\setcounter{equation}{9}
\begin{equation}\label{eq:8.10}
\begin{aligned}
t_d\colon{}&\mathcal M_1(X_{d-1})\otimes
\operatorname{Hom}_{\mathcal A}(X_{d-2},X_{d-1})\otimes\cdots
\otimes\operatorname{Hom}_{\mathcal A}(X_0,X_1)\\
&\longrightarrow\mathcal M_2(X_0)[\deg t-d+1].
\end{aligned}
\end{equation}

\SourceParagraph{ch08-8.3.4-009}{622}
随后，在预同态的分次向量空间 $\operatorname{Hom}_{\operatorname{Mod}(\mathcal A)}(\mathcal M_1,\mathcal M_2)$ 上，按下式定义 $m_1^{\operatorname{Mod}(\mathcal A)}$：
\[
\begin{aligned}
&\bigl(m_1^{\operatorname{Mod}(\mathcal A)}(t)_d\bigr)(b,a_{d-1},\ldots,a_1)\\
&\quad=\sum_n(-1)^*m_{n+1}^{\mathcal M_2}
\bigl(t_{d-n}(b,a_{d-1},\ldots,a_{n+1}),a_n,\ldots,a_1\bigr)\\
&\qquad+\sum_n(-1)^*t_{n+1}
\bigl(m_{d-n}^{\mathcal M_1}(b,a_{d-1},\ldots,a_{n+1}),a_n,\ldots,a_1\bigr)\\
&\qquad+\sum_{m,n}(-1)^*t_{d-m+1}
\bigl(b,a_{d-1},\ldots,a_{n+m+1},
m_m^{\mathcal A}(a_{n+m},\ldots,a_{n+1}),a_n,\ldots,a_1\bigr).
\end{aligned}
\]
这里，$*$ 表示表达式
\[
\deg a_{n+1}+\cdots+\deg a_{d-1}+\deg b-d+n+1.
\]
此外，还需要定义复合：给定 $t^1\colon\mathcal M_0\to\mathcal M_1$ 和 $t^2\colon\mathcal M_1\to\mathcal M_2$，令
\[
\begin{aligned}
&\bigl(m_2^{\operatorname{Mod}(\mathcal A)}(t^2,t^1)\bigr)_d
(b,a_{d-1},\ldots,a_1)\\
&\quad=\sum_n(-1)^*t^2_{n+1}
\bigl(t^1_{d-n}(b,a_{d-1},\ldots,a_{n+1}),a_n,\ldots,a_1\bigr).
\end{aligned}
\]

\SourceParagraph{ch08-8.3.4-010}{622}
读者可以在低次情形将这些式子展开，以再次体会它们的含义。幸运的是，此时所有高阶乘法映射 $m_d^{\operatorname{Mod}(\mathcal A)}$（$d>2$）均为零，所以得到的是一个微分分次范畴。若 $m_1^{\operatorname{Mod}(\mathcal A)}(t)=0$，则称 $t$ 为\emph{同态}。这个范畴是严格幺的，其中 $\operatorname{id}\in\operatorname{Hom}_{\operatorname{Mod}(\mathcal A)}(\mathcal M,\mathcal M)$ 由
\[
\operatorname{id}_1(b)=(-1)^{\deg b}b,\qquad
\operatorname{id}_d=0\quad(d>1)
\]
给出。所谓\emph{同构}，是指在 $H^*(\operatorname{Mod}(\mathcal A))$ 中为同构的同态。

\SourceParagraph{ch08-8.3.4-011}{622}
现在可以定义称为 \emph{Yoneda 嵌入}的 $A_\infty$ 函子 $\operatorname{Yon}\colon\mathcal A\to\operatorname{Mod}(\mathcal A)$。它的定义很简单：对于 $Y\in\operatorname{Ob}(\mathcal A)$，令 $\operatorname{Yon}(Y)$ 为模
\[
X\longmapsto\operatorname{Hom}_{\mathcal A}(X,Y).
\]
当然，后者已经是由 $m_1^{\mathcal A}$ 给出的链复形。为了定义模结构，需要给出\BookEqRef{eq:8.9}{8.9}中的映射。但在此情况下，$\mathcal M(X_{d-1})=\operatorname{Hom}_{\mathcal A}(X_{d-1},Y)$，且 $\mathcal M(X_0)=\operatorname{Hom}_{\mathcal A}(X_0,Y)$，所以这些映射直接由 $m_d^{\mathcal A}$ 给出。这样便定义了函子 $\operatorname{Yon}$ 的对象映射，将 $\mathcal A$ 的对象映为 $\operatorname{Mod}(\mathcal A)$ 的对象。不过，还必须定义高阶映射
\[
\begin{aligned}
\operatorname{Yon}_d\colon{}&\operatorname{Hom}_{\mathcal A}(X_{d-1},X_d)
\otimes\cdots\otimes\operatorname{Hom}_{\mathcal A}(X_0,X_1)\\
&\longrightarrow
\operatorname{Hom}_{\operatorname{Mod}(\mathcal A)}
(\operatorname{Yon}(X_0),\operatorname{Yon}(X_d))[1-d].
\end{aligned}
\]

\SourceParagraph{ch08-8.3.4-012}{622}
因此，给定 $a_1,\ldots,a_d$，需要给出一个预同态，即一组映射
\[
\begin{aligned}
t_n\colon&\operatorname{Yon}(X_0)(Y_{n-1})\otimes
\operatorname{Hom}_{\mathcal A}(Y_{n-2},Y_{n-1})\otimes\cdots
\otimes\operatorname{Hom}_{\mathcal A}(Y_0,Y_1)\\
&\longrightarrow\operatorname{Yon}(X_d)(Y_0)
\left[2-d-n+\sum_i\deg a_i\right],
\end{aligned}
\]
也就是
\[
\begin{aligned}
(\operatorname{Yon}_d)_n\colon&\operatorname{Hom}_{\mathcal A}(Y_{n-1},X_0)
\otimes\operatorname{Hom}_{\mathcal A}(Y_{n-2},Y_{n-1})\otimes\cdots
\otimes\operatorname{Hom}_{\mathcal A}(Y_0,Y_1)\\
&\longrightarrow\operatorname{Hom}_{\mathcal A}(Y_0,X_d)
\left[2-d-n+\sum_i\deg a_i\right].
\end{aligned}
\]
有一个显然的选择：
\[
\begin{aligned}
&\bigl(\operatorname{Yon}_d(a_d,\ldots,a_1)\bigr)_n(b,b_{n-1},\ldots,b_1)\\
&\hspace{3em}=m_{d+n}^{\mathcal A}(a_d,\ldots,a_1,b,b_{n-1},\ldots,b_1).
\end{aligned}
\]
这定义了一个 $A_\infty$ 函子 $\operatorname{Yon}$，读者可以自行验证。

\SourceParagraph{ch08-8.3.4-013}{623}
类似于普通范畴论中的情形，若存在同构 $t\colon\operatorname{Yon}(Y)\to\mathcal M$，就称对象 $Y\in\operatorname{Ob}(\mathcal A)$ \emph{拟表示} $A_\infty$ 模 $\mathcal M$。这种处理的优点在于，$\operatorname{Mod}(\mathcal A)$ 中的代数操作容易定义；当这些操作的结果可拟表示时，它们便转化为 $\mathcal A$ 中的概念。有两个例子对我们至关重要。

\SourceParagraph{ch08-8.3.4-014}{623}
首先是\emph{平移函子}：给定 $A_\infty$ 模 $\mathcal M$，定义 $S\mathcal M$ 为
\[
(S\mathcal M)(Y)=\mathcal M(Y)[1].
\]
若对所有 $Y\in\operatorname{Ob}(\mathcal A)$，$S(\operatorname{Yon}(Y))$ 均可由某个 $SY\in\operatorname{Ob}(\mathcal A)$ 拟表示，便得到平移函子 $S\colon\mathcal A\to\mathcal A$。例如，在 Fukaya 范畴中，只需将分次加 $1$，就能实现平移函子。

\SourceParagraph{ch08-8.3.4-015}{623}
对我们更重要的是\emph{映射锥}构造。设 $Y_0,Y_1\in\operatorname{Ob}(\mathcal A)$，$c\in\operatorname{Hom}_{\mathcal A}^0(Y_0,Y_1)$，满足 $m_1^{\mathcal A}(c)=0$。可以将 $c$ 的抽象映射锥定义为如下 $A_\infty$ 模 $\operatorname{Cone}(c)$：
\[
\operatorname{Cone}(c)(X):=
\operatorname{Hom}_{\mathcal A}(X,Y_0)[1]\oplus
\operatorname{Hom}_{\mathcal A}(X,Y_1),
\]
\[
\begin{aligned}
&m_d^{\operatorname{Cone}(c)}((b_0,b_1),a_{d-1},\ldots,a_1)\\
&\quad=\bigl(m_d^{\mathcal A}(b_0,a_{d-1},\ldots,a_1),\,
m_d^{\mathcal A}(b_1,a_{d-1},\ldots,a_1)\\
&\hspace{6em}
+m_{d+1}^{\mathcal A}(c,b_0,a_{d-1},\ldots,a_1)\bigr).
\end{aligned}
\]
若 $\mathcal A$ 中存在一个拟表示 $\operatorname{Cone}(c)$ 的对象，也将它记为 $\operatorname{Cone}(c)$。

\SourceParagraph{ch08-8.3.4-016}{623}
模 $\operatorname{Cone}(c)$ 带有典范的预同态
\[
\begin{aligned}
\iota&\in\operatorname{Hom}_{\operatorname{Mod}(\mathcal A)}^0
(\operatorname{Yon}(Y_1),\operatorname{Cone}(c)),\\
\pi&\in\operatorname{Hom}_{\operatorname{Mod}(\mathcal A)}^1
(\operatorname{Cone}(c),\operatorname{Yon}(Y_0)),
\end{aligned}
\]
由
\[
\iota_1(b_1)=(0,(-1)^{\deg b_1}b_1),\qquad
\pi_1(b_0,b_1)=(-1)^{\deg b_0-1}b_0
\]
给出，而这些预同态的所有高阶部分均为零。由此在 $H^*(\operatorname{Mod}(\mathcal A))$ 中得到一个“三角形”：
\setcounter{equation}{10}
\begin{figure}[H]
\begin{equation}\label{eq:8.11}
\BookImage{02_Unnumbered_Diagrams/Chapter_08/Ch08_Diagram_PDFp623_01_Bookp613.pdf}
\end{equation}
\end{figure}

\SourceParagraph{ch08-8.3.4-017}{624}
这里，$[\cdot]$ 表示同态在 $H^*(\operatorname{Mod}(\mathcal A))$ 中的上同调类。由于 $m_1^{\operatorname{Mod}(\mathcal A)}$ 作用于 $\operatorname{Yon}_1(c)$、$\iota$ 或 $\pi$ 时均为零，这些上同调类有意义，读者可以自行验证。请记住，$[\pi]$ 的次数为一。

\SourceParagraph{ch08-8.3.4-018}{624}
现在可以定义 $\mathcal A$ 中\emph{正合三角形}的概念。它是 $H^*(\mathcal A)$ 中如下形式的图表：
\begin{figure}[H]
\centering
\BookImage{02_Unnumbered_Diagrams/Chapter_08/Ch08_Diagram_PDFp624_01_Bookp614.pdf}
\end{figure}
它在 Yoneda 嵌入下变为与\BookEqRef{eq:8.11}{8.11}同构的图表，其中 $c_1=c$。

\BookStatementNumber{15}
\begin{definition}\label{def:8.15}
\SourceParagraph{ch08-8.3.4-019}{624}
若 $A_\infty$ 范畴 $\mathcal A$ 满足以下条件，则称它为\emph{三角} $A_\infty$ 范畴：
\begin{enumerate}
\item $H^0(\mathcal A)$ 中的每个态射 $[c_1]$ 都可以扩充为一个正合三角形。
\item 存在平移函子 $S\colon\mathcal A\to\mathcal A$。
\item 对每个 $Y\in\operatorname{Ob}(\mathcal A)$，均存在对象 $\widetilde Y\in\operatorname{Ob}(\mathcal A)$，使 $S\widetilde Y\cong Y$ 在 $H^0(\mathcal A)$ 中成立。
\end{enumerate}
\end{definition}

\SourceParagraph{ch08-8.3.4-020}{624}
三角 $A_\infty$ 范畴的一个例子是 $D^b_\infty(X)$：对于它的对象 $M$，平移函子就是\BookSectionRef{sec:4.4.3}{4.4.3}节定义的通常的复形平移 $M\mapsto M[1]$。此外，这个范畴中的映射锥也按照\BookSectionRef{sec:4.4.3}{4.4.3}节的方式定义。

\SourceParagraph{ch08-8.3.4-021}{624}
另一方面，尚不知道 Fukaya 范畴是否是三角的。至少预计需要把浸入而非嵌入的 Lagrange 子流形也加入 Fukaya 范畴，而这会给 Floer 同调的定义带来技术困难。所以，仍需解决这一问题。解决办法是\emph{扭曲构造}：给定 $A_\infty$ 范畴 $\mathcal A$，必要时先加入对象的形式平移，再定义新的 $A_\infty$ 范畴 $\operatorname{Tw}(\mathcal A)$，即\emph{扭曲复形范畴}；它是三角的。

\BookStatementNumber{16}
\begin{definition}\label{def:8.16}
\SourceParagraph{ch08-8.3.4-022}{624}
设 $\mathcal A$ 是一个 $A_\infty$ 范畴。一个\emph{扭曲复形}，由 $\mathcal A$ 的一列对象 $X_1,\ldots,X_n$（在需要时包括它们的形式平移），以及一个由一次态射组成的严格下三角矩阵 $\Delta=(\delta_{i,j})_{1\leq i,j\leq n}$ 构成，其中 $\delta_{i,j}\in\operatorname{Hom}^1(X_j,X_i)$，并满足
\[
\sum_{k\geq1}m_k(\Delta,\ldots,\Delta)=0.
\]
这里按通常的矩阵乘法规则，理解 $m_k$ 作用于态射矩阵的含义。因而，利用严格下三角条件，上式等价于对每对 $p>q$ 都有
\[
\sum_{\substack{k\geq1\\q=i_0<i_1<\cdots<i_k=p}}
m_k(\delta_{i_k,i_{k-1}},\ldots,\delta_{i_1,i_0})=0.
\]

\SourceParagraph{ch08-8.3.4-023}{625}
注意，这个看似无穷的和因此只有有限个非零项。记 $X=(X_1,\ldots,X_n)$，$\Delta=\Delta^X$，于是将扭曲对象记为 $(X,\Delta^X)$。
\end{definition}

\SourceParagraph{ch08-8.3.4-024}{625}
$\mathcal A$ 的扭曲范畴是以 $\mathcal A$ 的扭曲复形为对象的 $A_\infty$ 范畴，其态射空间为
\[
\operatorname{Hom}_{\operatorname{Tw}(\mathcal A)}((X,\Delta^X),(Y,\Delta^Y))
=\bigoplus_{i,j}\operatorname{Hom}_{\mathcal A}(X_i,Y_j).
\]
若 $(X^0,\Delta^{X^0}),\ldots,(X^d,\Delta^{X^d})$ 是一列扭曲复形，并且
\[
f^i:=(f_{j,k}^i)\in\operatorname{Hom}_{\operatorname{Tw}(\mathcal A)}
((X^{i-1},\Delta^{X^{i-1}}),(X^i,\Delta^{X^i})),
\]
则映射 $m_d^{\operatorname{Tw}(\mathcal A)}$ 为
\[
\begin{aligned}
&m_d^{\operatorname{Tw}(\mathcal A)}(f^d,\ldots,f^1)\\
&\quad=\sum_{j_0,\ldots,j_d}
m_{d+j_0+\cdots+j_d}^{\mathcal A}\bigl(
\underbrace{\Delta^{X^d},\ldots,\Delta^{X^d}}_{j_d},f^d,\\
&\hspace{6em}
\underbrace{\Delta^{X^{d-1}},\ldots,\Delta^{X^{d-1}}}_{j_{d-1}},
f^{d-1},\ldots,f^1,
\underbrace{\Delta^{X^0},\ldots,\Delta^{X^0}}_{j_0}\bigr).
\end{aligned}
\]
需要作一点说明：首先，求和遍历所有 $j_0,\ldots,j_d\geq0$；其次，$\Delta^{X^i}$ 在表达式中出现 $j_i$ 次；第三，与前面一样，应将该表达式理解为矩阵乘法。我们将验证这些 $m_d$ 满足 $A_\infty$ 关系的工作留给读者！

\SourceParagraph{ch08-8.3.4-025}{625}
最简单的例子是直和，此时 $\Delta^X=0$，所以对象 $((X_1,\ldots,X_n),0)$ 表示 $X_1\oplus\cdots\oplus X_n$。一个稍不平凡的扭曲复形例子是映射锥构造。假设 $A_\infty$ 范畴 $\mathcal A$ 具有平移函子 $S$，Fukaya 范畴就是如此。给定满足 $m_1^{\mathcal A}(c)=0$ 的 $c\in\operatorname{Hom}_{\mathcal A}^0(Y_0,Y_1)$，得到扭曲复形
\[
C=\left((SY_0,Y_1),\begin{pmatrix}0&0\\-S(c)&0\end{pmatrix}\right),
\]
其中，利用平移函子 $S$，把 $c$ 视为定义了 $\operatorname{Hom}_{\mathcal A}^1(SY_0,Y_1)$ 中的一个元素 $S(c)$。注意，存在典范态射
\[
\operatorname{Hom}^0(Y_1,C)=
\operatorname{Hom}^0(Y_1,SY_0)\oplus\operatorname{Hom}^0(Y_1,Y_1)
\]
中的 $(0,\operatorname{id}_{Y_1})$，以及
\[
\begin{aligned}
\operatorname{Hom}^1(C,Y_0)
&=\operatorname{Hom}^1(SY_0,Y_0)\oplus\operatorname{Hom}^1(Y_1,Y_0)\\
&=\operatorname{Hom}^0(Y_0,Y_0)\oplus\operatorname{Hom}^1(Y_1,Y_0)
\end{aligned}
\]
中的 $(\operatorname{id}_{Y_0},0)$。读者同样可以验证，这些映射在 $\operatorname{Tw}(\mathcal A)$ 中给出一个正合三角形，即 $C$ 在 $\operatorname{Tw}(\mathcal A)$ 中表示 $c$ 的映射锥。

\SourceParagraph{ch08-8.3.4-026}{625}
事实上，稍作推广就可看出，$\operatorname{Tw}(\mathcal A)$ 本身是三角的。给定闭态射
\[
c\in\operatorname{Hom}_{\operatorname{Tw}(\mathcal A)}^0
((X^0,\Delta^{X^0}),(X^1,\Delta^{X^1})),
\]
将 $\operatorname{Cone}(c)$ 构造为
\[
C=\left((SX^0,X^1),
\begin{pmatrix}\Delta^{SX^0}&0\\-S(c)&\Delta^{X^1}\end{pmatrix}\right),
\]
其中左上块为将 $X^0$ 连同其扭曲数据一并平移所得的矩阵。

\SourceParagraph{ch08-8.3.4-027}{626}
可以将 $\operatorname{Tw}(\mathcal A)$ 视为 $\mathcal A$ 的导出范畴，并把它写为 $D^b(\mathcal A)$。

\SourceParagraph{ch08-8.3.4-028}{626}
遗憾的是，这还不完全是我们所需要的范畴。$X$ 上凝聚层的导出范畴 $D^b(X)$ 还有一个 $\operatorname{Fuk}(\check X)$ 或 $D^b(\operatorname{Fuk}(\check X))$ 未必具有的性质：若 $F$ 是 $D^b(X)$ 的对象，$p\colon F\to F$ 是满足 $p^2=p$ 的态射，则可以把 $p$ 的像看作 $D^b(X)$ 中的另一个对象。这样的态射 $p$ 称为\emph{幂等态射}，而 $D^b(X)$ 是\emph{分裂闭}的；实质上，这意味着 $D^b(X)$ 包含所有幂等态射的像。

\SourceParagraph{ch08-8.3.4-029}{626}
一般而言，对于给定的线性范畴，可以构造一个包含所有幂等态射的像的更大范畴，即 \emph{Karoubi 完备化}。这可以推广到 $A_\infty$ 范畴，从 $A_\infty$ 范畴 $\mathcal A$ 得到一个分裂闭范畴。我们略去构造细节；感兴趣的读者可参阅\BookCite{ref421}第 4 节。该构造与前面看到的构造具有相似的特点。特别地，以 $D^\pi(\mathcal A)$ 表示将这一过程施加于 $D^b(\mathcal A)$ 所得的结果。可以证明它是一个分裂闭的三角 $A_\infty$ 范畴。完成这一步后，我们所希望的同调镜像对称的严格表述就成为：

\BookStatementNumber{17}
\begin{conjecture}\label{conj:8.17}
\SourceParagraph{ch08-8.3.4-030}{626}
设 $X$ 和 $\check X$ 为一对互为镜像的 Calabi--Yau 流形。则 $D^b_\infty(X)$ 与 $D^\pi(\operatorname{Fuk}(\check X))$ 作为 $A_\infty$ 范畴拟等价。
\end{conjecture}

\section{椭圆曲线}\label{sec:8.4}
\setcounter{equation}{11}
\setcounter{figure}{23}

\SourceParagraph{ch08-8.4-intro-001}{626}
现在证明椭圆曲线 HMS 的一个受限版本。为简单起见，不尝试比较完整的导出范畴与 Fukaya 范畴，不过，这里的工作已包含其中大部分。椭圆曲线情形已在\BookCite{ref397,ref396}中详细研究；这里采用稍有不同的方法，受到 Abouzaid\BookCite{ref002}启发，并在\BookCite{ref001}中发展，利用 Čech 方法定义导出范畴的 $A_\infty$ 结构。

\SourceParagraph{ch08-8.4-intro-002}{626}
将构造椭圆曲线的退化族 $\mathcal X\to D$。$D$ 可以看作坐标为 $q$ 的单位圆盘；在概形范畴中，则取 $D=\operatorname{Spec}\mathbb C[[q]]$。此族为\BookSectionRef{sec:6.2.4}{6.2.4}中所考虑族的紧化。使用这样的族的优点，是参数 $q$ 对应于镜像 Fukaya 范畴中的 Kähler 参数 $q$，也就是 Novikov 环的变量。特别地，应该将两侧都看作由 $q$ 参数化的范畴族。此外，采用显式退化，使这里的方法与 Gross、Siebert\BookCite{ref205,ref206}提出的思想很好地契合。虽然不在这里展开这些思想，它们在这一情形仍很有帮助。

\SourceParagraph{ch08-8.4-intro-003}{626}
实际上，将在构造退化 $\mathcal X\to D$ 的同时，构造线丛 $\mathcal L$，其在每条纤维上的限制都丰沛，即正。随后，计算涉及 $\mathcal L$ 的正、负张量幂的全部高阶乘法映射。

\subsection{退化：Tate 曲线}\label{sec:8.4.1}

\SourceParagraph{ch08-8.4.1-001}{627}
所用退化构造方法，是 Mumford 关于 Abel 簇退化的构造\BookCite{ref373}的一个特例。这一椭圆曲线退化通常称为 Tate 曲线。从下述数据出发：选择线丛的正整数次数 $d$，考虑连续分段线性函数 $\varphi\colon\mathbb R\to\mathbb R$，在每个区间 $[i,i+1]$ 上斜率为 $i$，$i\in\mathbb Z$。具体地，可以令它在 $[0,1]$ 上取值零；此时，
\[
\varphi(x)=ix-\frac{i(i+1)}2,\qquad x\in[i,i+1],
\]
且满足周期性条件
\begin{equation}\label{eq:8.12}
\varphi(x+d)=\varphi(x)+dx+\frac{d(d-1)}2.
\end{equation}
令 $\Delta=\{(x,y)\in\mathbb R^2\mid y\geq\varphi(x)\}$。
\begin{figure}[H]
\centering
\BookImage{01_Numbered_Figures/Chapter_08/Ch08_Fig24_PDFp627_Bookp617.pdf}
\caption{$\varphi$ 的图像。阴影区域为 $\Delta$。}\label{fig:8.24}
\end{figure}

\SourceParagraph{ch08-8.4.1-002}{627}
正如凸多胞体决定射影环面簇，无界多面体也可以决定环面簇。这可通过 $\Delta$ 的法扇描述。对于每个面 $\sigma\subseteq\Delta$，定义 $\Delta$ 在 $\sigma$ 处的法锥
\[
\check\sigma=\left\{n\in(\mathbb R^2)^\vee\ \middle|\
\begin{gathered}
n|_\sigma\text{ 为常数},\\
\langle n,m\rangle\geq\langle n,m'\rangle
\quad\forall m\in\Delta,\ m'\in\sigma
\end{gathered}\right\}.
\]
法扇为各面对应的锥的集合 $\{\check\sigma\mid\sigma\subseteq\Delta\text{ 为面}\}$，位于 $(\mathbb R^2)^\vee$ 中。容易看出，若 $\sigma$ 是端点为 $(i,\varphi(i))$、$(i+1,\varphi(i+1))$ 的区间，则 $\check\sigma$ 为 $(-i,1)$ 生成的射线。因此，$\Delta$ 的法扇 $\Sigma$ 如\BookRef{fig:8.25}{图25}所示。它包含无穷多个锥，超出了通常对扇的有限性要求，不过，这并不妨碍构造：它仍定义环面簇 $X_\Sigma$，由这些锥给出的无穷多个仿射环面坐标片覆盖。

\SourceParagraph{ch08-8.4.1-003}{628}
此外，$\mathbb R^2$ 中的 $(0,1)$ 在 $\Sigma$ 每条射线的本原生成元上取值一。将对应于 $(0,1)$ 的单项式记为 $q$，则它为 $X_\Sigma$ 上的正则函数，定义映射 $g\colon X_\Sigma\to\mathbb A^1$。
\begin{figure}[H]
\centering
\BookImage{01_Numbered_Figures/Chapter_08/Ch08_Fig25_PDFp628_Bookp618.pdf}
\caption{$\Delta$ 的法扇。}\label{fig:8.25}
\end{figure}

\SourceParagraph{ch08-8.4.1-004}{628}
由于 $q$ 在 $X_\Sigma$ 的每个环面除子上，即每条射线对应的除子上，以一阶消失，纤维 $g^{-1}(0)$ 为无穷多个 $\mathbb P^1$ 组成的链。另一方面，不难看出，$g$ 的其他每条纤维，均为 $\mathbb C^*$ 的副本。

\SourceParagraph{ch08-8.4.1-005}{628}
像任意环面簇一样，$X_\Sigma$ 有由扇的锥，等价地由 $\Delta$ 的面标记的开覆盖。特别地，有开覆盖 $\{U_w\mid w\text{ 为 }\Delta\text{ 的顶点}\}$，其中
\[
U_w=\operatorname{Spec}\mathbb C[\check w^\vee\cap\mathbb Z^2].
\]
它是对应于扇中锥 $\check w$ 的开子集。也可以将 $\check w^\vee$ 看作 $\Delta$ 在 $w$ 处的切楔，见\BookRef{fig:8.26}{图26}。同样，若 $\sigma$ 为 $\Delta$ 的边，则得到开子集 $U_\sigma=\operatorname{Spec}\mathbb C[\check\sigma^\vee\cap\mathbb Z^2]$。
\begin{figure}[H]
\centering
\BookImage{01_Numbered_Figures/Chapter_08/Ch08_Fig26_PDFp628_Bookp618.pdf}
\BookNumberOnlyCaption\label{fig:8.26}
\end{figure}

\SourceParagraph{ch08-8.4.1-006}{628}
现在，$\mathbb Z$ 在 $X_\Sigma$ 上作用，可以通过它在扇 $\Sigma$ 上的作用描述。作用生成元在 $(\mathbb R^2)^\vee$ 上为矩阵
\[
\begin{pmatrix}1&d\\0&1\end{pmatrix},
\]
将射线 $\mathbb R_{\geq0}(i,1)$ 映为 $\mathbb R_{\geq0}(i+d,1)$。在大环面轨道 $(\mathbb C^*)^2\subseteq X_\Sigma$ 上，取对应于 $\mathbb R^2$ 标准基 $(1,0),(0,1)$ 的坐标 $z,q$，它给出作用
\[
(z,q)\longmapsto(zq^d,q).
\]

\SourceParagraph{ch08-8.4.1-007}{629}
因此，固定非零 $q$ 后，作用在由 $z$ 参数化的 $\mathbb C^*$ 上为 $z\mapsto zq^d$。对于 $0<|q|<1$，此纤维的 $\mathbb Z$ 商为 $\mathbb C^*/q^{d\mathbb Z}$，等同于
\[
\mathbb C/\left\langle1,\frac d{2\pi\mathrm i}\log q\right\rangle.
\]
注意，$\exp(2\pi\mathrm i\,\cdot)$ 将 $\mathbb C/\mathbb Z$ 与 $\mathbb C^*$ 识别；只要 $|q|\neq1$，这一商显然为椭圆曲线。

\SourceParagraph{ch08-8.4.1-008}{629}
另一方面，$q=0$ 时，作用将 $\mathbb P^1$ 的无穷链平移 $d$ 个位置。因此，其 $\mathbb Z$ 商为 $d$ 条有理曲线组成的环，如\BookRef{fig:8.27}{图27}所示。
\begin{figure}[H]
\centering
\BookImage{01_Numbered_Figures/Chapter_08/Ch08_Fig27_PDFp629_Bookp619.pdf}
\BookNumberOnlyCaption\label{fig:8.27}
\end{figure}

\SourceParagraph{ch08-8.4.1-009}{629}
现在取 $\mathcal X=g^{-1}(D)/\mathbb Z$，$D$ 为开单位圆盘。由 $q$ 诱导的 $g\colon\mathcal X\to D$，为椭圆曲线向所谓 $I_d$ 型 Kodaira 纤维的退化，即向 $d$ 条有理曲线组成的环退化。

\SourceParagraph{ch08-8.4.1-010}{629}
这里，就像 Fukaya 范畴中使用作为形式幂级数环的 Novikov 环一样，采用此商的形式版本更方便。这一点技术性较强，但它使我们基本上可以完全围绕奇异纤维的拓扑工作，也使 $\mathbb Z$ 商能够在更适合代数几何的范畴中构造。

\SourceParagraph{ch08-8.4.1-011}{629}
确实，很难对代数簇直接取 $\mathbb Z$ 这类群的商。这里先注意，$g^{-1}(0)/\mathbb Z$ 作为代数簇有意义；如前所述，它为 $d$ 条有理曲线组成的环。另一方面，令 $X_\Sigma^k$ 为方程 $q^{k+1}=0$ 定义的子概形，即 $g^{-1}(0)$ 的“$k$ 阶加厚”。可以验证，$X_\Sigma^k/\mathbb Z$ 也作为概形有意义，为 $g^{-1}(0)/\mathbb Z$ 的 $k$ 阶加厚；具体细节对这里不重要。

\SourceParagraph{ch08-8.4.1-012}{630}
随后构造形式概形 $\widehat X_\Sigma$，即沿 $g^{-1}(0)$ 的形式完备化，记为 $\widehat X_\Sigma=\varprojlim X_\Sigma^k$。形式概形的定义见\BookCite{ref222}第 II 章第 9 节。它是赋环空间，底层拓扑空间为 $X_\Sigma^k$ 的拓扑空间，对所有 $k$ 均相同；等价地，就是 $g^{-1}(0)$。对于开集 $U\subseteq g^{-1}(0)$，其结构层满足
\[
\Gamma(U,\mathcal O_{\widehat X_\Sigma})
=\varprojlim_k\Gamma(U,\mathcal O_{X_\Sigma^k}).
\]
逆极限的定义以及环完备化的更多细节，见例如\BookCite{ref025}。

\SourceParagraph{ch08-8.4.1-013}{630}
本质上，这允许考虑 $g^{-1}(0)$ 的“形式邻域”中的函数，它们是关于 $q$ 的形式幂级数。例如，对于前述开集 $U_w$，$\widehat U_w=g^{-1}(0)\cap U_w$ 在拓扑上是 $g^{-1}(0)$ 的开子集，且
\[
\begin{aligned}
\Gamma(\widehat U_w,\mathcal O_{\widehat X_\Sigma})
&=\varprojlim_n\mathbb C[\check w^\vee\cap\mathbb Z^2]/(q^n)\\
&=: \mathbb C[\check w^\vee\cap\mathbb Z^2]
\widehat\otimes_{\mathbb C[q]}\mathbb C[[q]].
\end{aligned}
\]
右侧的完备张量积由左侧逆极限定义。于是，$\widehat{\mathcal X}:=\widehat X_\Sigma/\mathbb Z$ 也作为形式概形有意义，有由 $q$ 诱导的形式概形映射 $\widehat g\colon\widehat{\mathcal X}\to\widehat{\mathbb A}^1$。这里，$\widehat{\mathbb A}^1$ 为由点零及环 $\mathbb C[[q]]$ 组成的赋环空间，应看作 $g\colon\mathcal X\to D$ 的形式版本。

\SourceParagraph{ch08-8.4.1-014}{630}
实际上，将在 $\widehat{\mathcal X}$ 上计算所需的乘法映射，结果均为关于 $q$ 的形式幂级数。不过，首先需要构造线丛。

\SourceParagraph{ch08-8.4.1-015}{630}
选择 $\widehat{\mathcal X}$ 上的相对丰沛线丛。其实，$\mathcal L$ 的选择已经包含在 $\Delta$ 中：多面体不仅决定环面簇，也决定其上的线丛。这里，$X_\Sigma$ 上有线丛 $\mathcal L=\mathcal O_\Delta(1)$，其截面基与 $\Delta$ 中的整格点一一对应。可以通过环面簇线丛的通常转移映射描述它：$\mathcal L|_{U_w}$ 平凡，自然识别为 $\mathcal O_{U_w}\cdot z^w$。因此，在 $U_w\cap U_{w'}$ 上，从 $\mathcal O_{U_w}$ 到 $\mathcal O_{U_{w'}}$ 的转移映射为乘以 $z^{w-w'}$。容易看出，
\[
\Delta=\bigcap_{w\text{ 为 }\Delta\text{ 顶点}}(\check w^\vee+w),
\]
所以每个整格点 $m\in\Delta$ 都给出 $\mathcal L$ 的良好定义的截面 $z^m$。注意，将 $\mathcal L$ 限制到 $g^{-1}(0)$ 的任一不可约分支，就得到 $\mathcal O_{\mathbb P^1}(1)$。

\SourceParagraph{ch08-8.4.1-016}{630}
希望 $\mathcal L$ 在 $\widehat{\mathcal X}$ 上诱导线丛 $\widehat{\mathcal L}$。为此，仅知道 $\mathbb Z$ 在 $X_\Sigma$ 上的作用还不够，必须知道它如何提升为 $\mathcal L$ 上的作用。只需描述它在
\[
\bigoplus_{n=0}^\infty H^0(X_\Sigma,\mathcal L^{\otimes n})
\]
上的作用。此截面空间有由 $C(\Delta)\subseteq\mathbb R^3$ 的整格点给出的典范基，其中
\[
C(\Delta)=\overline{\{(rm,r)\mid m\in\Delta,\ r>0\}}.
\]
第三坐标为整数 $n$ 的整格点，对应于 $H^0(X_\Sigma,\mathcal L^{\otimes n})$ 的基元素。

\SourceParagraph{ch08-8.4.1-017}{631}
所需作用应满足下列性质。
\begin{enumerate}
\item 在 $\mathbb Z^3\subseteq\mathbb R^3$ 上线性，因为希望它通过 $\mathcal L$ 总空间的环面自同构作用。
\item 保持第三坐标，即不改变截面的次数。
\item 与 $X_\Sigma$ 上的作用相容，因此在 $\mathbb Z^2\times\{0\}\subseteq\mathbb Z^3$ 上的作用，由矩阵 $\begin{pmatrix}1&0\\d&1\end{pmatrix}$ 生成，它是 $(\mathbb R^2)^\vee$ 上作用的转置。
\item 将 $C(\Delta)$ 映到自身。
\end{enumerate}
前三个条件要求生成元为 $\begin{pmatrix}1&0&*\\d&1&*\\0&0&1\end{pmatrix}$。第四个条件进一步要求保持 $C(\Delta)$ 的边界，由周期性\BookEqRef{eq:8.12}{8.12}，唯一的选择为
\[
T=\begin{pmatrix}
1&0&d\\
d&1&\dfrac{d(d-1)}2\\
0&0&1
\end{pmatrix}.
\]

\SourceParagraph{ch08-8.4.1-018}{631}
对于给定 $n\geq0$，将此生成元写为仿射变换
\[
T_n(x,y):=\left(x+nd,xd+y+\frac{d(d-1)}2n\right),
\]
即 $T(x,y,n)$ 的前两坐标。这描述了它在 $\mathcal L^{\otimes n}$ 截面上的作用，截面基由 $n\Delta$ 的整格点代表。

\SourceParagraph{ch08-8.4.1-019}{631}
像先前一样，现对 $\mathcal L$ 取 $\mathbb Z$ 商，这只在 $g^{-1}(0)$ 的形式邻域 $\widehat X_\Sigma$ 中有意义。由此得到 $\widehat{\mathcal X}$ 上的线丛 $\widehat{\mathcal L}$。

\SourceParagraph{ch08-8.4.1-020}{631}
最简单的描述，是写出 $H^0(\widehat{\mathcal X},\widehat{\mathcal L}^{\otimes n})$ 的 $\mathbb C[[q]]$ 截面基，即形式完备化上 $\mathcal L^{\otimes n}$ 的 $\mathbb Z$ 不变截面的基。对于 $n\geq1$，任取 $p\in\tfrac1n\mathbb Z$，定义
\[
\vartheta_{n,p}:=\sum_{s=-\infty}^\infty
z^{T_n^s(np,n\varphi(p))}.
\]
注意，在任意给定开集 $U_w$ 上，此式在形式意义下成立，即
\[
\vartheta_{n,p}\in
\bigl(\mathbb C[\check w^\vee\cap\mathbb Z^2]
\widehat\otimes_{\mathbb C[q]}\mathbb C[[q]]\bigr)\cdot z^{nw}.
\]
对于任意给定 $k>0$，定义此和式的项，除有限个外，在
\[
\bigl(\mathbb C[\check w^\vee\cap\mathbb Z^2]
\otimes_{\mathbb C[q]}\mathbb C[q]/(q^k)\bigr)\cdot z^{nw}
\]
中均为零，由此即可看出其良好定义。这些函数称为 theta 函数。

\SourceParagraph{ch08-8.4.1-021}{632}
还有 $\vartheta_{n,p}=\vartheta_{n,p+d}$，因为
\[
T_n(np,n\varphi(p))=(n(p+d),n\varphi(p+d)).
\]
不难看出，各 $\vartheta_{n,p}$ 作为 $\mathbb C[[q]]$ 模生成 $\mathcal L^{\otimes n}|_{\widehat X_\Sigma}$ 的 $\mathbb Z$ 不变截面空间。不过，必须容许关于 $q$ 的形式幂级数，才能存在这种非零不变截面；对于 $n>0$，$X_\Sigma$ 上的 $\mathcal L^{\otimes n}$ 没有这样的非零截面。这再次说明，为取商，必须先作完备化。更一般地，若 $(x,y)\in\mathbb Z^2\cap n\Delta$，即 $y\geq n\varphi(x/n)$，则
\begin{equation}\label{eq:8.13}
\sum_{s=-\infty}^\infty z^{T_n^s(x,y)}
=q^{y-n\varphi(x/n)}\vartheta_{n,x/n}.
\end{equation}

\SourceParagraph{ch08-8.4.1-022}{632}
定义
\begin{equation}\label{eq:8.14}
\operatorname{ord}_n(x,y)=y-n\varphi(x/n),
\end{equation}
其中 $n\in\mathbb Z\setminus\{0\}$、$(x,y)\in\mathbb Z^2$。若 $n=0$、$x=0$，则自然取 $\operatorname{ord}_0(0,y)=y$。这一概念的动机来自\BookEqRef{eq:8.13}{8.13}。

\SourceParagraph{ch08-8.4.1-023}{632}
这一不变截面基的元素，与整仿射流形 $B=\mathbb R/d\mathbb Z$ 上坐标属于 $\tfrac1n\mathbb Z$ 的点一一对应，见\BookSectionRef{sec:6.1.2}{6.1.2}；马上还将回到这一点。将这些点的集合记为 $B((1/n)\mathbb Z)$。

\SourceParagraph{ch08-8.4.1-024}{632}
若读者不满意上述较粗略的商构造描述，也可以如下处理。考虑 $\mathbb C[[q]]$ 代数
\[
\widehat R=\bigoplus_{n=0}^\infty
H^0(\widehat{\mathcal X},\widehat{\mathcal L}^{\otimes n}).
\]
仅需知道其 $\mathbb C[[q]]$ 模生成元集合
\[
\{1\}\cup\{\vartheta_{n,p}\mid p\in B((1/n)\mathbb Z),\ n\geq1\},
\]
以及乘法映射，也就是代数的结构常数，即可定义它。乘法马上将说明。随后考虑概形 $\operatorname{Proj}\widehat R$，它定义在 $\operatorname{Spec}\mathbb C[[q]]$ 上，参见\BookCite{ref222}第 II 章第 2 节。这是通常意义的概形，而非形式概形。存在丰沛线丛时，可以从形式概形得到通常概形，这正是 Grothendieck 存在性定理，见 EGA III，5.4.5\BookCite{ref211}。所以，上述构造也可看作这一一般结论的特例。

\subsection{theta 函数的乘法与三角形}\label{sec:8.4.2}

\SourceParagraph{ch08-8.4.2-001}{633}
现在转向 theta 函数及其乘法的几何解释。前面对 theta 函数基的描述，已经产生一个仿射流形 $B$，因此，自然应尝试将当前构造放入\BookChapterRef{6}的框架。

\SourceParagraph{ch08-8.4.2-002}{633}
实际上，将 $\mathcal X\to D$ 的纤维看作 $\check X(B)$，见\BookSectionRef{sec:6.2.1}{6.2.1}。纤维带有线丛 $\mathcal L$ 的限制，其次数必为 $d$。另一方面，$\check X(B)$ 有典范辛形式 $\omega$，满足 $\int_{\check X(B)}\omega=d$，因此 $\omega$ 代表 $c_1(\mathcal L)$。

\SourceParagraph{ch08-8.4.2-003}{633}
读者或许认为，这与\BookSectionRef{sec:6.3}{6.3}的思路相反：那里，在 $X(B)$ 上考虑线丛，在 $\check X(B)$ 上考虑 Lagrange 子流形。确实如此，不过这一交换有其原因。讨论 Fukaya 范畴时，至少在椭圆曲线情形，可以使用固定的概复结构，即 $X(B)$ 的复结构；它的全纯圆盘很容易描述。选择在 Fukaya 一侧使用 $X(B)$，使圆盘易于描述，却使 Lagrange 子流形较难描述。同样，在导出范畴一侧使用 $\check X(B)$，层较难描述，却从辛形式获得一个正则给定的上同调类。因此，只考虑第一 Chern 类由 $\omega$ 的整倍数代表的线丛，更具体地，只考虑 $\mathcal L$ 的幂。在 $X(B)$ 上，则只考虑非常特殊的 Lagrange 子流形：由展开映射的倍数诱导的截面，见\BookSectionRef{sec:6.1.2}{6.1.2}，其局部系统平凡。

\SourceParagraph{ch08-8.4.2-004}{633}
按照\BookSectionRef{sec:6.3}{6.3}的思路，更显式地描述与 $\mathcal L^{\otimes n}$ 镜像对偶的 $X(B)$ 上的 Lagrange 子流形。首先，指定 $X(B)$ 的辛形式。由\BookRef{prop:6.14}{命题6.14}，可以通过 $B$ 上的凸多值函数指定；取 $K=y^2/2$ 就很好，得到 $\omega=\mathrm dy\wedge\mathrm dx$。这里暂时沿用第 6 章的约定，以 $y$ 为 $B$ 的坐标，$x$ 为 $X(B)$ 的纤维坐标。任何实曲线都是 Lagrange 的，所以 $\omega$ 在这里唯一的作用，是计算全纯圆盘的面积。

\SourceParagraph{ch08-8.4.2-005}{633}
对于任意 $n\in\mathbb Z$，令 $L_n$ 为图像 $x=-ny$ 在 $X(B)$ 中的像，沿 $y$ 增加的方向定向。负号来自 Fukaya 范畴的符号约定。由于 $x$ 的周期为一、$y$ 的周期为 $d$，这一图像定义良好，见\BookRef{fig:8.28}{图28}。首先注意，若 $f\colon X(B)\to B$ 为投影，$n_1\neq n_2$，则
\[
f(L_{n_1}\cap L_{n_2})
=f(L_0\cap L_{n_2-n_1})
=B\left(\frac1{n_2-n_1}\mathbb Z\right).
\]
由\BookSectionRef{sec:6.3}{6.3}，期待 $L_n$ 与次数为 $nd$ 的线丛互为镜像；至少除一个符号外如此，符号由 Fukaya 范畴的约定决定，这里并非关键。因此，自然期待 $\widehat{\mathcal L}^{\otimes n}$ 与 $L_n$ 互为镜像。实际上，作为理解这一对偶的第一步，已经在上同调层面，找到两侧间的一个自然同构。
\begin{figure}[H]
\centering
\BookImage{01_Numbered_Figures/Chapter_08/Ch08_Fig28_PDFp634_Bookp624.pdf}
\BookNumberOnlyCaption\label{fig:8.28}
\end{figure}

\SourceParagraph{ch08-8.4.2-006}{634}
更准确地，设 $n_1<n_2$。则
\[
\begin{aligned}
\operatorname{Ext}^i(\widehat{\mathcal L}^{\otimes n_1},
\widehat{\mathcal L}^{\otimes n_2})
&=H^i(\widehat{\mathcal X},\widehat{\mathcal L}^{\otimes(n_2-n_1)})\\
&=\begin{cases}
\displaystyle\bigoplus_{p\in B(\frac1{n_2-n_1}\mathbb Z)}[p]\mathbb C[[q]],&i=0,\\
0,&i=1.
\end{cases}
\end{aligned}
\]
特别地，复形 $\operatorname{Hom}_{D^b_\infty(\widehat{\mathcal X})}
(\widehat{\mathcal L}^{\otimes n_1},\widehat{\mathcal L}^{\otimes n_2})$，拟同构于
\begin{equation}\label{eq:8.15}
\cdots\longrightarrow0\longrightarrow
\bigoplus_{p\in B(\frac1{n_2-n_1}\mathbb Z)}[p]\mathbb C[[q]]
\longrightarrow0\longrightarrow\cdots,
\end{equation}
其中非零项位于次数零。

\SourceParagraph{ch08-8.4.2-007}{634}
另一方面，若 $n_1>n_2$，则
\[
\begin{aligned}
\operatorname{Ext}^i(\widehat{\mathcal L}^{\otimes n_1},
\widehat{\mathcal L}^{\otimes n_2})
&=H^i(\widehat{\mathcal X},\widehat{\mathcal L}^{\otimes(n_2-n_1)})\\
&=\begin{cases}
0,&i=0,\\
\displaystyle\bigoplus_{p\in B(\frac1{n_2-n_1}\mathbb Z)}[p]\mathbb C[[q]],&i=1.
\end{cases}
\end{aligned}
\]
最后一个等式来自 Serre 对偶，后面将给出更显式的描述。因此，态射复形拟同构于仅在次数一非零的复形
\begin{equation}\label{eq:8.16}
\cdots\longrightarrow0\longrightarrow
\bigoplus_{p\in B(\frac1{n_2-n_1}\mathbb Z)}[p]\mathbb C[[q]]
\longrightarrow0\longrightarrow\cdots.
\end{equation}

\SourceParagraph{ch08-8.4.2-008}{634}
与 Fukaya 一侧比较，需要在 $L_n$ 上选择分次。它们在 $X(B)$ 的万有覆盖 $\mathbb R^2$ 中为直线，因此可以取常数分次。按照\BookRef{fig:8.18}{图18}，根据直线的角度，在 $(-1,1)$ 中选择相应分次。于是，由图18可见，$n_1<n_2$ 时，$L_{n_1}\cap L_{n_2}$ 的所有点指标为零；$n_1>n_2$ 时，指标为一。它们当然与 $B(\tfrac1{n_2-n_1}\mathbb Z)$ 一一对应。因此，张量上 $\Lambda_{\mathrm{nov}}$ 后，这两种情形的 Floer 复形，分别与\BookEqRef{eq:8.15}{8.15}、\BookEqRef{eq:8.16}{8.16}一致。

\SourceParagraph{ch08-8.4.2-009}{635}
现在可以比较乘法映射。首先只在 $n_1,n_2>0$ 时，考虑复合
\[
\begin{aligned}
\operatorname{Hom}(\widehat{\mathcal L}^{\otimes(n_0+n_1)},
\widehat{\mathcal L}^{\otimes(n_0+n_1+n_2)})
\otimes\operatorname{Hom}(\widehat{\mathcal L}^{\otimes n_0},
\widehat{\mathcal L}^{\otimes(n_0+n_1)})\\
\longrightarrow
\operatorname{Hom}(\widehat{\mathcal L}^{\otimes n_0},
\widehat{\mathcal L}^{\otimes(n_0+n_1+n_2)}).
\end{aligned}
\]
因为这时可以使用 theta 函数的显式描述。引入 Čech 复形计算高阶乘法映射后，再处理其余情形。

\SourceParagraph{ch08-8.4.2-010}{635}
不妨取 $n_0=0$，要计算的乘积便为
\[
H^0(\widehat{\mathcal L}^{\otimes n_1})
\otimes H^0(\widehat{\mathcal L}^{\otimes n_2})
\longrightarrow H^0(\widehat{\mathcal L}^{\otimes(n_1+n_2)}).
\]
取 $p_1\in B(\tfrac1{n_1}\mathbb Z)$、$p_2\in B(\tfrac1{n_2}\mathbb Z)$，选择其在 $\mathbb R$ 上的提升。暂记
\[
N=n_1+n_2,\qquad
P_\alpha=\frac{n_1p_1+n_2(p_2+\alpha d)}N,
\qquad\alpha=s_2-s_1.
\]
则逐步计算得
\[
\begin{aligned}
\vartheta_{n_1,p_1}\vartheta_{n_2,p_2}
&=\left(\sum_{s_1\in\mathbb Z}z^{T_{n_1}^{s_1}(n_1p_1,n_1\varphi(p_1))}\right)
\left(\sum_{s_2\in\mathbb Z}z^{T_{n_2}^{s_2}(n_2p_2,n_2\varphi(p_2))}\right)\\
&=\sum_{s_1,s_2\in\mathbb Z}
z^{(n_1(p_1+s_1d),n_1\varphi(p_1+s_1d))}
z^{(n_2(p_2+s_2d),n_2\varphi(p_2+s_2d))}\\
&=\sum_{s_1,s_2\in\mathbb Z}
z^{(N(P_\alpha+s_1d),\ n_1\varphi(p_1+s_1d)+n_2\varphi(p_2+s_2d))}\\
&=\sum_{s_1,s_2\in\mathbb Z}
z^{(N(P_\alpha+s_1d),\ N\varphi(P_\alpha+s_1d))}\\
&\qquad\cdot
q^{n_1\varphi(p_1+s_1d)+n_2\varphi(p_2+s_2d)-N\varphi(P_\alpha+s_1d)}.
\end{aligned}
\]

\SourceParagraph{ch08-8.4.2-011}{635}
利用\BookEqRef{eq:8.12}{8.12}，可以验证
\[
\begin{aligned}
&n_1\varphi(p_1+s_1d)+n_2\varphi(p_2+s_2d)
-N\varphi(P_{s_2-s_1}+s_1d)\\
&\qquad=n_1\varphi(p_1)+n_2\varphi(p_2+(s_2-s_1)d)
-N\varphi(P_{s_2-s_1}).
\end{aligned}
\]
对于 $p_1\in\tfrac1{n_1}\mathbb Z$、$p_2\in\tfrac1{n_2}\mathbb Z$，定义
\[
\deg(p_1,p_2)=n_1\varphi(p_1)+n_2\varphi(p_2)
-(n_1+n_2)\varphi\left(\frac{n_1p_1+n_2p_2}{n_1+n_2}\right).
\]

\SourceParagraph{ch08-8.4.2-012}{636}
由 $\varphi$ 的凸性，此数始终非负。因此，
\[
\begin{aligned}
\vartheta_{n_1,p_1}\vartheta_{n_2,p_2}
&=\sum_{\alpha\in\mathbb Z}\sum_{s\in\mathbb Z}
z^{T_N^s(NP_\alpha,N\varphi(P_\alpha))}
q^{\deg(p_1,p_2+\alpha d)}\\
&=\sum_{\alpha\in\mathbb Z}
\vartheta_{N,P_\alpha}q^{\deg(p_1,p_2+\alpha d)}.
\end{aligned}
\]

\SourceParagraph{ch08-8.4.2-013}{636}
几何上，可以如下解释。$p_1,p_2$ 被看作 $B$ 的万有覆盖 $\mathbb R$ 上的点，只定义到模 $d$。因此，对于 $\alpha\in\mathbb Z$，可以取提升 $p_1,p_2+\alpha d$，连接它们的有向区间映到 $B$ 上连接 $p_1,p_2$ 的仿射线段。对所有 $\alpha$ 求和，就相当于对端点为 $p_1\bmod d\mathbb Z$、$p_2\bmod d\mathbb Z$ 的全部仿射线段求和。注意，每条线段的贡献位于一个 $B$ 中的点，它是提升端点的加权平均。

\SourceParagraph{ch08-8.4.2-014}{636}
\BookStatementNumber{18}
\begin{note}[numbered]\label{note:8.18}
这是射影环面簇的一个非常简单现象的推广。若紧格多胞体 $\Delta$ 定义射影环面簇 $(\mathbb P_\Delta,\mathcal O_{\mathbb P_\Delta}(1))$，则 $\Delta(\tfrac1n\mathbb Z)$ 的点，构成 $H^0(\mathbb P_\Delta,\mathcal O_{\mathbb P_\Delta}(n))$ 的一组基。两个基截面 $m_1\in\Delta(\tfrac1{n_1}\mathbb Z)$、$m_2\in\Delta(\tfrac1{n_2}\mathbb Z)$ 的乘积，给出加权平均
\[
\frac{n_1m_1+n_2m_2}{n_1+n_2}
\in\Delta\left(\frac1{n_1+n_2}\mathbb Z\right).
\]
\end{note}

\SourceParagraph{ch08-8.4.2-015}{636}
现在可以与 Fukaya 一侧联系。不妨取 $n_0=0$，考虑乘法映射
\[
\begin{aligned}
\operatorname{Hom}_{\operatorname{Fuk}(\check X)}(L_{n_1},L_{n_1+n_2})
\otimes\operatorname{Hom}_{\operatorname{Fuk}(\check X)}(L_0,L_{n_1})\\
\longrightarrow\operatorname{Hom}_{\operatorname{Fuk}(\check X)}(L_0,L_{n_1+n_2}).
\end{aligned}
\]
如前所述，$p_1,p_2$ 分别确定从 $L_0$ 到 $L_{n_1}$、从 $L_{n_1}$ 到 $L_{n_1+n_2}$ 的态射。给定 $\alpha$，将三条 Lagrange 曲线提升到 $\check X=X(B)=\mathbb C/\langle d,\mathrm i\rangle$ 的万有覆盖 $\mathbb C$，如\BookRef{fig:8.29}{图29}所示。加权平均出现于 $\widetilde L_0\cap\widetilde L_{n_1+n_2}$。因此，阴影三角形 $T$，作为全纯圆盘，对 Floer 乘积的贡献为 $q^{\int_T\omega}$，按照\BookSectionRef{sec:8.3.3}{8.3.3}，其符号为正。图中的三角形面积为
\[
\begin{aligned}
\int_T\omega
&=\frac12\left(\frac{n_1p_1+n_2(p_2+\alpha d)}{n_1+n_2}-p_1\right)
n_1(p_2+\alpha d-p_1)\\
&=\frac12\frac{n_1n_2((p_2+\alpha d)-p_1)^2}{n_1+n_2}\\
&=n_1\psi(p_1)+n_2\psi(p_2+\alpha d)
-(n_1+n_2)\psi\left(\frac{n_1p_1+n_2(p_2+\alpha d)}{n_1+n_2}\right),
\end{aligned}
\]
其中
\begin{equation}\label{eq:8.17}
\psi(x)=\frac{x(x-1)}2.
\end{equation}
这很容易验证。
\begin{figure}[H]
\centering
\BookImage{01_Numbered_Figures/Chapter_08/Ch08_Fig29_PDFp637_Bookp627.pdf}
\BookNumberOnlyCaption\label{fig:8.29}
\end{figure}

\SourceParagraph{ch08-8.4.2-016}{637}
此面积与 $\deg(p_1,p_2+\alpha d)$ 不同，乍看并未得到相同贡献。不过，令
\begin{equation}\label{eq:8.18}
A(x)=\psi(x)-\varphi(x),
\end{equation}
则 $A(x+d)=A(x)$。对于任意 $n_1\neq n_2$，可以在态射的上同调层面定义映射
\[
\begin{aligned}
\operatorname{Ext}^*(\widehat{\mathcal L}^{\otimes n_1},
\widehat{\mathcal L}^{\otimes n_2})\otimes_{\mathbb C[[q]]}\Lambda_{\mathrm{nov}}
\longrightarrow\operatorname{Hom}_{\operatorname{Fuk}(X(B))}(L_{n_1},L_{n_2}),
\end{aligned}
\]
它对 $p\in B(\tfrac1{n_2-n_1}\mathbb Z)$，给出
\begin{equation}\label{eq:8.19}
[p]\longmapsto[p]q^{-(n_2-n_1)A(p)}.
\end{equation}
这就解决了问题。在上述
\[
H^0(\mathcal L^{\otimes n_1})\otimes H^0(\mathcal L^{\otimes n_2})
\longrightarrow H^0(\mathcal L^{\otimes(n_1+n_2)})
\]
的乘积中，$[p_1]$ 与 $[p_2]$ 的乘积在 $[P_\alpha]$ 处的系数为 $q^{\deg(p_1,p_2+\alpha d)}$，其中仍取 $P_\alpha=\tfrac{n_1p_1+n_2(p_2+\alpha d)}{n_1+n_2}$。

\SourceParagraph{ch08-8.4.2-017}{638}
通过\BookEqRef{eq:8.19}{8.19}，将此乘积输送到
\[
\begin{aligned}
\operatorname{Hom}_{\operatorname{Fuk}(X(B))}^0(L_{n_1},L_{n_1+n_2})
\otimes\operatorname{Hom}_{\operatorname{Fuk}(X(B))}^0(L_0,L_{n_1})\\
\longrightarrow\operatorname{Hom}_{\operatorname{Fuk}(X(B))}^0(L_0,L_{n_1+n_2}),
\end{aligned}
\]
则 $[p_2+\alpha d]\otimes1$ 与 $[p_1]\otimes1$ 复合所得 $[P_\alpha]$ 的系数，为
\[
q^{n_1A(p_1)}q^{n_2A(p_2)}q^{\deg(p_1,p_2+\alpha d)}
q^{-(n_1+n_2)A(P_\alpha)}=q^{\operatorname{Area}(T)}.
\]
因此，theta 函数的乘法与 Floer 乘法一致。

\SourceParagraph{ch08-8.4.2-018}{638}
现在希望理解高阶乘积，为此必须具体使用 $D^b_\infty(\widehat{\mathcal X})$ 的 Čech 描述。不过，先给出 Fukaya 一侧的几何描述，使其与上述乘法解释相契合。

\subsection{热带 Morse 树}\label{sec:8.4.3}

\SourceParagraph{ch08-8.4.3-001}{638}
Fukaya 与 Oh 的一篇重要论文\BookCite{ref167}，将流形余切丛中截面的 Floer 同调，与该流形的 Morse 同调联系起来。粗略地说，其思想是用梯度流线产生的结构来逼近全纯圆盘。Kontsevich 与 Soibelman 在\BookCite{ref311}中利用这一思想，证明了 Abel 簇同调镜像对称的一个版本。对于椭圆曲线，这一步尤其简单；我们将采用这一思想的一个变体，定义所谓的\emph{热带 Morse 树}。以下叙述遵循 Abouzaid、Gross 与 Siebert 的合作工作\BookCite{ref001}。

\SourceParagraph{ch08-8.4.3-002}{638}
\BookStatementNumber{19}
\begin{definition}\label{def:8.19}
设 $B$ 为整仿射流形。给定一列互不相同的整数 $n_0,\ldots,n_d\in\mathbb Z$，以及任意度量带状树 $S$，可以如下用整数 $n_e$ 标记 $S$ 的各条边 $e$。若 $e$ 为与第 $i$ 个外部顶点相接的外部入边，则 $n_e=n_i-n_{i-1}$；否则，若 $e$ 从顶点 $v$ 出发，则 $n_e$ 为所有进入 $v$ 的边的标记之和。再给定点
\[
p_{i,i+1}\in B\left(\frac1{n_{i+1}-n_i}\mathbb Z\right),
\qquad
p_{0,d}\in B\left(\frac1{n_d-n_0}\mathbb Z\right),
\]
定义
\[
S_d^{\mathrm{trop}}(p_{0,d};p_{0,1},\ldots,p_{d-1,d})
\]
为 $B$ 上\emph{热带 Morse 树}的模空间。所谓热带 Morse 树，是从带状树出发的连续映射 $\phi\colon S\to B$，附带一组\emph{仿射位移向量}：即对 $S$ 的每条边 $e$，有截面 $v_e\in\Gamma(e,(\phi|_e)^*TB)$，满足下列性质。
\begin{enumerate}
\item 若 $v$ 为第 $i$ 个外部入顶点，则 $\phi(v)=p_{i-1,i}$；若 $v$ 为外部出顶点，则 $\phi(v)=p_{0,d}$。
\item 若 $e$ 为 $S$ 的一条边，则 $\phi(e)$ 在 $B$ 上局部为仿射线段。这条线段的斜率可以是无理数。若 $e$ 为外部边，也允许 $\phi(e)$ 为一点。
\SourceParagraph{ch08-8.4.3-003}{639}
\item 若 $v$ 为外部顶点，而 $e$ 为 $S$ 中唯一包含 $v$ 的边，则 $v_e(v)=0$。
\item 将 $S$ 的一条边 $e$ 与坐标为 $s$ 的 $[0,1]$ 识别，并将 $e$ 定向为从零到一。则 $v_e(s)$ 在 $\phi(s)$ 处与 $\phi(e)$ 相切，并且指向由 $e$ 的定向在 $\phi(e)$ 上诱导的方向。此外，利用仿射结构将 $(\phi|_e)^*TB$ 与 $e$ 上的平凡丛识别后，有
\[
\frac{d}{ds}v_e(s)=n_e\phi_*(\partial/\partial s).
\]
\item 若 $v$ 为 $S$ 的内部顶点，入边为 $e_1,\ldots,e_p$，出边为 $e_{\mathrm{out}}$，则
\[
v_{e_{\mathrm{out}}}(v)=\sum_{i=1}^p v_{e_i}(v).
\]
\item $S$ 中边 $e$ 的长度，等于
\[
\frac1{n_e}\log\left(\frac{v_e(1)}{v_e(0)}\right).
\]
这里应记住，带状树的每条边都带有长度，外部边的长度为无穷。由于 $v_e(0)$ 与 $v_e(1)$ 是同方向的成比例向量，所以其比值具有意义，是一个正数。有一种特殊情形：若 $e$ 为被 $\phi$ 收缩的外部边，则 $v_e(0)=v_e(1)=0$，但仍将其长度取为无穷。
\end{enumerate}
\end{definition}

\SourceParagraph{ch08-8.4.3-004}{639}
\BookStatementNumber{20}
\begin{note}[numbered]\label{note:8.20}
在 $B=\mathbb R^n$ 的情形，可以更清楚地解释条件(4)的含义：它告诉我们
\[
v_e(s)=v_e(0)+n_e\bigl(\phi(s)-\phi(0)\bigr).
\]
\SourceParagraph{ch08-8.4.3-005}{639}
条件(6)在这一定义中实际上没有发挥重要作用，只是给出从 $S_d^{\mathrm{trop}}(p_{0,d};p_{0,1},\ldots,p_{d-1,d})$ 到 $S_d$ 的映射。不过，稍后构造范畴的拟同构时，它将发挥更重要的作用。
\end{note}

\SourceParagraph{ch08-8.4.3-006}{639}
\BookStatementNumber{21}
\begin{exercise}\label{exer:8.21}
可以通过以下练习理解这一定义的各个条件；练习将说明，热带 Morse 树的概念本质上只是梯度树概念的 Legendre 变换。假设 $B$ 是由 $(\check B,\check K)$ 作 Legendre 变换得到的，如\BookSectionRef{sec:6.1.2}{6.1.2}所述。将 $\check K$ 看作 $\check B$ 的万有覆盖上的函数。给定 $n_0,\ldots,n_d\in\mathbb Z$，证明：$B$ 上的热带 Morse 树，恰好对应于 $\check B$ 上函数
\[
(n_0\check K+l_0,\ldots,n_d\check K+l_d)
\]
的一棵梯度树，其中度量由 $\check K$ 诱导，$l_0,\ldots,l_d$ 为 $\check B$ 的万有覆盖上某些斜率为整向量的仿射线性函数。此外，每条边上梯度流所经历的时间长度，恰好由定义的条件(6)给出。
\end{exercise}

\SourceParagraph{ch08-8.4.3-007}{640}
\BookStatementNumber{22}
\begin{example}\label{ex:8.22}
前面关于 theta 函数乘法的描述，已经给出了热带 Morse 树的例子。设 $B=\mathbb R/d\mathbb Z$，$T$ 为具有两条入边的带状树，且 $n_1,n_2>0$。则\BookRef{fig:8.30}{图30}给出一棵热带 Morse 树，表示 $\vartheta_{n_1,p_{0,1}}$ 与 $\vartheta_{n_2,p_{1,2}}$ 的乘积对 $\vartheta_{n_3,p_{0,2}}$ 的一项贡献，其中 $n_3=n_1+n_2$。注意，外部出边必须收缩为一点，否则仿射位移向量不可能在出顶点处为零。

\SourceParagraph{ch08-8.4.3-008}{640}
另一方面，若 $n_1>0$、$n_2<0$，且 $n_3=n_1+n_2<0$，则\BookRef{fig:8.31}{图31}给出了热带 Morse 树的另一个例子。此时，根据条件(4)，带负权的入边必须收缩为一点。注意，两种情形中的树映射，都可以看作经由万有覆盖分解；所以 $p_{0,1}$、$p_{1,2}$ 与 $p_{0,2}$ 可以是 $B$ 中相应点在 $\mathbb R$ 中的任意提升。
\setcounter{figure}{29}
\begin{figure}[H]
\centering
\BookImage{01_Numbered_Figures/Chapter_08/Ch08_Fig30_PDFp640_Bookp630.pdf}
\caption{一棵热带 Morse 树。定义中条件(3)与(5)要求 $n_1(p_{0,2}-p_{0,1})+n_2(p_{0,2}-p_{1,2})=0$；等价地，$p_{0,2}=(n_1p_{0,1}+n_2p_{1,2})/(n_1+n_2)$。}\label{fig:8.30}
\end{figure}
\begin{figure}[H]
\centering
\BookImage{01_Numbered_Figures/Chapter_08/Ch08_Fig31_PDFp640_Bookp630.pdf}
\caption{$v_e$ 在 $p_{0,2}$ 处为零的条件，同样推出 $p_{0,2}=(n_1p_{0,1}+n_2p_{1,2})/(n_1+n_2)$。}\label{fig:8.31}
\end{figure}
\end{example}

\subsection{全纯多边形与热带 Morse 树}\label{sec:8.4.4}

\SourceParagraph{ch08-8.4.4-001}{640}
现在解释，一棵热带 Morse 树如何产生 $X(B)$ 中若干多边形的并。若 $\dim B\geq2$，这个并未必是真正的全纯子流形，因为它是分段线性对象；但若 $\dim B=1$，它总是一个全纯多边形。

\SourceParagraph{ch08-8.4.4-002}{640}
当 $B=\mathbb R/d\mathbb Z$ 时，$X(B)\to B$ 的 Lagrange 子流形 $L_n$，可以推广到任意 $r$ 维整仿射流形 $B$。取 $B$ 上的局部整仿射坐标 $y_1,\ldots,y_r$，以及 $TB$ 上如\BookSectionRef{sec:6.2.1}{6.2.1}所述的纤维坐标 $x_1,\ldots,x_r$，便可以用局部公式
\[
\sigma_n(y_1,\ldots,y_r)=(y_1,\ldots,y_r,-ny_1,\ldots,-ny_r)
\]
定义截面 $\sigma_n\colon B\to X(B)$。
\SourceParagraph{ch08-8.4.4-003}{641}
容易验证，在整仿射坐标变换下，这一定义是良定义的。令
\[
L_n=\sigma_n(B).
\]
在这个一般情形下，我们实际上并未在 $X(B)$ 上赋予辛结构，因而谈论 $L_n$ 是否为 Lagrange 子流形并无意义。

\SourceParagraph{ch08-8.4.4-004}{641}
下面说明，热带 Morse 树 $\phi\colon S\to B$ 如何定义一个分段线性圆盘。$S$ 的任意边 $e$，都标记为某个 $j>i$ 对应的 $n_e=n_j-n_i$。考虑映射
\[
R_e\colon e\times[0,1]\longrightarrow X(B),
\qquad
R_e(s,t)=\sigma_{n_i}(\phi(s))-t\,v_e(s),
\]
其中 $t\,v_e(s)$ 看作 $\phi(s)$ 处的切向量。将 $e$ 的两个顶点记为 $v_{\mathrm{in}}$ 与 $v_{\mathrm{out}}$。根据定义，有
\[
R_e(e\times\{0\})\subseteq L_{n_i}.
\]
若
\[
R_e(v_{\mathrm{in}}\times\{1\})\subseteq L_{n_j},
\]
即
\[
-v_e(v_{\mathrm{in}})
=\sigma_{n_j}(\phi(v_{\mathrm{in}}))
-\sigma_{n_i}(\phi(v_{\mathrm{in}}))\pmod\Lambda,
\]
则由\BookRef{def:8.19}{定义8.19}的条件(4)，模 $\Lambda$ 有
\[
\begin{aligned}
-\frac{d}{ds}v_e(s)
&=-(n_j-n_i)(\phi|_e)_*\frac\partial{\partial s}\\
&=\frac{d}{ds}\bigl(\sigma_{n_j}(\phi(s))-\sigma_{n_i}(\phi(s))\bigr).
\end{aligned}
\]
因此，$R_e(v_{\mathrm{in}}\times\{1\})\subseteq L_{n_j}$ 推出 $R_e(e\times\{1\})\subseteq L_{n_j}$。

\SourceParagraph{ch08-8.4.4-005}{641}
再考察 $S$ 的一个顶点 $v$。若 $v$ 为外部入顶点 $v_i$，则
\[
\phi(v_i)\in B\left(\frac1{n_i-n_{i-1}}\mathbb Z\right),
\]
所以
\[
\sigma_{n_i}(\phi(v_i))=\sigma_{n_{i-1}}(\phi(v_i))
\in L_{n_{i-1}}\cap L_{n_i}.
\]
特别地，$R_e(v_i\times\{1\})\subseteq L_{n_i}$，从而对于与 $v_i$ 相接的唯一一条边 $e$，有
\[
R_e(e\times\{0\})\subseteq L_{n_{i-1}},
\qquad
R_e(e\times\{1\})\subseteq L_{n_i}.
\]

\SourceParagraph{ch08-8.4.4-006}{641}
若 $v$ 为任意内部顶点，入边为 $e_1,\ldots,e_p$，出边为 $e_{\mathrm{out}}$，且 $e_j$ 的权为 $n_{i_j}-n_{i_{j-1}}$，其中 $i_0<\cdots<i_p$，则可以归纳地假设
\[
R_{e_j}(e_j\times\{0\})\subseteq L_{n_{i_{j-1}}},
\qquad
R_{e_j}(e_j\times\{1\})\subseteq L_{n_{i_j}}.
\]
\SourceParagraph{ch08-8.4.4-007}{642}
\BookRef{def:8.19}{定义8.19}的条件(5)告诉我们，实际上
\[
\begin{aligned}
R_{e_{\mathrm{out}}}(v\times\{1\})
&=\sigma_{n_{i_0}}(\phi(v))-v_{e_{\mathrm{out}}}(v)\\
&=\sigma_{n_{i_0}}(\phi(v))-\sum_{j=1}^p v_{e_j}(v)\\
&=\sigma_{n_{i_p}}(\phi(v))\in L_{n_{i_p}}.
\end{aligned}
\]
因此
\[
R_{e_{\mathrm{out}}}(e_{\mathrm{out}}\times\{0\})\subseteq L_{n_{i_0}},
\qquad
R_{e_{\mathrm{out}}}(e_{\mathrm{out}}\times\{1\})\subseteq L_{n_{i_p}}.
\]
在顶点 $v$ 处，矩形 $R_{e_i}(e_i\times[0,1])$（$1\leq i\leq p$）与 $R_{e_{\mathrm{out}}}(e_{\mathrm{out}}\times[0,1])$ 未必恰好拼接在一起。不过，在 $T_{\phi(v)}B$ 中，有一个多边形 $P_v$，其依次定向的边向量为 $-v_{e_1}(v),\ldots,-v_{e_p}(v)$ 与 $v_{e_{\mathrm{out}}}(v)$，如\BookRef{fig:8.32}{图32}所示。

\SourceParagraph{ch08-8.4.4-008}{642}
将矩形 $e\times[0,1]$ 与多边形区域 $P_v$ 拼接起来，形成树 $S$ 的一个加厚，便得到拓扑圆盘 $\operatorname{Fat}(S)$，见\BookRef{fig:8.33}{图33}。对于外部边，用三角形代替矩形可能更直观，因为这样的矩形中，对应于外部顶点的那条边总是被收缩。也可以略去被收缩的外部边。将映射 $R_e$ 拼接起来，就得到连续映射 $R\colon\operatorname{Fat}(S)\to X(B)$，也就是从圆盘到 $X(B)$ 的映射；这里不拘泥于拼接的具体细节。根据\BookSectionRef{sec:6.3}{6.3}中的观察，$R$ 的像由全纯三角形、全纯四边形，以及位于 $X(B)\to B$ 的纤维中的非全纯多边形组成。
\begin{figure}[H]
\centering
\BookImage{01_Numbered_Figures/Chapter_08/Ch08_Fig32_PDFp642_Bookp632.pdf}
\BookNumberOnlyCaption\label{fig:8.32}
\end{figure}
\begin{figure}[H]
\centering
\BookImage{01_Numbered_Figures/Chapter_08/Ch08_Fig33_PDFp643_Bookp633.pdf}
\BookNumberOnlyCaption\label{fig:8.33}
\end{figure}

\SourceParagraph{ch08-8.4.4-009}{642}
\BookStatementNumber{23}
\begin{example}\label{ex:8.23}
在椭圆曲线上应用上述构造，总能得到真正的全纯多边形，并将其表示为三角形与四边形的并。此时不存在二维的 $P_v$。例如，\BookRef{fig:8.30}{图30}与\BookRef{fig:8.31}{图31}都给出三角形，前者对应于\BookRef{fig:8.29}{图29}中的三角形。

\SourceParagraph{ch08-8.4.4-010}{643}
反过来，在椭圆曲线的情形，给定任意以 Lagrange 子流形 $L_{n_{i_0}},\ldots,L_{n_{i_d}}$ 为边界的全纯多边形，都能构造一棵产生该多边形的热带 Morse 树。具体细节留给读者，\BookRef{fig:8.34}{图34}与\BookRef{fig:8.35}{图35}给出了一些例子。
\begin{figure}[H]
\centering
\BookImage{01_Numbered_Figures/Chapter_08/Ch08_Fig34_PDFp644_Bookp634.pdf}
\caption{以 $p_{0,1}$ 与 $p_{1,2}$ 为顶点的外部边，其映像长度为零。}\label{fig:8.34}
\end{figure}
\begin{figure}[H]
\centering
\BookImage{01_Numbered_Figures/Chapter_08/Ch08_Fig35_PDFp644_Bookp634.pdf}
\caption{左图中，$S$ 有一个四价顶点，从 $p_{1,2}$ 出发的边被收缩。Morse 树只作示意性绘制，因为这些边在 $B$ 上彼此重叠。这棵 Morse 树是一族单参数树中的一个元素；右图给出另一个元素，以及相应的全纯多边形。}\label{fig:8.35}
\end{figure}
\end{example}

\subsection{\v{C}ech 计算}\label{sec:8.4.5}

\SourceParagraph{ch08-8.4.5-001}{643}
现在可以陈述本节的主要结果。考虑如下定义的 $A_\infty$ 预范畴 $D(\widehat X)$ 与 $\check D(X(B))$。$D(\widehat X)$ 的对象，是 $\widehat X$ 上线丛 $\widehat{\mathcal L}$ 的任意张量幂。若对任意 $i\ne j$ 都有 $n_i\ne n_j$，则称线丛列 $(\widehat{\mathcal L}^{\otimes n_1},\ldots,\widehat{\mathcal L}^{\otimes n_d})$ 为横截的。$\operatorname{Hom}_{D(\widehat X)}(\widehat{\mathcal L}^{\otimes n_1},\widehat{\mathcal L}^{\otimes n_2})$ 取为计算 $\widehat{\mathcal L}^{\otimes(n_2-n_1)}$ 上同调的一个 \v{C}ech 复形。因此，$m_1^{D(\widehat X)}$ 为 \v{C}ech 微分，$m_2^{D(\widehat X)}$ 则为\BookSectionRef{sec:8.2.1}{8.2.1}所述的复合；所有高阶乘法映射都为零。

\SourceParagraph{ch08-8.4.5-002}{643}
$\check D(X(B))$ 的对象，为 Lagrange 子流形 $L_i$（$i\in\mathbb Z$），按\BookSectionRef{sec:8.3.3}{8.3.3}中的讨论赋予定向与分次，并带有非平凡的自旋结构。态射与高阶乘法映射仍按\BookSectionRef{sec:8.3.3}{8.3.3}中的通常方式给出；横截序列同样是满足任意 $i\ne j$ 都有 $n_i\ne n_j$ 的 $(L_{n_1},\ldots,L_{n_d})$。

\SourceParagraph{ch08-8.4.5-003}{643}
本节的目标是证明：
\BookStatementNumber{24}
\begin{theorem}\label{thm:8.24}
$D(\widehat X)\otimes_{\mathbb C[[q]]}\Lambda_{\mathrm{nov}}$ 与 $\check D(X(B))$ 拟等价。这里的张量积，只是指将所有态射空间与 $\Lambda_{\mathrm{nov}}$ 作张量积。
\end{theorem}

\SourceParagraph{ch08-8.4.5-004}{645}
先说明所选的 \v{C}ech 开覆盖，从而描述 $D(\widehat X)$ 的态射空间。

\SourceParagraph{ch08-8.4.5-005}{645}
已经看到，$X_\Sigma$ 有自然的开覆盖 $\{U_w\mid w\text{ 为 }\Delta\text{ 的顶点}\}$。每个 $U_w$ 限制为 $f^{-1}(0)$ 的一个开子集，从而定义 $\widehat X_\Sigma$ 的开集 $\widehat U_w$。如前所述，这是一个赋环空间，其底拓扑空间与
\[
\operatorname{Spec}\mathbb C[\check w^\vee\cap\mathbb Z^2]/(q)
\]
的底拓扑空间相同，而结构层则由环
\[
\varprojlim_n\mathbb C[\check w^\vee\cap\mathbb Z^2]/(q^n)
=\mathbb C[\check w^\vee\cap\mathbb Z^2]
\widehat\otimes_{\mathbb C[q]}\mathbb C[[q]]
\]
所对应的层给出。开集 $\widehat U_w$ 下降为商 $\widehat X=\widehat X_\Sigma/\mathbb Z$ 的开集，因此得到 $\widehat X$ 的一个开覆盖，其指标可以取为 $\Delta$ 的顶点模 $\mathbb Z$ 作用的一组代表元，等价地，也可以用 $B(\mathbb Z)$ 的元素标记。此外，若 $w_1,w_2$ 为 $\Delta$ 的一条边 $\sigma$ 的两个顶点，则
\[
\widehat U_{w_1}\cap\widehat U_{w_2}=\widehat U_\sigma;
\]
若 $w_1,w_2$ 对应于 $B(\mathbb Z)$ 中不相邻的点，则 $\widehat U_{w_1}\cap\widehat U_{w_2}=\varnothing$。见\BookRef{fig:8.36}{图36}。
\begin{figure}[H]
\centering
\BookImage{01_Numbered_Figures/Chapter_08/Ch08_Fig36_PDFp645_Bookp635.pdf}
\BookNumberOnlyCaption\label{fig:8.36}
\end{figure}

\SourceParagraph{ch08-8.4.5-006}{645}
现在必须作一些选择。需要为这些开集选择一组具体代表元，并用有序指标集标记。取具有通常次序的指标集 $I=\{0,\ldots,d-1\}$，并记
\[
\widehat U_i:=\widehat U_{w_i},
\qquad w_i:=(i,\varphi(i)).
\]
于是，对于 $i<j$，除非 $j=i+1$，或 $i=0$ 且 $j=d-1$，否则 $\widehat U_{i,j}:=\widehat U_i\cap\widehat U_j=\varnothing$。

\SourceParagraph{ch08-8.4.5-007}{646}
如下描述 \v{C}ech 复形。回忆，在典范识别下有
\[
\widehat{\mathcal L}^{\otimes n}|_{\widehat U_w}
=\mathcal O_{\widehat U_w}\cdot z^{nw},
\]
这对于正的或负的 $n$ 都成立。因此
\[
\Gamma(\widehat U_w,\widehat{\mathcal L}^{\otimes n})
=\Gamma_w^n
:=\bigl(\mathbb C[\check w^\vee\cap\mathbb Z^2]
\widehat\otimes_{\mathbb C[q]}\mathbb C[[q]]\bigr)\cdot z^{nw}.
\]
特别地，有识别
\setcounter{equation}{19}
\begin{equation}\label{eq:8.20}
\Gamma(\widehat U_i,\widehat{\mathcal L}^{\otimes n})=\Gamma_{w_i}^n.
\end{equation}
若 $\sigma$ 为 $\Delta$ 中以 $w_i,w_j$ 为端点的一条边，则
\[
\begin{aligned}
\Gamma(\widehat U_\sigma,\widehat{\mathcal L}^{\otimes n})
&=\Gamma(\widehat U_{w_i}\cap\widehat U_{w_j},\widehat{\mathcal L}^{\otimes n})\\
&=\Gamma_\sigma^n
:=\bigl(\mathbb C[\check\sigma^\vee\cap\mathbb Z^2]
\widehat\otimes_{\mathbb C[q]}\mathbb C[[q]]\bigr)\cdot z^{nw_i}\\
&=\bigl(\mathbb C[\check\sigma^\vee\cap\mathbb Z^2]
\widehat\otimes_{\mathbb C[q]}\mathbb C[[q]]\bigr)\cdot z^{nw_j}.
\end{aligned}
\]
所以，若 $\sigma_i$ 为端点是 $w_i,w_{i+1}$ 的区间，则对 $0\leq i<d-1$，典范地有
\begin{equation}\label{eq:8.21}
\Gamma(\widehat U_{i,i+1},\widehat{\mathcal L}^{\otimes n})=\Gamma_{\sigma_i}^n,
\end{equation}
而
\begin{equation}\label{eq:8.22}
\Gamma(\widehat U_{0,d-1},\widehat{\mathcal L}^{\otimes n})=\Gamma_{\sigma_{d-1}}^n.
\end{equation}

\SourceParagraph{ch08-8.4.5-008}{646}
这造成了一种不对称，以下将反复遇到。具体地，对于 $0\leq i<d-1$，限制映射
\[
\Gamma(\widehat U_i,\widehat{\mathcal L}^{\otimes n})
\longrightarrow\Gamma(\widehat U_{i,i+1},\widehat{\mathcal L}^{\otimes n}),
\]
\[
\Gamma(\widehat U_{i+1},\widehat{\mathcal L}^{\otimes n})
\longrightarrow\Gamma(\widehat U_{i,i+1},\widehat{\mathcal L}^{\otimes n})
\]
就是显然的包含 $\Gamma_{w_i}^n,\Gamma_{w_{i+1}}^n\subseteq\Gamma_{\sigma_i}^n$；限制映射
\[
\Gamma(\widehat U_{d-1},\widehat{\mathcal L}^{\otimes n})
\longrightarrow\Gamma(\widehat U_{0,d-1},\widehat{\mathcal L}^{\otimes n})
\]
也是如此。另一方面，限制映射
\[
\Gamma(\widehat U_0,\widehat{\mathcal L}^{\otimes n})
\longrightarrow\Gamma(\widehat U_{0,d-1},\widehat{\mathcal L}^{\otimes n})
\]
由复合
\[
\Gamma_{w_0}^n\xrightarrow{T_n}\Gamma_{w_d}^n
\subseteq\Gamma_{\sigma_{d-1}}^n
\]
给出。

\SourceParagraph{ch08-8.4.5-009}{646}
于是，\v{C}ech 复形为
\[
\begin{aligned}
C^0(\widehat{\mathcal L}^{\otimes n})
&=\bigoplus_{i\in I}\Gamma(\widehat U_i,\widehat{\mathcal L}^{\otimes n})\\
&\longrightarrow C^1(\widehat{\mathcal L}^{\otimes n})\\
&=\left(\bigoplus_{0\leq i<d-1}
\Gamma(\widehat U_{i,i+1},\widehat{\mathcal L}^{\otimes n})\right)
\oplus\Gamma(\widehat U_{0,d-1},\widehat{\mathcal L}^{\otimes n}).
\end{aligned}
\]
将 \v{C}ech 复形的元素写作 $\sum_i(\tau_i,f_i)$，其中
\[
\tau_i\in\{w_j,\sigma_j\mid0\leq j\leq d-1\},
\qquad
f_i\in\Gamma(\widehat U_{\tau_i},\widehat{\mathcal L}^{\otimes n}).
\]
\SourceParagraph{ch08-8.4.5-010}{647}
具体地，有
\begin{equation}\label{eq:8.23}
\begin{aligned}
m_1^{D(\widehat X)}\left(\sum_{i=0}^{d-1}(w_i,z^{m_i})\right)
&=\sum_{i=0}^{d-2}(\sigma_i,z^{m_{i+1}}-z^{m_i})\\
&\quad +(\sigma_{d-1},z^{m_{d-1}}-z^{T_n(m_0)}).
\end{aligned}
\end{equation}
另一方面，可以如下列出 $m_2^{D(\widehat X)}$ 的全部情形；其中第二与第四行中 $0\leq j<d-1$：
\begin{equation}\label{eq:8.24}
\begin{aligned}
m_2\bigl((w_j,z^{m_j}),(w_i,z^{m_i})\bigr)
&=\begin{cases}(w_i,z^{m_i+m_j}),&i=j,\\0,&\text{其他情形};\end{cases}\\[2pt]
m_2\bigl((w_i,z^{m_i}),(\sigma_j,z^{m_j})\bigr)
&=\begin{cases}(\sigma_i,-z^{m_i+m_j}),&i=j,\\0,&\text{其他情形};\end{cases}\\[2pt]
m_2\bigl((w_i,z^{m_i}),(\sigma_{d-1},z^{m_j})\bigr)
&=\begin{cases}(\sigma_{d-1},-z^{T_n(m_i)+m_j}),&i=0,\\0,&\text{其他情形};\end{cases}\\[2pt]
m_2\bigl((\sigma_j,z^{m_j}),(w_i,z^{m_i})\bigr)
&=\begin{cases}(\sigma_j,z^{m_i+m_j}),&i=j+1,\\0,&\text{其他情形};\end{cases}\\[2pt]
m_2\bigl((\sigma_{d-1},z^{m_j}),(w_i,z^{m_i})\bigr)
&=\begin{cases}(\sigma_{d-1},z^{m_i+m_j}),&i=d-1,\\0,&\text{其他情形}.
\end{cases}
\end{aligned}
\end{equation}

\SourceParagraph{ch08-8.4.5-011}{647}
描述了 $A_\infty$ 预范畴 $D(\widehat X)$ 中态射空间的结构后，接着注意，$\check D(X(B))$ 可以与\emph{热带 Morse 范畴}识别。

\SourceParagraph{ch08-8.4.5-012}{647}
\BookStatementNumber{25}
\begin{definition}\label{def:8.25}
设 $B$ 为一维整仿射流形。\emph{热带 Morse 范畴} $\operatorname{TMC}(B)$ 的对象集就是整数集；横截序列为满足任意 $i\ne j$ 都有 $n_i\ne n_j$ 的 $(n_1,\ldots,n_d)$。当 $n_2>n_1$ 时，$\operatorname{Hom}_{\operatorname{TMC}(B)}(n_1,n_2)$ 为
\[
\bigoplus_{p\in B(\frac1{n_2-n_1}\mathbb Z)}[p]\cdot\Lambda_{\mathrm{nov}}
\longrightarrow0;
\]
当 $n_2<n_1$ 时，它为
\[
0\longrightarrow
\bigoplus_{p\in B(\frac1{n_2-n_1}\mathbb Z)}[p]\cdot\Lambda_{\mathrm{nov}}.
\]
高阶乘法映射由计数热带 Morse 树给出：给定横截序列 $(n_0,\ldots,n_d)$ 及
\[
p_{i-1,i}\in B\left(\frac1{n_i-n_{i-1}}\mathbb Z\right),
\]
则 $m_d(p_{d-1,d},\ldots,p_{0,1})$ 中
\[
p_{0,d}\in B\left(\frac1{n_d-n_0}\mathbb Z\right)
\]
对应的项为
\[
\sum_\phi(-1)^{s(\phi)}q^{\deg(\phi)}[p_{0,d}],
\]
\SourceParagraph{ch08-8.4.5-013}{648}
其中 $\phi$ 遍历 $S_d^{\mathrm{trop}}(p_{0,d};p_{0,1},\ldots,p_{d-1,d})$ 中需要计数的孤立元素。

\SourceParagraph{ch08-8.4.5-014}{648}
对于符号 $s(\phi)$，就这里的目的而言，只需赋予热带 Morse 树所定义的全纯圆盘相应的符号即可。这在二维情形中可行；若在更高维工作，显然需要更精细的处理。

\SourceParagraph{ch08-8.4.5-015}{648}
注意，定义中特意要求 $\dim B=1$。确实，可以在任意维数下写出这一定义，但在高维，它是否定义了 $A_\infty$（预）范畴，并不清楚。问题在于，热带 Morse 树的模空间未必具有正确维数，需要某种虚计数或扰动技术，才能把上述描述转化为一个 $A_\infty$ 范畴。问题本质上来自：我们使用了非常特殊的函数的梯度流线，而非一般选择的函数。然而，在一维中，不必担心这些问题；上一小节已经看到，此时热带 Morse 树等价于 $X(B)$ 中的全纯圆盘，而在实二维中计数圆盘，即使不假设将 Lagrange 子流形作一般位置的扰动，也确实给出 $A_\infty$ 结构，参见例如\BookCite{ref003}。

\SourceParagraph{ch08-8.4.5-016}{648}
为了定义次数，将 $\phi\colon S\to B$ 提升到万有覆盖，得到 $\widetilde\phi\colon S\to\widetilde B=\mathbb R$。令
\[
\deg(\phi)=\sum_{i=1}^d(n_i-n_{i-1})
\varphi(\widetilde\phi(v_{i-1,i}))
-(n_d-n_0)\varphi(\widetilde\phi(v_{0,d})).
\]
\end{definition}

\SourceParagraph{ch08-8.4.5-017}{648}
\BookStatementNumber{26}
\begin{proposition}\label{prop:8.26}
$\check D(X(B))$ 与 $\operatorname{TMC}(B)$ 拟同构。
\end{proposition}

\SourceParagraph{ch08-8.4.5-018}{648}
\begin{proof}
当然，$\check D(X(B))$ 的对象 $L_n$ 对应于 $\operatorname{TMC}(B)$ 的对象 $n$。定义映射
\[
F_1\colon\operatorname{Hom}_{\check D(X(B))}(L_{n_0},L_{n_1})
\longrightarrow\operatorname{Hom}_{\operatorname{TMC}(B)}(n_0,n_1)
\]
为
\begin{equation}\label{eq:8.25}
F_1([x])=[x]q^{(n_1-n_0)A(x)},
\end{equation}
其中 $A(x)$ 由\BookEqRef{eq:8.18}{8.18}定义。所有高阶映射 $F_d$ 都取为零。容易验证
\[
m_d^{\operatorname{TMC}(B)}(F_1(x_d),\ldots,F_1(x_1))
=F_1\bigl(m_d^{\check D(X(B))}(x_d,\ldots,x_1)\bigr).
\]
事实上，这些映射分别通过计数热带 Morse 树与全纯圆盘定义，而我们已经看到，两种计数是一回事。此外，符号的定义也一致，唯一需要比较的是全纯圆盘的面积与热带 Morse 树的次数。由\BookRef{lem:8.27}{引理8.27}，在\BookEqRef{eq:8.25}{8.25}之下，热带 Morse 树的次数转化为相应圆盘的面积。
\end{proof}

\SourceParagraph{ch08-8.4.5-019}{649}
\BookStatementNumber{27}
\begin{lemma}\label{lem:8.27}
设 $\phi\colon S\to B$ 为 $B=\mathbb R/d\mathbb Z$ 上的热带 Morse 树，入顶点为 $v_{i-1,i}$，出顶点为 $v_{0,d}$；取提升 $\widetilde\phi\colon S\to\widetilde B=\mathbb R$，以及\BookSectionRef{sec:8.4.4}{8.4.4}给出的 $R\colon\operatorname{Fat}(S)\to X(B)$。假设 $\phi$ 是此类树的零维模空间中的一点。则
\[
\int_{\operatorname{Fat}(S)}R^*\omega
=\sum_{i=1}^d(n_i-n_{i-1})\psi(\widetilde\phi(v_{i-1,i}))
-(n_d-n_0)\psi(\widetilde\phi(v_{0,d})),
\]
其中 $\psi$ 如\BookEqRef{eq:8.17}{8.17}所定义。
\end{lemma}

\SourceParagraph{ch08-8.4.5-020}{649}
\begin{proof}
$\phi$ 属于零维模空间这一事实，意味着将 $R\colon\operatorname{Fat}(S)\to X(B)$ 提升到万有覆盖，得到 $\widetilde R\colon\operatorname{Fat}(S)\to\mathbb C$ 后，$\operatorname{Fat}(S)$ 的像是凸的。否则，如\BookRef{fig:8.22}{图22}那样，可以使此像在一族单参数圆盘中变化，而边界仍位于相应的 Lagrange 子流形中。因此，特别地，$\widetilde R$ 是单射，可以应用平面多边形面积的标准公式。若各顶点坐标为 $(y_i,x_i)$（$0\leq i\leq d$），则其面积为
\[
\frac12\sum_{i=0}^d(y_ix_{i+1}-y_{i+1}x_i),
\qquad (y_{d+1},x_{d+1})=(y_0,x_0).
\]

\SourceParagraph{ch08-8.4.5-021}{649}
现在，多边形各边都是线段，分别位于 $L_0,\ldots,L_d$ 的提升 $\widetilde L_0,\ldots,\widetilde L_d$ 中。假设 $\widetilde L_i$ 的方程为 $x=-n_iy+a_i$，其中 $a_i\in\mathbb Z$。若 $\widetilde\phi(v_{i-1,i})=p_{i-1,i}$，则 $p_{i-1,i}$ 是交点 $\widetilde L_{i-1}\cap\widetilde L_i$ 的第一坐标，因此
\[
-n_{i-1}p_{i-1,i}+a_{i-1}=-n_ip_{i-1,i}+a_i.
\]
同样，若 $\widetilde\phi(v_{0,d})=p_{0,d}$，则它为 $\widetilde L_0\cap\widetilde L_d$ 的第一坐标，所以
\[
-n_0p_{0,d}+a_0=-n_dp_{0,d}+a_d.
\]
不妨取 $a_0=0$。于是对于 $1\leq i\leq d$，可知
\[
a_i=\sum_{j=1}^i(n_j-n_{j-1})p_{j-1,j},
\qquad
a_d=(n_d-n_0)p_{0,d}.
\]

\SourceParagraph{ch08-8.4.5-022}{650}
令 $p_{d,d+1}:=p_{0,d}$，面积公式便给出
\[
\begin{aligned}
&\frac12\sum_{i=1}^d
\bigl(p_{i-1,i}(-n_ip_{i,i+1}+a_i)
-p_{i,i+1}(-n_ip_{i-1,i}+a_i)\bigr)\\
&\quad=\frac12\sum_{i=1}^d(p_{i-1,i}-p_{i,i+1})a_i\\
&\quad=\frac12\sum_{i=1}^d\sum_{j=1}^i
(n_j-n_{j-1})p_{j-1,j}(p_{i-1,i}-p_{i,i+1})\\
&\quad=\frac12\sum_{j=1}^d
(n_j-n_{j-1})p_{j-1,j}(p_{j-1,j}-p_{0,d})\\
&\quad=\sum_{j=1}^d\left(
(n_j-n_{j-1})\frac{p_{j-1,j}(p_{j-1,j}-1)}2
-(n_j-n_{j-1})\frac{p_{j-1,j}(p_{0,d}-1)}2\right)\\
&\quad=\sum_{j=1}^d(n_j-n_{j-1})\psi(p_{j-1,j})
-(n_d-n_0)\psi(p_{0,d}).
\end{aligned}
\]
最后一个等式，使用了 $a_d$ 的两种表达式相等这一事实。
\end{proof}

\SourceParagraph{ch08-8.4.5-023}{650}
很自然地会想，直接比较热带 Morse 范畴给出的 $A_\infty$ 结构，与\BookSectionRef{sec:8.2.1}{8.2.1}的构造在 $D(\widehat X)$ 上给出的 $A_\infty$ 结构。这需要写出 \v{C}ech 复形的投影与同伦。然而，这样得到的线丛一侧的描述，虽然在某种程度上类似于热带 Morse 一侧的图景，却不能完全吻合。问题本质上在于，\v{C}ech 覆盖的离散性，使我们无法细致看清热带 Morse 树在 $B$ 的线段 $[i,i+1]$ 内发生的变化。因此，最方便的办法，是更直接地写出一个给出拟同构的 $A_\infty$ 函子。这仍是对\BookCite{ref002}中思想的改编。

\SourceParagraph{ch08-8.4.5-024}{650}
首先，定义一点额外结构。设 $\mathcal P$ 为将 $B$ 分解成长度为一的线段的多面体分解；即 $\mathcal P$ 由单位区间 $[i,i+1]$ 及顶点 $\{i\}$ 组成。沿用一个略有滥用的记法，将区间 $[d-1,d]$ 在 $B$ 中的像记为 $[0,d-1]$；写成 $[d-1,0]$ 可能更合理，但希望把较小的指标放在左边。对每个 $0\leq i<d$，选取一般位置的点 $\alpha_i\in(i,i+1)$。于是可以定义对偶胞腔复形 $\check{\mathcal P}$，由顶点 $\{\alpha_i\}$、区间 $[\alpha_i,\alpha_{i+1}]$（$0\leq i<d-1$），以及 $[\alpha_{d-1},d+\alpha_0]$ 组成；记住，这些区间都按模 $d\mathbb Z$ 来看。$\mathcal P$ 的胞腔与 $\check{\mathcal P}$ 的胞腔之间，显然存在一个反转包含关系的一一对应，记为 $\sigma\mapsto\sigma^\vee$。例如，对于 $0<i\leq d-1$，有 $i^\vee=[\alpha_{i-1},\alpha_i]$。

\SourceParagraph{ch08-8.4.5-025}{651}
对任意充分小的 $\epsilon>0$，将构造 $\check{\mathcal P}$ 的一个微小变形 $\check{\mathcal P}_\epsilon$。其一维胞腔为
\[
\begin{gathered}
\{[\alpha_i-\epsilon,\alpha_{i+1}-\epsilon]\mid0\leq i<d-2\}\\
\cup\{[\alpha_{d-2}-\epsilon,\alpha_{d-1}+\epsilon],
[\alpha_{d-1}+\epsilon,\alpha_0+d-\epsilon]\},
\end{gathered}
\]
零维胞腔当然就是这些区间的端点。见\BookRef{fig:8.37}{图37}。给定胞腔 $\sigma\in\mathcal P$，将其在 $\check{\mathcal P}_\epsilon$ 中的对偶胞腔的内部记为 $\sigma_\epsilon^\vee$。
\begin{figure}[H]
\centering
\BookImage{01_Numbered_Figures/Chapter_08/Ch08_Fig37_PDFp651_Bookp641.pdf}
\caption{内圈同时表示 $\mathcal P$ 与 $\check{\mathcal P}$；外圈表示 $\check{\mathcal P}_{\epsilon_1}$ 与 $\check{\mathcal P}_{\epsilon_2}$，其中 $\epsilon_2>\epsilon_1$。}\label{fig:8.37}
\end{figure}

\SourceParagraph{ch08-8.4.5-026}{651}
这些分解具有如下性质。若 $\epsilon_1<\epsilon_2$，则对于 $0\leq j<d-1$，
\[
i_{\epsilon_1}^\vee\cap[j,j+1]_{\epsilon_2}^\vee\ne\varnothing
\quad\Longleftrightarrow\quad i=j,
\]
而
\[
i_{\epsilon_1}^\vee\cap[0,d-1]_{\epsilon_2}^\vee\ne\varnothing
\quad\Longleftrightarrow\quad i=0.
\]
同时，对于 $0\leq j<d-1$，
\[
[j,j+1]_{\epsilon_1}^\vee\cap i_{\epsilon_2}^\vee\ne\varnothing
\quad\Longleftrightarrow\quad i=j+1,
\]
而
\[
[0,d-1]_{\epsilon_1}^\vee\cap i_{\epsilon_2}^\vee\ne\varnothing
\quad\Longleftrightarrow\quad i=d-1.
\]
将这些关系与 $m_2^{D(\widehat X)}$ 的公式\BookEqRef{eq:8.24}{8.24}比较，就会看到，这一组看起来颇为古怪的多面体分解，是为了在进行乘法时强制遵守 \v{C}ech 开覆盖指标集的次序。

\SourceParagraph{ch08-8.4.5-027}{652}
接着，在 $B=\mathbb R/d\mathbb Z$ 上定义热带 Morse 树的一个变体。
\BookStatementNumber{28}
\begin{definition}\label{def:8.28}
$T_d$ 是\emph{灌木树}的模空间。灌木树是带状树，其所有边都被赋予有限长度，唯有出边的长度仍为无穷；并且满足：若 $v_{\mathrm{adj}}$ 为与出顶点相邻的唯一顶点，则从 $v_{\mathrm{adj}}$ 到任意入顶点的距离，都与所选入顶点无关。称此条件为\emph{长度条件}。

\SourceParagraph{ch08-8.4.5-028}{652}
$T_1$ 只有一个点，对应于一条长度为无穷的线段，也就是出边。
\end{definition}

\SourceParagraph{ch08-8.4.5-029}{652}
以下设 $\widetilde B=\mathbb R\to B$ 为 $B$ 的万有覆盖，$\pi_1\colon\mathbb Z^2\to\mathbb Z$ 为到第一坐标的投影。

\SourceParagraph{ch08-8.4.5-030}{652}
\BookStatementNumber{29}
\begin{definition}\label{def:8.29}
假设给定数据
\[
\begin{gathered}
r\geq1,\qquad n_0,\ldots,n_r\in\mathbb Z,\\
\tau_{i-1,i}\subsetneq B\quad(1\leq i\leq r),\\
\widetilde\tau_{i-1,i}\subset\widetilde B
\text{ 为 }\tau_{i-1,i}\text{ 的一个提升},\\
m_{i-1,i}\in\mathbb Z^2\quad(1\leq i\leq r),\\
p_{0,r}\in B\left(\frac1{n_r-n_0}\mathbb Z\right).
\end{gathered}
\]
定义
\[
T_r^{\mathrm{trop}}\bigl(p_{0,r};(\tau_{0,1},m_{0,1}),\ldots,
(\tau_{r-1,r},m_{r-1,r})\bigr)
\]
为如下数据的模空间：从灌木树 $S\in T_r$ 出发的连续映射 $\phi\colon S\to B$，以及 $S$ 各边 $e$ 上的一组仿射位移向量 $v_e\in\Gamma(e,(\phi|_e)^*TB)$，满足以下性质。
\begin{enumerate}
\item 若 $v_i$ 为第 $i$ 个入顶点，则 $\phi(v_i)\in\tau_{i-1,i}$。若 $v_{\mathrm{out}}$ 为外部出顶点，则 $\phi(v_{\mathrm{out}})=p_{0,r}$。
\item 满足\BookRef{def:8.19}{定义8.19}中的公理(2)。
\item 若 $e_i$ 为 $S$ 中唯一包含 $v_i$ 的边，设 $\widetilde v_i\in\widetilde B$ 为 $\phi(v_i)$ 在 $\widetilde\tau_{i-1,i}$ 中的提升，则
\[
v_{e_i}(v_i)=(n_i-n_{i-1})\widetilde v_i-\pi_1(m_{i-1,i}).
\]
若 $e_{\mathrm{out}}$ 为 $S$ 中唯一包含 $v_{\mathrm{out}}$ 的边，则 $v_{e_{\mathrm{out}}}(v_{\mathrm{out}})=0$。
\item 满足\BookRef{def:8.19}{定义8.19}中的公理(4)。
\item 满足\BookRef{def:8.19}{定义8.19}中的公理(5)。
\item 满足\BookRef{def:8.19}{定义8.19}中的公理(6)。
\end{enumerate}
\end{definition}

\SourceParagraph{ch08-8.4.5-031}{653}
\BookStatementNumber{30}
\begin{note}[numbered]\label{note:8.30}
\begin{enumerate}
\item 这些树与先前定义的热带 Morse 树的主要区别，在于入顶点不再是对应于 Lagrange 子流形交点的指定点，而是被要求位于某些胞腔中；初始仿射位移向量也不再为零，而是由提升 $\widetilde\tau_{i-1,i}$ 与 $m_{i-1,i}$ 这些附加数据指定。注意，从第 $i$ 个入顶点出发，若 $n_i-n_{i-1}>0$，就远离 $\pi_1(m_{i-1,i})/(n_i-n_{i-1})$；若 $n_i-n_{i-1}<0$，则朝向这个点。
\SourceParagraph{ch08-8.4.5-032}{653}
\item 尽管记号中没有写出提升 $\widetilde\tau_{i-1,i}$，模空间确实依赖于这些提升。若将 $\widetilde\tau_{i-1,i}$ 换成 $\widetilde\tau_{i-1,i}+d$，则须同时将 $m_{i-1,i}$ 换成 $T_{n_i-n_{i-1}}(m_{i-1,i})$。请验证，这一变换给出相同的模空间。
\SourceParagraph{ch08-8.4.5-033}{653}
\item 按(2)所述改变提升 $\widetilde\tau_{i-1,i}$ 及 $m_{i-1,i}$ 后，对于给定的 $\phi\colon S\to B$，可以找到提升 $\widetilde\phi\colon S\to\widetilde B$，使第 $i$ 个入顶点的像总是位于 $\widetilde\tau_{i-1,i}$ 中。于是，若 $S$ 的边 $e$ 标记为 $n_j-n_{i-1}$，其中 $j>i$，则对于任意 $p\in e$，实际上有
\[
v_e(p)=(n_j-n_{i-1})\widetilde\phi(p)
-\pi_1\left(\sum_{s=i}^jm_{s-1,s}\right).
\]
根据定义，这在起始处成立，随后可归纳证明。由于仿射位移向量在出顶点 $v_{\mathrm{out}}$ 处必须为零，实际上得到
\[
\widetilde\phi(v_{\mathrm{out}})
=\frac{\pi_1(\sum_{i=1}^rm_{i-1,i})}{n_r-n_0},
\]
它模 $d\mathbb Z$ 后就是 $p_{0,r}$。
\end{enumerate}
\end{note}

\SourceParagraph{ch08-8.4.5-034}{653}
\BookStatementNumber{31}
\begin{definition}\label{def:8.31}
给定
\[
\phi\in T_r^{\mathrm{trop}}\bigl(p_{0,r};(\tau_{0,1},m_{0,1}),\ldots,
(\tau_{r-1,r},m_{r-1,r})\bigr),
\]
取\BookRef{note:8.30}{笔记8.30}(3)中的提升 $\widetilde\phi\colon S\to\widetilde B$；这可能需要改变各个 $\widetilde\tau_{i-1,i}$ 与 $m_{i-1,i}$。于是可以将 $\phi$ 的次数定义为
\[
\deg\phi=\operatorname{ord}_{n_r-n_0}
\left(\sum_{i=1}^rm_{i-1,i}\right).
\]
\end{definition}

\SourceParagraph{ch08-8.4.5-035}{653}
\BookStatementNumber{32}
\begin{proposition}\label{prop:8.32}
设 $\sigma_{0,1},\ldots,\sigma_{d-1,d}\in\mathcal P$，$\epsilon_1,\ldots,\epsilon_d$ 充分小，且
\[
\phi\in T_d^{\mathrm{trop}}\bigl(p_{0,d};
((\sigma_{0,1})_{\epsilon_1}^\vee,m_{0,1}),\ldots,
((\sigma_{d-1,d})_{\epsilon_d}^\vee,m_{d-1,d})\bigr),
\]
并且对每个 $i$，都有
\[
z^{m_{i-1,i}}\in
\Gamma(\widehat U_{\sigma_{i-1,i}},
\widehat{\mathcal L}^{\otimes(n_i-n_{i-1})}).
\]
则 $\deg\phi\geq0$。
\end{proposition}

\SourceParagraph{ch08-8.4.5-036}{653}
\begin{proof}
取\BookRef{note:8.30}{笔记8.30}(3)中的提升 $\widetilde\phi\colon S\to\widetilde B$。由关于 $z^{m_{i-1,i}}$ 的假设，若第 $i$ 个入顶点在 $\widetilde B$ 中的像为 $\widetilde v_i$，则 $m_{i-1,i}$ 位于某条直线上方：该直线由 $(n_i-n_{i-1})\Delta$ 的一条边张成，且这条边包含点
\[
\bigl((n_i-n_{i-1})\widetilde v_i,
(n_i-n_{i-1})\varphi(\widetilde v_i)\bigr).
\]
视 $\widetilde v_i$ 而定，这样的直线有一条或两条；后一情形中，$m_{i-1,i}$ 位于一个楔形区域内。将归纳证明：对于 $S$ 中标记为 $n_e=n_j-n_{i-1}$ 的边上的任意点 $p$，$\sum_{s=i}^jm_{s-1,s}$ 位于某条直线上方，该直线由 $(n_j-n_{i-1})\Delta$ 的一条边张成，且这条边包含点
\[
\bigl((n_j-n_{i-1})\widetilde\phi(p),
(n_j-n_{i-1})\varphi(\widetilde\phi(p))\bigr).
\]

\SourceParagraph{ch08-8.4.5-037}{654}
先假设这在 $S$ 的某条边 $e$ 的起始顶点 $v_{\mathrm{in}}$ 处成立。由\BookRef{note:8.30}{笔记8.30}(3)，仿射位移向量为
\[
v_e(v_{\mathrm{in}})=(n_j-n_{i-1})\widetilde\phi(v_{\mathrm{in}})
-\pi_1\left(\sum_{s=i}^jm_{s-1,s}\right).
\]
令
\[
x:=(n_j-n_{i-1})^{-1}\pi_1\left(\sum_{s=i}^jm_{s-1,s}\right).
\]
若 $n_j-n_{i-1}>0$，该向量指向远离 $x$ 的方向；若 $n_j-n_{i-1}<0$，则指向 $x$。前一情形中，$\widetilde\phi(e)$ 必须远离 $x$，所以由 $v_{\mathrm{in}}$ 处的归纳假设及\BookRef{fig:8.38}{图38}，显然归纳假设继续在任意 $p\in e$ 处成立。若 $n_j-n_{i-1}<0$，$\widetilde\phi(e)$ 朝向 $x$，但永不越过它。由\BookRef{fig:8.39}{图39}，同样可见归纳假设继续在任意 $p\in e$ 处成立。

\SourceParagraph{ch08-8.4.5-038}{654}
最后，为说明归纳假设在 $v_{\mathrm{in}}$ 处成立，注意 $\sum_{s=i}^jm_{s-1,s}$，正是所有进入 $v_{\mathrm{in}}$ 的边上相应的和的总和。既然归纳假设在每条入边的 $v_{\mathrm{in}}$ 处成立，它也就在边 $e$ 上成立。

\SourceParagraph{ch08-8.4.5-039}{654}
再考察 $S$ 的最后一个顶点经 $\widetilde\phi$ 的像 $\widetilde v_{\mathrm{out}}$。由于
\[
(n_d-n_0)\widetilde v_{\mathrm{out}}
=\pi_1\left(\sum_{i=1}^dm_{i-1,i}\right),
\]
归纳假设意味着 $\sum_{i=1}^dm_{i-1,i}$ 恰好位于
\[
\bigl((n_d-n_0)\widetilde v_{\mathrm{out}},
(n_d-n_0)\varphi(\widetilde v_{\mathrm{out}})\bigr)
\]
的正上方。而 $\operatorname{ord}_{n_d-n_0}(\sum_{i=1}^dm_{i-1,i})$，恰好是这两点之间的竖直距离，因此非负。
\end{proof}

\SourceParagraph{ch08-8.4.5-040}{654}
可以启发式地计算
\[
\dim T_r^{\mathrm{trop}}\bigl(p_{0,r};(\tau_{0,1},m_{0,1}),\ldots,
(\tau_{r-1,r},m_{r-1,r})\bigr).
\]
取该空间中内部顶点均为三价的一个元素，考察允许怎样使它变化。不难看到，在施加长度条件之前，除 $v_{\mathrm{adj}}$ 以外，每个内部顶点的像都可以自由移动；而 $v_{\mathrm{adj}}$ 的像可以自由移动，当且仅当 $\deg p_{0,r}=1$。此外，第 $i$ 个入顶点的像可以在 $\tau_{i-1,i}$ 内变化。移动这些顶点，会改变各边的长度，并希望这些长度的变化彼此独立；随后施加长度条件，就增加了 $r-1$ 个约束。因此总维数为
\[
\begin{aligned}
&\text{内部顶点数}-1+\deg p_{0,r}
+\sum_{i=1}^r\dim\tau_{i-1,i}-(r-1)\\
&\qquad=\deg p_{0,r}-1+\sum_{i=1}^r\dim\tau_{i-1,i}.
\end{aligned}
\]
\begin{figure}[H]
\centering
\BookImage{01_Numbered_Figures/Chapter_08/Ch08_Fig38_PDFp655_Bookp645.pdf}
\caption{$n>0$。虚线表示 $\Gamma_\sigma^n$ 与 $\Gamma_w^n$。注意，左图中若 $m\in\Gamma_\sigma^n$ 位于 $nv$ 的左侧，则对于 $v$ 右侧的 $w$，还有 $m\in\Gamma_w^n$。右图中若 $m\in\Gamma_\sigma^n$ 位于 $nv$ 的右侧，则对于 $v$ 左侧的 $w$，还有 $m\in\Gamma_w^n$。}\label{fig:8.38}
\end{figure}
\begin{figure}[H]
\centering
\BookImage{01_Numbered_Figures/Chapter_08/Ch08_Fig39_PDFp655_Bookp645.pdf}
\caption{$n<0$。}\label{fig:8.39}
\end{figure}

\SourceParagraph{ch08-8.4.5-041}{655}
事实上，若对每个 $i$ 都有 $\dim\tau_{i-1,i}=1$，则这个公式总是正确的：一旦第一个顶点及所有内部顶点的位置固定，长度条件就确定了其余所有入顶点的位置。

\SourceParagraph{ch08-8.4.5-042}{655}
以下关心的是模空间
\[
T_d^{\mathrm{trop}}
:=T_d^{\mathrm{trop}}\bigl(p_{0,d};
((\sigma_{0,1})_{\epsilon_1}^\vee,m_{0,1}),\ldots,
((\sigma_{d-1,d})_{\epsilon_d}^\vee,m_{d-1,d})\bigr),
\]
其中给定任意一列 $\sigma_{i-1,i}\in\mathcal P$、$m_{i-1,i}\in\mathbb Z^2$（$i=1,\ldots,d$），并且 $0<\epsilon_1<\cdots<\epsilon_d$。始终约定，集合 $(\sigma_{i-1,i})_{\epsilon_i}^\vee$ 的提升，是以下集合之一：
\SourceParagraph{ch08-8.4.5-043}{656}
\[
\begin{aligned}
&(\alpha_j-\epsilon_i,\alpha_{j+1}-\epsilon_i),
&&0\leq j<d-2,\\
&(\alpha_{d-2}-\epsilon_i,\alpha_{d-1}+\epsilon_i),\\
&(\alpha_{d-1}+\epsilon_i,\alpha_0+d-\epsilon_i),\\
&\{\alpha_j-\epsilon_i\},&&0\leq j\leq d-2,\\
&\{\alpha_{d-1}+\epsilon_i\}.
\end{aligned}
\]

\SourceParagraph{ch08-8.4.5-044}{656}
总可以找到 $0<\epsilon_1<\cdots<\epsilon_d$，使 $T_d^{\mathrm{trop}}$ 具有预期维数。事实上，令 $I=\{i\mid\dim\sigma_{i-1,i}^\vee=0\}$，并令 $T_d'$ 为上述模空间的变体：对每个 $i\in I$，将 $(\sigma_{i-1,i})_{\epsilon_i}^\vee$ 换成其附近的一个区间。有取值映射
\[
\operatorname{ev}\colon T_d'\longrightarrow B^{\#I},
\]
在 $I$ 所标记的入顶点处对 $\phi\in T_d'$ 取值。于是
\[
T_d^{\mathrm{trop}}=
\operatorname{ev}^{-1}\left(\prod_{i\in I}(\sigma_{i-1,i})_{\epsilon_i}^\vee\right).
\]
根据前面的论证，$T_d'$ 具有预期维数。如果 $\operatorname{ev}$ 的像的维数不等于 $\#I$，则对于一般选择的 $\epsilon_i$，$T_d^{\mathrm{trop}}$ 为空；否则，对于充分一般的 $\epsilon_i$，纤维的维数就是预期维数。

\SourceParagraph{ch08-8.4.5-045}{656}
实际上，不难看到，当这个模空间为零维时，还可以得到稍强的结论：随着 $\epsilon_i$ 变化，$T_d^{\mathrm{trop}}$ 中的树连续变化，只是在某些一维模空间的边界点出现时，可能发生跳跃。特别地，对任意给定的上述数据 $(\sigma_{i-1,i},m_{i-1,i})$，可以找到某个 $\epsilon$，使任意 $0<\epsilon_1<\cdots<\epsilon_d<\epsilon$ 所对应的模空间 $T_d^{\mathrm{trop}}$ 都自然同构。将用这些模空间定义所需的 $A_\infty$ 函子。实际上，这里有一个微妙之处，在本次叙述中希望暂且避开：这些模空间即使为零维，也可能是无限的，并且无论 $\epsilon$ 多小，都可能出现跳跃。不过，若将 $\phi$ 的次数限制在某个给定上界内，零维模空间就是有限的。这使我们能将下文的映射 $F_d$ 定义到 $q$ 的任意给定幂次，从而在 $\mathbb C[[q]]$ 上定义 $F_d$。以下不再讨论这一细节。关于更自然的解决办法，见\BookCite{ref002}。

\SourceParagraph{ch08-8.4.5-046}{656}
如下定义 $A_\infty$ 函子 $F\colon D(\widehat X)\to\operatorname{TMC}(B)$。在对象层面，$F(\widehat{\mathcal L}^{\otimes n})=n$。对每个 $d\geq1$ 及横截序列 $(\widehat{\mathcal L}^{\otimes n_0},\ldots,\widehat{\mathcal L}^{\otimes n_d})$，定义映射
\[
\begin{gathered}
F_d\colon C^*(\widehat{\mathcal L}^{\otimes(n_d-n_{d-1})})
\otimes\cdots\otimes C^*(\widehat{\mathcal L}^{\otimes(n_1-n_0)})\\
\longrightarrow\operatorname{Hom}_{\operatorname{TMC}(B)}^*(n_0,n_d)
\end{gathered}
\]
如下。

\SourceParagraph{ch08-8.4.5-047}{657}
选取胞腔 $\sigma_{0,1},\ldots,\sigma_{d-1,d}\in\mathcal P$，以及 $m_{i-1,i}\in\mathbb Z^2$，使得在识别\BookEqRef{eq:8.20}{8.20}、\BookEqRef{eq:8.21}{8.21}或\BookEqRef{eq:8.22}{8.22}之下，有
\[
z^{m_{i-1,i}}\in\Gamma(\widehat U_{\sigma_{i-1,i}},
\widehat{\mathcal L}^{\otimes(n_i-n_{i-1})}).
\]
这定义了
\[
C^*(\widehat{\mathcal L}^{\otimes(n_d-n_{d-1})})
\otimes\cdots\otimes C^*(\widehat{\mathcal L}^{\otimes(n_1-n_0)})
\]
中的元素
\[
(\sigma_{d-1,d},z^{m_{d-1,d}})\otimes\cdots\otimes
(\sigma_{0,1},z^{m_{0,1}}).
\]
给定 $p_{0,d}\in B(\frac1{n_d-n_0}\mathbb Z)$，此元素在 $\operatorname{Hom}_{\operatorname{TMC}(B)}^*(n_0,n_d)$ 中的像，其 $[p_{0,d}]$ 的系数为
\[
\begin{cases}
\displaystyle\sum_\phi(-1)^{s'(\phi)}q^{\deg\phi},
&\displaystyle\deg p_{0,d}-1+\sum_i\dim\sigma_{i-1,i}^\vee=0,\\[4pt]
0,&\displaystyle\deg p_{0,d}-1+\sum_i\dim\sigma_{i-1,i}^\vee\ne0,
\end{cases}
\]
其中求和遍历
\[
T_d^{\mathrm{trop}}\bigl(p_{0,d};
((\sigma_{0,1})_{\epsilon_1}^\vee,m_{0,1}),\ldots,
((\sigma_{d-1,d})_{\epsilon_d}^\vee,m_{d-1,d})\bigr)
\]
中的所有 $\phi$。这里 $0<\epsilon_1<\cdots<\epsilon_d<\epsilon$，而 $\epsilon$ 按上一段选取，使该模空间不依赖于各个 $\epsilon_i$ 的具体选择。此处略去符号 $s'(\phi)$ 的定义，因为它相当复杂。\BookCite{ref002}的附录 B 给出了如何处理这一符号的线索。

\SourceParagraph{ch08-8.4.5-048}{657}
为了理解 $F_d$ 的作用，先考察 $d=1$ 的情形。给定胞腔 $\sigma\in\mathcal P$ 及 $z^m\in\Gamma(\widehat U_\sigma,\widehat{\mathcal L}^{\otimes(n_1-n_0)})$，分为 $n_1>n_0$ 与 $n_1<n_0$ 两种情形。考虑可能出现的树时，用提升 $\widetilde\phi\colon S\to\widetilde B$ 来理解最方便。此时，这些树只是从 $\sigma_\epsilon^\vee$ 出发，远离或朝向 $\pi_1(m)/(n_1-n_0)$ 的线段。不过，若 $n_1>n_0$，仿射位移向量的长度递增，因而永远不能达到零，除非起始顶点 $v$ 满足 $\widetilde\phi(v)=\pi_1(m)/(n_1-n_0)$。此时得到一棵退化树，唯一的边映到点 $\pi_1(m)/(n_1-n_0)$。这只可能在
\[
p_{0,1}=\pi_1(m)/(n_1-n_0)\pmod{d\mathbb Z},
\qquad p_{0,1}\in\sigma_\epsilon^\vee
\]
时出现，而且是唯一的贡献。它给出
\[
q^{\operatorname{ord}_{n_1-n_0}(m)}[p_{0,1}].
\]
$q$ 的指数，表示 $z^m$ 在 $\Gamma(\widehat U_\sigma,\widehat{\mathcal L}^{\otimes(n_1-n_0)})$ 中能被 $q$ 整除的最高次数。注意，这只可能在 $\dim\sigma=0$ 时发生；因为若 $\sigma_\epsilon^\vee$ 为一点，所选的 $\epsilon$ 已充分小，使该点不可能是分母为 $n_1-n_0$ 的有理点。

\SourceParagraph{ch08-8.4.5-049}{657}
若 $n_1<n_0$，仿射位移向量递减，并恰好在线段到达 $\pi_1(m)/(n_1-n_0)$ 时为零。因此，总存在这样的树。此类树的模空间维数为 $\dim\sigma_\epsilon^\vee$；若这个维数为一，则 $F_1$ 在该 \v{C}ech 上链上为零。否则 $\dim\sigma=1$，得到与 $n_1>n_0$ 情形相同的贡献公式。不过，规定
\SourceParagraph{ch08-8.4.5-050}{658}
\[
s'(\phi)=\begin{cases}
1,&\sigma=[0,d-1],\\
0,&\sigma\ne[0,d-1].
\end{cases}
\]

\SourceParagraph{ch08-8.4.5-051}{658}
\begin{proof}[定理8.24的证明梗概]
$A_\infty$ 函子的次数一关系为
\[
m_1^{\operatorname{TMC}(B)}\circ F_1
=F_1\circ m_1^{D(\widehat X)}.
\]
当然，$m_1^{\operatorname{TMC}(B)}=0$。所以，只需在 $n_1<n_0$ 的情形，对上述 $(\{i\},z^m)$ 验证 $F_1\circ m_1^{D(\widehat X)}=0$。$m_1^{D(\widehat X)}$ 的结果有两项，并取决于 $i$：
\[
\begin{aligned}
&m_1^{D(\widehat X)}\bigl((\{i\},z^m)\bigr)\\
&\quad=\begin{cases}
([i-1,i],z^m)+([i,i+1],-z^m),&0<i<d-1,\\
([0,d-1],-z^{T_n(m)})+([0,1],-z^m),&i=0,\\
([d-2,d-1],z^m)+([0,d-1],z^m),&i=d-1.
\end{cases}
\end{aligned}
\]
容易看到，对这些结果应用 $F_1$，总得到零。同样容易看出，$F_1$ 诱导链复形的拟同构，因此只要 $F$ 是 $A_\infty$ 函子，它就是一个拟同构。

\SourceParagraph{ch08-8.4.5-052}{658}
下面简述 $F$ 为 $A_\infty$ 函子的论证。记住，$D(\widehat X)$ 中只有 $m_1$ 与 $m_2$ 非零。对每个 $d$，需验证如下关系：
\[
\begin{aligned}
&\sum_r\sum_{\substack{s_1,\ldots,s_r\geq1\\s_1+\cdots+s_r=d}}
m_r^{\operatorname{TMC}(B)}
\bigl(F_{s_r}(\cdots),\ldots,F_{s_1}(\cdots)\bigr)\\
&\quad=\sum_q(-1)^{\deg a_1+\cdots+\deg a_q-q}
F_d\bigl(a_d,\ldots,a_{q+2},m_1^{D(\widehat X)}(a_{q+1}),
a_q,\ldots,a_1\bigr)\\
&\qquad+\sum_q(-1)^{\deg a_1+\cdots+\deg a_q-q}
F_{d-1}\bigl(a_d,\ldots,a_{q+3},m_2^{D(\widehat X)}(a_{q+2},a_{q+1}),
a_q,\ldots,a_1\bigr).
\end{aligned}
\]
固定 $(n_0,\ldots,n_d)$、$a_i=(\sigma_{i-1,i},z^{m_{i-1,i}})$（$i=1,\ldots,d$），以及 $p_{0,d}\in B(\frac1{n_d-n_0}\mathbb Z)$，希望证明上述关系。选取充分小的 $0<\epsilon_1<\cdots<\epsilon_d<\epsilon$；究竟需要多小，将在以下论证中明确。

\SourceParagraph{ch08-8.4.5-053}{658}
由于次数的限制，除非模空间
\[
T_d^{\mathrm{trop}}\bigl(p_{0,d};
((\sigma_{0,1})_{\epsilon_1}^\vee,m_{0,1}),\ldots,
((\sigma_{d-1,d})_{\epsilon_d}^\vee,m_{d-1,d})\bigr)
\]
为一维，否则 $A_\infty$ 函子关系中 $[p_{0,d}]$ 的所有系数都为零。考察这个模空间的边界。自然，带上正确符号后，边界的总计数为零。下面考察该模空间的元素可能怎样退化。

\SourceParagraph{ch08-8.4.5-054}{658}
(1) 某些边的长度可以趋于 $\infty$。这意味着，在出边以外的某条边上，仿射位移向量趋于零。由于从顶点 $v_{\mathrm{adj}}$ 到任意入顶点的长度都相同，只有当每条从 $v_{\mathrm{adj}}$ 到入顶点的路径，都包含一条长度趋于无穷的边时，这种情形才会发生。若 $\epsilon_i$ 选得充分一般，容易验证，每条这样的路径上恰有一条此类边。因此，可以将极限树看作若干棵\BookRef{def:8.29}{定义8.29}中那种树的出顶点，粘接到一棵热带 Morse 树的入顶点上。这样的退化对第一个项作出贡献。

\SourceParagraph{ch08-8.4.5-055}{659}
(2) 某个入顶点移动到其所在胞腔的边界。选择的一般性同样保证，这只会发生在一个胞腔上。这样的退化对第二个项作出贡献。

\SourceParagraph{ch08-8.4.5-056}{659}
(3) 某条入边的长度可以趋于零。为了使从 $v_{\mathrm{adj}}$ 回到任意入顶点的长度都相同，至少两条相邻入边的长度必须同时趋于零。一般选择又保证，这种情形只会成对发生。因此，存在某个 $i$，使第 $i$ 与第 $i+1$ 个入顶点 $v_i,v_{i+1}$ 满足 $\phi(v_i)=\phi(v_{i+1})$；从而可以将 $\phi$ 看作以下模空间中的树：
\[
\begin{aligned}
T_{d-1}^{\mathrm{trop}}\bigl(&p_{0,d};
((\sigma_{0,1})_{\epsilon_1}^\vee,m_{0,1}),\ldots,\\
&((\sigma_{i-1,i})_{\epsilon_i}^\vee
\cap(\sigma_{i,i+1})_{\epsilon_{i+1}}^\vee,
m_{i-1,i}+m_{i,i+1}),\\
&\ldots,((\sigma_{d-1,d})_{\epsilon_d}^\vee,m_{d-1,d})\bigr).
\end{aligned}
\]
事实上，若 $v'$ 为与 $v_i,v_{i+1}$ 相邻的顶点，$e$ 为从 $v'$ 出发的出边，则可选择提升 $\widetilde\phi\colon S\to\widetilde B$，使 $\widetilde\phi(v')=\widetilde\phi(v_i)=\widetilde\phi(v_{i+1})$，并且
\[
v_e(v')=(n_{i+1}-n_{i-1})\widetilde\phi(v')
-\pi_1(m_{i-1,i}+m_{i,i+1}).
\]
这就解释了为什么出现 $m_{i-1,i}+m_{i,i+1}$。若交集为空，这种退化当然不可能发生。若交集非空，而 $(\sigma_{i-1,i})_{\epsilon_i}^\vee$ 或 $(\sigma_{i,i+1})_{\epsilon_{i+1}}^\vee$ 中有一个是零维的，则交集就是这个零维胞腔，此时的贡献恰好对应于\BookEqRef{eq:8.24}{8.24}后四种情形所描述的 $m_2$。注意，条件 $\epsilon_i<\epsilon_{i+1}$ 强制遵守 \v{C}ech 开覆盖指标集的次序。

\SourceParagraph{ch08-8.4.5-057}{659}
另一方面，若交集为一维，有两种可能。首先，可以有 $\sigma_{i-1,i}\ne\sigma_{i,i+1}$，此时交集非常小，大小约为 $\epsilon$。可以假设 $\epsilon$ 充分小，使这个模空间实际上为空，因而没有贡献。若 $\sigma_{i-1,i}=\sigma_{i,i+1}$，则
\[
(\sigma_{i-1,i})_{\epsilon_i}^\vee
\cap(\sigma_{i,i+1})_{\epsilon_{i+1}}^\vee
\]
非常接近 $(\sigma_{i-1,i})_{\epsilon_i}^\vee$，因此可以用后者代替该交集，而不影响模空间。注意，此时 $\dim\sigma_{i-1,i}=\dim\sigma_{i,i+1}=0$，给出的正是\BookEqRef{eq:8.24}{8.24}的第一种情形。所以，这些退化对 $A_\infty$ 函子关系的第三个项作出贡献。

\SourceParagraph{ch08-8.4.5-058}{659}
仔细处理所有符号后，便得到 $F$ 是 $A_\infty$ 函子，从而证明\BookRef{thm:8.24}{定理8.24}。
\end{proof}

\SourceParagraph{ch08-8.4.5-059}{660}
为说明这个证明，考虑涉及 $F_2$ 的关系中可能出现的一种情形。该关系为
\begin{equation}\label{eq:8.26}
\begin{aligned}
&m_2^{\operatorname{TMC}}\bigl(F_1(\sigma_{1,2},z^{m_{1,2}}),
F_1(\sigma_{0,1},z^{m_{0,1}})\bigr)\\
&\quad=F_2\bigl((\sigma_{1,2},z^{m_{1,2}}),
m_1^{D(\widehat X)}(\sigma_{0,1},z^{m_{0,1}})\bigr)\\
&\qquad+(-1)^{\dim\sigma_{0,1}-1}
F_2\bigl(m_1^{D(\widehat X)}(\sigma_{1,2},z^{m_{1,2}}),
(\sigma_{0,1},z^{m_{0,1}})\bigr)\\
&\qquad+F_1\bigl(m_2^{D(\widehat X)}
((\sigma_{1,2},z^{m_{1,2}}),(\sigma_{0,1},z^{m_{0,1}}))\bigr).
\end{aligned}
\end{equation}
考虑 $n_1-n_0$ 与 $n_2-n_1$ 都为正，且 $\dim\sigma_{1,2}=\dim\sigma_{0,1}=0$ 的情形，此时模空间 $T_2^{\mathrm{trop}}$ 为一维。\BookRef{fig:8.40}{图40}表示树可能退化的若干方式。
\begin{figure}[H]
\centering
\BookImage{01_Numbered_Figures/Chapter_08/Ch08_Fig40_PDFp660_Bookp650.pdf}
\BookNumberOnlyCaption\label{fig:8.40}
\end{figure}

\SourceParagraph{ch08-8.4.5-060}{660}
树(1)表示单参数族中一棵一般树的一个具体情形，其中起始点 $\phi(v_{0,1})$ 与 $\phi(v_{1,2})$ 位于
\[
p_{0,1}:=\frac{\pi_1(m_{0,1})}{n_1-n_0},
\qquad
p_{1,2}:=\frac{\pi_1(m_{1,2})}{n_2-n_1}
\]
之间。这些点彼此排斥，两条线段在 $p_{0,2}$ 处相遇。可以改变两个端点，但不能独立地改变它们。当一个端点趋向例如 $p_{0,1}$ 时，另一个端点必须趋向 $p_{1,2}$，从而得到\BookRef{fig:8.40}{图40}中(2)所示的退化情形。此时得到一棵通常的热带 Morse 树，对乘积 $m_2^{\operatorname{TMC}}([p_{1,2}],[p_{0,1}])$ 作出贡献；而对于 $i=1,2$，
\[
F_1(\sigma_{i-1,i},z^{m_{i-1,i}})
=q^{\operatorname{ord}_{n_i-n_{i-1}}(m_{i-1,i})}[p_{i-1,i}].
\]
这对 $F_2$ 关系的左侧作出贡献。若反过来，使端点远离 $p_{0,1}$ 与 $p_{1,2}$，则到某一时刻，一个端点会位于相应胞腔的边界，如\BookRef{fig:8.40}{图40}的(3)所示。视首先到达 $(\sigma_{0,1})_{\epsilon_1}^\vee$ 还是 $(\sigma_{1,2})_{\epsilon_2}^\vee$ 的边界而定，所得贡献对应于 $F_2$ 关系右侧的第一项或第二项。

\SourceParagraph{ch08-8.4.5-061}{661}
当 $(\sigma_{0,1})_{\epsilon_1}^\vee$ 与 $(\sigma_{1,2})_{\epsilon_2}^\vee$ 有交集时，还可能出现另一种行为。此时，重叠部分或者很小，或者很大；后一情形发生于 $\sigma_{0,1}=\sigma_{1,2}$。在后一情形中，若 $p_{0,1},p_{1,2}\in(\sigma_{0,1})_{\epsilon_1}^\vee\cap(\sigma_{1,2})_{\epsilon_2}^\vee$，则\BookRef{fig:8.40}{图40}(4)所示的树族有两种退化方式：(5)端点趋向 $p_{0,1}$ 与 $p_{1,2}$，再次对\BookEqRef{eq:8.26}{8.26}的左侧作出贡献；或者(6)端点趋向 $p_{0,2}$，使这些边的长度趋于零，对\BookEqRef{eq:8.26}{8.26}的最后一项作出贡献。

\SourceParagraph{ch08-8.4.5-062}{661}
当然，还有许多其他可能情形，也需要一定工作来保证符号正确；这里略去这一论证。

% 第8.5节；首次阅读翻译，按用户要求不进行额外复核。
\section{Seidel 关于四次 K3 曲面的结果}\label{sec:8.5}
\setcounter{equation}{26}

\SourceParagraph{ch08-8.5-001}{661}
在最后一节，简要讨论 Paul Seidel\BookCite{ref420}的工作：他在四次 K3 曲面这一极不平凡的情形中证明了同调镜像对称。准确地说，设 $X$ 为 $\mathbb P^3$ 中的非奇异四次曲面，配备适当的辛形式，例如 Fubini--Study 度量的 Kähler 形式在 $X$ 上的限制。另一方面，考虑镜像族；这里关注的是它的复结构。取 $\mathbb P^3\times\operatorname{Spec}\mathbb C[[q]]$ 中由
\begin{equation}\label{eq:8.27}
q(x_0^4+x_1^4+x_2^4+x_3^4)+x_0x_1x_2x_3=0
\end{equation}
定义的族，再对下列群的作用取商：
\[
\Gamma_{16}:=
\left\{(a_0,a_1,a_2,a_3)\ \middle|\
 a_i\in\mathbb C,\ a_i^4=1,\ \prod_{i=0}^3a_i=1\right\}
\big/\left\{(a,a,a,a)\mid a^4=1\right\}.
\]
它在 $\mathbb P^3$ 上以对角方式作用，即 $(a_0,a_1,a_2,a_3)$ 将
\[
(x_0,x_1,x_2,x_3)\longmapsto
(a_0x_0,a_1x_1,a_2x_2,a_3x_3).
\]
将上述族对 $\Gamma_{16}$ 的商记为 $\check X\to\operatorname{Spec}\mathbb C[[q]]$。其一般纤维的有理双点经过保典消解，得到域 $\mathbb C((q))$ 上的 K3 曲面 $\widehat X_q$，因而可以考虑范畴 $D^b(\widehat X_q)$，即 $\widehat X_q$ 上凝聚层的有界导出范畴。可将此范畴与 $\Lambda_{\mathrm{nov}}$ 作张量积；后者包含 $\mathbb C((q))$ 的代数闭包。也就是说，将所有态射空间都与 $\Lambda_{\mathrm{nov}}$ 作张量积。Seidel 的结果大体上可表述如下。（先声明一下：本书给出的定义，与 Seidel 论文\BookCite{ref420}所需的定义，在技术上有一些差异。精确陈述请见\BookCite{ref420}。）

\SourceParagraph{ch08-8.5-002}{662}
\BookStatementNumber{33}
\begin{theorem}\label{thm:8.33}
存在 $\psi\in\operatorname{Aut}(\mathbb C[[q]])$，以及三角 $A_\infty$ 范畴之间的等价
\[
D^\pi\operatorname{Fuk}(X)
\cong\psi^*D^b(\widehat X_q)\otimes\Lambda_{\mathrm{nov}}.
\]
\end{theorem}

\SourceParagraph{ch08-8.5-003}{662}
应将 $\psi$ 看作镜像映射，它在两侧的 $q$ 之间变换坐标。因为 $q$ 不是族 $\check X\to\operatorname{Spec}\mathbb C[[q]]$ 的典范坐标，所以预计这一变换并不平凡。不过，目前还不知道这个坐标变换 $\psi$ 是否确实对应于典范坐标。

\SourceParagraph{ch08-8.5-004}{662}
Seidel 的证明与这里给出的椭圆曲线情形的证明很不一样，而且相当复杂。这里只能满足于一个极简略的证明思路。

\SourceParagraph{ch08-8.5-005}{662}
基本想法是在两侧分别找到一小组对象，它们分裂生成相应的范畴。若通过逐次取映射锥、分裂出直和项以及取同构，能够从一组对象得到范畴 $\mathcal A$ 中的任意对象，就称这组对象分裂生成 $\mathcal A$。这些对象所生成的范畴，可以用显式方法研究并加以比较。

\SourceParagraph{ch08-8.5-006}{662}
在凝聚层一侧，Beilinson 定理，即\BookRef{thm:4.72}{定理4.72}，说明 $D^b(\mathbb P^3)$ 由一小组局部自由层生成，即 $\Omega^i(i)[i]$，$0\leq i\leq3$。事实上，这些层在四次 K3 曲面上的限制，分裂生成该曲面的有界导出范畴。在我们的情形中，$\check X$ 被描述为一个 K3 曲面族 $\mathcal X$ 对 $\Gamma_{16}$ 的商。因而，利用导出 McKay 对应，可以用商空间保典消解的导出范畴所对应的、覆盖空间上的 $\Gamma_{16}$ 等变导出范畴来研究问题。覆盖空间上的每个分裂生成元都有 16 个等变版本，从而得到 $D^b(\widehat X_q)$ 的一组 64 个分裂生成元。

\SourceParagraph{ch08-8.5-007}{662}
在另一侧，所考虑的 Fermat 族\BookEqRef{eq:8.27}{8.27}，作为 $\operatorname{Spec}\mathbb C[q]$ 上的族，除了 $q=0$ 处的退化纤维外，还有四个奇异纤维，每个纤维都有 16 个奇异点。每个这样的奇异点，都会在此族中某个选定的一般纤维内诱导一个作为消灭闭链的 Lagrange 球面。这样便在该纤维中得到 64 个 Lagrange 球面；适当地选择它们，就得到一组对象，用来与导出范畴中的对象比较。可以证明，这 64 个 Lagrange 子流形分裂生成 $D^\pi\operatorname{Fuk}(X)$。最精细的一步，是计算这些对象所生成的子范畴；这需要大量计算，特别需要计算若干乘法映射。最终，由此可以比较这两个范畴。

\SourceParagraph{ch08-8.5-008}{663}
事实上，可以证明，两侧的范畴都是同一个范畴的非平凡变形。在两侧，$q$ 都是变形参数。通过 Hochschild 上同调的计算，关于 Hochschild 上同调的一些讨论见\BookSectionRef{sec:2.2.3}{2.2.3}，能够证明，在所考虑的对称性约束下，只有一个单参数变形族。因此，这两个族必定相同，但无法给出显式的镜像映射。

